%%%%%%%% ICML 2026 EXAMPLE LATEX SUBMISSION FILE %%%%%%%%%%%%%%%%%

\documentclass{article}

% Recommended, but optional, packages for figures and better typesetting:
\usepackage{microtype}
\usepackage{graphicx}
\usepackage{subcaption}
\usepackage{booktabs} % for professional tables

% hyperref makes hyperlinks in the resulting PDF.
% If your build breaks (sometimes temporarily if a hyperlink spans a page)
% please comment out the following usepackage line and replace
% \usepackage{icml2026} with \usepackage[nohyperref]{icml2026} above.
\usepackage{hyperref}

%Package Add By Liying
\usepackage{multirow}

% Attempt to make hyperref and algorithmic work together better:
\newcommand{\theHalgorithm}{\arabic{algorithm}}

% Use the following line for the initial blind version submitted for review:
\usepackage{icml2026}

% For preprint, use
% \usepackage[preprint]{icml2026}

% If accepted, instead use the following line for the camera-ready submission:
% \usepackage[accepted]{icml2026}

\usepackage{amsmath}
\usepackage{amssymb}
\usepackage{mathtools}
\usepackage{amsthm}

\newcommand{\minisection}[1]{\vspace{0.0in} \noindent {\bf #1}}

% if you use cleveref..
\usepackage[capitalize,noabbrev]{cleveref}

%%%%%%%%%%%%%%%%%%%%%%%%%%%%%%%%
% THEOREMS
%%%%%%%%%%%%%%%%%%%%%%%%%%%%%%%%
\theoremstyle{plain}
\newtheorem{theorem}{Theorem}[section]
\newtheorem{proposition}[theorem]{Proposition}
\newtheorem{lemma}[theorem]{Lemma}
\newtheorem{corollary}[theorem]{Corollary}
\theoremstyle{definition}
\newtheorem{definition}[theorem]{Definition}
\newtheorem{assumption}[theorem]{Assumption}
\theoremstyle{remark}
\newtheorem{remark}[theorem]{Remark}

% Todonotes is useful during development; simply uncomment the next line
%    and comment out the line below the next line to turn off comments
%\usepackage[disable,textsize=tiny]{todonotes}
\usepackage[textsize=tiny]{todonotes}

%\newcommand{\SM}[1]{{\color{orange}{\bf SM:} #1}}
%\newcommand{\YX}[1]{{\color{brown}{\bf YX:} #1}}
\newcommand{\JW}[1]{{\color{magenta}{\bf JW:} #1}}

% The \icmltitle you define below is probably too long as a header.
% Therefore, a short form for the running title is supplied here:
\icmltitlerunning{Revisiting Sharpness in Low-rank Subspaces for Continual Learning}

\begin{document}

\twocolumn[
  \icmltitle{
  Revisiting Sharpness in Low-rank Subspaces for Continual Learning}
  % Adapter Subspace Sharpness Matters for LoRA-Based Continual Learning: A PAC-Bayesian Perspective}

  % It is OKAY to include author information, even for blind submissions: the
  % style file will automatically remove it for you unless you've provided
  % the [accepted] option to the icml2026 package.

  % List of affiliations: The first argument should be a (short) identifier you
  % will use later to specify author affiliations Academic affiliations
  % should list Department, University, City, Region, Country Industry
  % affiliations should list Company, City, Region, Country

  % You can specify symbols, otherwise they are numbered in order. Ideally, you
  % should not use this facility. Affiliations will be numbered in order of
  % appearance and this is the preferred way.
  \icmlsetsymbol{equal}{*}

  \begin{icmlauthorlist}
    \icmlauthor{Firstname1 Lastname1}{equal,yyy}
    \icmlauthor{Firstname2 Lastname2}{equal,yyy,comp}
    \icmlauthor{Firstname3 Lastname3}{comp}
    \icmlauthor{Firstname4 Lastname4}{sch}
    \icmlauthor{Firstname5 Lastname5}{yyy}
    \icmlauthor{Firstname6 Lastname6}{sch,yyy,comp}
    \icmlauthor{Firstname7 Lastname7}{comp}
    %\icmlauthor{}{sch}
    \icmlauthor{Firstname8 Lastname8}{sch}
    \icmlauthor{Firstname8 Lastname8}{yyy,comp}
    %\icmlauthor{}{sch}
    %\icmlauthor{}{sch}
  \end{icmlauthorlist}

  \icmlaffiliation{yyy}{Department of XXX, University of YYY, Location, Country}
  \icmlaffiliation{comp}{Company Name, Location, Country}
  \icmlaffiliation{sch}{School of ZZZ, Institute of WWW, Location, Country}

  \icmlcorrespondingauthor{Firstname1 Lastname1}{first1.last1@xxx.edu}
  \icmlcorrespondingauthor{Firstname2 Lastname2}{first2.last2@www.uk}

  % You may provide any keywords that you find helpful for describing your
  % paper; these are used to populate the "keywords" metadata in the PDF but
  % will not be shown in the document
  \icmlkeywords{Machine Learning, ICML}

  \vskip 0.3in
]

% this must go after the closing bracket ] following \twocolumn[ ...

% This command actually creates the footnote in the first column listing the
% affiliations and the copyright notice. The command takes one argument, which
% is text to display at the start of the footnote. The \icmlEqualContribution
% command is standard text for equal contribution. Remove it (just {}) if you
% do not need this facility.

% Use ONE of the following lines. DO NOT remove the command.
% If you have no special notice, KEEP empty braces:
\printAffiliationsAndNotice{}  % no special notice (required even if empty)
% Or, if applicable, use the standard equal contribution text:
% \printAffiliationsAndNotice{\icmlEqualContribution}

\begin{abstract}
Low-Rank Adaptation (LoRA) enables parameter-efficient continual learning (PECL) for large pre-trained models by confining task updates to a low-rank, \emph{task-permissible subspace}.
While sharpness-aware optimization improves generalization in full-model learning, its role under subspace constrained updates remains unclear. This prompts the question of whether controlling \emph{subspace sharpness} (sharpness within the {task-permissible subspace}) suffices to improve generalization across tasks, and how it relates to the sharpness of the full model.
In this paper we define the subspace sharpness and derive a sequential hierarchical PAC-Bayesian bound that augments empirical risk with a subspace sharpness term and a hierarchical KL term for cross-task posterior complexity. We compare subspace-restricted and full-parameter perturbation scopes and evaluate multiple sharpness types within the task-permissible subspace.
Curvature diagnostics further show that lower subspace sharpness is associated with higher continual accuracy and less forgetting across tasks. 
\end{abstract}



% \begin{abstract}
% % Low-Rank Adaptation (LoRA) offers a scalable route to Parameter Efficient Continual Learning (PECL) with large pre-trained models, yet shared low-rank subspaces make it difficult to balance stability and plasticity.
% % In LoRA-based continual learning, each task is learned 
% Low-Rank Adaptation (LoRA) enables parameter-efficient continual learning (PECL) for large pre-trained models by confining task updates to a low-rank, \emph{task-permissible subspace}. While sharpness-aware optimization improves generalization in full-model learning, its role under subspace constrained updates remains unclear. This prompts the question of whether controlling \emph{subspace sharpness} ( sharpness within the {task-permissible subspace}) suffices to improve generalization across tasks, and how it relates to the sharpness of the full model.

% % Low-Rank Adaptation (LoRA) offers a scalable route for parameter efficient continual learning (PECL) with large pre-trained models
% % by restricting task updates to a task-permissible subspace.
% % by updating weights only within a low-rank, task permissible subspace.
% While sharpness-aware optimization improves generalization in full-model learning, its role under subspace constrained updates remains unclear. This prompts the question of whether controlling \emph{subspace sharpness} ( sharpness within the {task-permissible subspace}) suffices to improve generalization across tasks, and how it relates to the sharpness of the full model.
% % The key question is whether suppressing curvature within task permissible subspace is enough to improve continual generalization and stability, and how it relates to the sharpness of the full model.
% % remains unclear. 
% % Parameter-efficient continual learning with LoRA confines each task update to a task-permissible low-rank subspace, 
% % It motivate the question of whether controlling curvature within this subspace suffices to improve continual generalization and stability.
% % creating a geometric mismatch with sharpness objectives defined over full-parameter perturbations. 
% % The key question is whether suppressing sharpness within task permissible subspace is enough to improve continual generalization and stability, and how it relates to the sharpness of the full model.
% % has not been theoretically established. 
% % role inside the low-rank subspace for CL has not been theoretically established. 
% % We address this gap by revisiting \emph{low-rank subspace sharpness} for CL. 
% % We define task-conditional sharpness restricted to the LoRA update subspace and derive a sequential hierarchical PAC-Bayesian bound for PECL. 
% %restricted to the update subspace
% % We revisit sharpness in LoRA based continual learning by defining task subspace sharpness and derive a sequential hierarchical PAC Bayesian bound.
% % The bound yields an empirical risk term with a subspace sharpness term and a hierarchical KL term for cross-task posterior complexity.
% % We derive a sequential hierarchical PAC-Bayesian  bound that augments empirical risk with an subspace sharpness term and a hierarchical KL term that tracks posterior complexity across tasks. 
% In this paper we define the subspace sharpness and derive a sequential hierarchical PAC-Bayesian bound that augments empirical risk with a subspace sharpness term and a hierarchical KL term for cross-task posterior complexity. We compare subspace-restricted and full-parameter perturbation scopes and evaluate multiple sharpness types within the task-permissible subspace.
% % We analyze perturbation design in LoRA-based PECL along two axes: \emph{where} to perturb and \emph{how} to perturb. 
% % We discuss how these choices interact with LoRA organization over the task sequence. 
% % Using a unified protocol 
% % Across LoRA variants and perturbation type, we demonstrate that controlling sharpness within the adapter subspace improves accuracy over full-parameter perturbations. 
% % Restricting perturbations to the trainable LoRA coordinates consistently improves performance, whereas full parameter perturbations degrade accuracy at matched perturbation scale. Within the LoRA restricted scope, more curvature aware objectives perform better, and first order sharpness minimization yields the strongest gains. 
% Curvature diagnostics further show that lower subspace sharpness is associated with higher continual accuracy and less forgetting across tasks. 
% % Curvature diagnostics link reduced subspace sharpness to improved stability, with flatter subspace solutions exhibiting higher continual accuracy and less forgetting across tasks.
% % TO DO.
% % in the LoRA-Based CL.
% % Our results provide a theory-to-practice account and design guidance for PECL.
% \end{abstract}

% \begin{abstract}
%   This document provides a basic paper template and submission guidelines.
%   Abstracts must be a single paragraph, ideally between 4--6 sentences long.
%   Gross violations will trigger corrections at the camera-ready phase.
% \end{abstract}

% \section{Electronic Submission}

% Submission to ICML 2026 will be entirely electronic, via a web site
% (not email). Information about the submission process and \LaTeX\ templates
% are available on the conference web site at:
% \begin{center}
%   \texttt{http://icml.cc/}
% \end{center}

% The guidelines below will be enforced for initial submissions and
% camera-ready copies. Here is a brief summary:
% \begin{itemize}
%   \item Submissions must be in PDF\@.
%   \item If your paper has appendices, submit the appendix together with the
%         main body and the references \textbf{as a single file}. Reviewers will not
%         look for appendices as a separate PDF file. So if you submit such an extra
%         file, reviewers will very likely miss it.
%   \item Page limit: The main body of the paper has to be fitted to 8 pages,
%         excluding references and appendices; the space for the latter two is not
%         limited in pages, but the total file size may not exceed 10MB. For the
%         final version of the paper, authors can add one extra page to the main
%         body.
%   \item \textbf{Do not include author information or acknowledgements} in your
%         initial submission.
%   \item Your paper should be in \textbf{10 point Times font}.
%   \item Make sure your PDF file only uses Type-1 fonts.
%   \item Place figure captions \emph{under} the figure (and omit titles from
%         inside the graphic file itself). Place table captions \emph{over} the
%         table.
%   \item References must include page numbers whenever possible and be as
%         complete as possible. Place multiple citations in chronological order.
%   \item Do not alter the style template; in particular, do not compress the
%         paper format by reducing the vertical spaces.
%   \item Keep your abstract brief and self-contained, one paragraph and roughly
%         4--6 sentences. Gross violations will require correction at the
%         camera-ready phase. The title should have content words capitalized.
% \end{itemize}

% \subsection{Submitting Papers}

% \textbf{Anonymous Submission:} ICML uses double-blind review: no identifying
% author information may appear on the title page or in the paper
% itself. \cref{author info} gives further details.

% \medskip

% Authors must provide their manuscripts in \textbf{PDF} format.
% Furthermore, please make sure that files contain only embedded Type-1 fonts
% (e.g.,~using the program \texttt{pdffonts} in linux or using
% File/DocumentProperties/Fonts in Acrobat). Other fonts (like Type-3)
% might come from graphics files imported into the document.

% Authors using \textbf{Word} must convert their document to PDF\@. Most
% of the latest versions of Word have the facility to do this
% automatically. Submissions will not be accepted in Word format or any
% format other than PDF\@. Really. We're not joking. Don't send Word.

% Those who use \textbf{\LaTeX} should avoid including Type-3 fonts.
% Those using \texttt{latex} and \texttt{dvips} may need the following
% two commands:

% {\footnotesize
% \begin{verbatim}
% dvips -Ppdf -tletter -G0 -o paper.ps paper.dvi
% ps2pdf paper.ps
% \end{verbatim}}
% It is a zero following the ``-G'', which tells dvips to use
% the config.pdf file. Newer \TeX\ distributions don't always need this
% option.

% Using \texttt{pdflatex} rather than \texttt{latex}, often gives better
% results. This program avoids the Type-3 font problem, and supports more
% advanced features in the \texttt{microtype} package.

% \textbf{Graphics files} should be a reasonable size, and included from
% an appropriate format. Use vector formats (.eps/.pdf) for plots,
% lossless bitmap formats (.png) for raster graphics with sharp lines, and
% jpeg for photo-like images.

% The style file uses the \texttt{hyperref} package to make clickable
% links in documents. If this causes problems for you, add
% \texttt{nohyperref} as one of the options to the \texttt{icml2026}
% usepackage statement.

% \subsection{Submitting Final Camera-Ready Copy}

% The final versions of papers accepted for publication should follow the
% same format and naming convention as initial submissions, except that
% author information (names and affiliations) should be given. See
% \cref{final author} for formatting instructions.

% The footnote, ``Preliminary work. Under review by the International
% Conference on Machine Learning (ICML). Do not distribute.'' must be
% modified to ``\textit{Proceedings of the
%   $\mathit{43}^{rd}$ International Conference on Machine Learning},
% Seoul, South Korea, PMLR 306, 2026.
% Copyright 2026 by the author(s).''

% For those using the \textbf{\LaTeX} style file, this change (and others) is
% handled automatically by simply changing
% $\mathtt{\backslash usepackage\{icml2026\}}$ to
% $$\mathtt{\backslash usepackage[accepted]\{icml2026\}}$$
% Authors using \textbf{Word} must edit the
% footnote on the first page of the document themselves.

% Camera-ready copies should have the title of the paper as running head
% on each page except the first one. The running title consists of a
% single line centered above a horizontal rule which is $1$~point thick.
% The running head should be centered, bold and in $9$~point type. The
% rule should be $10$~points above the main text. For those using the
% \textbf{\LaTeX} style file, the original title is automatically set as running
% head using the \texttt{fancyhdr} package which is included in the ICML
% 2026 style file package. In case that the original title exceeds the
% size restrictions, a shorter form can be supplied by using

% \verb|\icmltitlerunning{...}|

% just before $\mathtt{\backslash begin\{document\}}$.
% Authors using \textbf{Word} must edit the header of the document themselves.

% \section{Format of the Paper}

% All submissions must follow the specified format.

% \subsection{Dimensions}

% The text of the paper should be formatted in two columns, with an
% overall width of 6.75~inches, height of 9.0~inches, and 0.25~inches
% between the columns. The left margin should be 0.75~inches and the top
% margin 1.0~inch (2.54~cm). The right and bottom margins will depend on
% whether you print on US letter or A4 paper, but all final versions
% must be produced for US letter size.
% Do not write anything on the margins.

% The paper body should be set in 10~point type with a vertical spacing
% of 11~points. Please use Times typeface throughout the text.

% \subsection{Title}

% The paper title should be set in 14~point bold type and centered
% between two horizontal rules that are 1~point thick, with 1.0~inch
% between the top rule and the top edge of the page. Capitalize the
% first letter of content words and put the rest of the title in lower
% case.
% You can use TeX math in the title (we suggest sparingly),
% but no custom macros, images, or other TeX commands.
% Please make sure that accents, special characters, etc., are entered using
% TeX commands and not using non-English characters.

% \subsection{Author Information for Submission}
% \label{author info}

% ICML uses double-blind review, so author information must not appear. If
% you are using \LaTeX\/ and the \texttt{icml2026.sty} file, use
% \verb+\icmlauthor{...}+ to specify authors and \verb+\icmlaffiliation{...}+
% to specify affiliations. (Read the TeX code used to produce this document for
% an example usage.) The author information will not be printed unless
% \texttt{accepted} is passed as an argument to the style file. Submissions that
% include the author information will not be reviewed.

% \subsubsection{Self-Citations}

% If you are citing published papers for which you are an author, refer
% to yourself in the third person. In particular, do not use phrases
% that reveal your identity (e.g., ``in previous work \cite{langley00}, we
% have shown \ldots'').

% Do not anonymize citations in the reference section. The only exception are manuscripts that are
% not yet published (e.g., under submission). If you choose to refer to
% such unpublished manuscripts \cite{anonymous}, anonymized copies have
% to be submitted
% as Supplementary Material via OpenReview\@. However, keep in mind that an ICML
% paper should be self contained and should contain sufficient detail
% for the reviewers to evaluate the work. In particular, reviewers are
% not required to look at the Supplementary Material when writing their
% review (they are not required to look at more than the first $8$ pages of the submitted document).

% \subsubsection{Camera-Ready Author Information}
% \label{final author}

% If a paper is accepted, a final camera-ready copy must be prepared.
% %
% For camera-ready papers, author information should start 0.3~inches below the
% bottom rule surrounding the title. The authors' names should appear in 10~point
% bold type, in a row, separated by white space, and centered. Author names should
% not be broken across lines. Unbolded superscripted numbers, starting 1, should
% be used to refer to affiliations.

% Affiliations should be numbered in the order of appearance. A single footnote
% block of text should be used to list all the affiliations. (Academic
% affiliations should list Department, University, City, State/Region, Country.
% Similarly for industrial affiliations.)

% Each distinct affiliations should be listed once. If an author has multiple
% affiliations, multiple superscripts should be placed after the name, separated
% by thin spaces. If the authors would like to highlight equal contribution by
% multiple first authors, those authors should have an asterisk placed after their
% name in superscript, and the term ``\textsuperscript{*}Equal contribution"
% should be placed in the footnote block ahead of the list of affiliations. A
% list of corresponding authors and their emails (in the format Full Name
% \textless{}email@domain.com\textgreater{}) can follow the list of affiliations.
% Ideally only one or two names should be listed.

% A sample file with author names is included in the ICML2026 style file
% package. Turn on the \texttt{[accepted]} option to the stylefile to
% see the names rendered. All of the guidelines above are implemented
% by the \LaTeX\ style file.

% \subsection{Abstract}

% The paper abstract should begin in the left column, 0.4~inches below the final
% address. The heading `Abstract' should be centered, bold, and in 11~point type.
% The abstract body should use 10~point type, with a vertical spacing of
% 11~points, and should be indented 0.25~inches more than normal on left-hand and
% right-hand margins. Insert 0.4~inches of blank space after the body. Keep your
% abstract brief and self-contained, limiting it to one paragraph and roughly 4--6
% sentences. Gross violations will require correction at the camera-ready phase.

% \subsection{Partitioning the Text}

% You should organize your paper into sections and paragraphs to help readers
% place a structure on the material and understand its contributions.

% \subsubsection{Sections and Subsections}

% Section headings should be numbered, flush left, and set in 11~pt bold type
% with the content words capitalized. Leave 0.25~inches of space before the
% heading and 0.15~inches after the heading.

% Similarly, subsection headings should be numbered, flush left, and set in 10~pt
% bold type with the content words capitalized. Leave
% 0.2~inches of space before the heading and 0.13~inches afterward.

% Finally, subsubsection headings should be numbered, flush left, and set in
% 10~pt small caps with the content words capitalized. Leave
% 0.18~inches of space before the heading and 0.1~inches after the heading.

% Please use no more than three levels of headings.

% \subsubsection{Paragraphs and Footnotes}

% Within each section or subsection, you should further partition the paper into
% paragraphs. Do not indent the first line of a given paragraph, but insert a
% blank line between succeeding ones.

% You can use footnotes\footnote{Footnotes should be complete sentences.}
% to provide readers with additional information about a topic without
% interrupting the flow of the paper. Indicate footnotes with a number in the
% text where the point is most relevant. Place the footnote in 9~point type at
% the bottom of the column in which it appears. Precede the first footnote in a
% column with a horizontal rule of 0.8~inches.\footnote{Multiple footnotes can
%   appear in each column, in the same order as they appear in the text,
%   but spread them across columns and pages if possible.}

% \begin{figure}[ht]
%   \vskip 0.2in
%   \begin{center}
%     \centerline{\includegraphics[width=\columnwidth]{icml_numpapers}}
%     \caption{
%       Historical locations and number of accepted papers for International
%       Machine Learning Conferences (ICML 1993 -- ICML 2008) and International
%       Workshops on Machine Learning (ML 1988 -- ML 1992). At the time this
%       figure was produced, the number of accepted papers for ICML 2008 was
%       unknown and instead estimated.
%     }
%     \label{icml-historical}
%   \end{center}
% \end{figure}

% \subsection{Figures}

% You may want to include figures in the paper to illustrate your approach and
% results. Such artwork should be centered, legible, and separated from the text.
% Lines should be dark and at least 0.5~points thick for purposes of
% reproduction, and text should not appear on a gray background.

% Label all distinct components of each figure. If the figure takes the form of a
% graph, then give a name for each axis and include a legend that briefly
% describes each curve. Do not include a title inside the figure; instead, the
% caption should serve this function.

% Number figures sequentially, placing the figure number and caption \emph{after}
% the graphics, with at least 0.1~inches of space before the caption and
% 0.1~inches after it, as in \cref{icml-historical}. The figure caption should be
% set in 9~point type and centered unless it runs two or more lines, in which
% case it should be flush left. You may float figures to the top or bottom of a
% column, and you may set wide figures across both columns (use the environment
% \texttt{figure*} in \LaTeX). Always place two-column figures at the top or
% bottom of the page.

% \subsection{Algorithms}

% If you are using \LaTeX, please use the ``algorithm'' and ``algorithmic''
% environments to format pseudocode. These require the corresponding stylefiles,
% algorithm.sty and algorithmic.sty, which are supplied with this package.
% \cref{alg:example} shows an example.

% \begin{algorithm}[tb]
%   \caption{Bubble Sort}
%   \label{alg:example}
%   \begin{algorithmic}
%     \STATE {\bfseries Input:} data $x_i$, size $m$
%     \REPEAT
%     \STATE Initialize $noChange = true$.
%     \FOR{$i=1$ {\bfseries to} $m-1$}
%     \IF{$x_i > x_{i+1}$}
%     \STATE Swap $x_i$ and $x_{i+1}$
%     \STATE $noChange = false$
%     \ENDIF
%     \ENDFOR
%     \UNTIL{$noChange$ is $true$}
%   \end{algorithmic}
% \end{algorithm}


% \subsection{Tables}

% You may also want to include tables that summarize material. Like figures,
% these should be centered, legible, and numbered consecutively. However, place
% the title \emph{above} the table with at least 0.1~inches of space before the
% title and the same after it, as in \cref{sample-table}. The table title should
% be set in 9~point type and centered unless it runs two or more lines, in which
% case it should be flush left.

% % Note use of \abovespace and \belowspace to get reasonable spacing
% % above and below tabular lines.

% \begin{table}[t]
%   \caption{Classification accuracies for naive Bayes and flexible
%     Bayes on various data sets.}
%   \label{sample-table}
%   \begin{center}
%     \begin{small}
%       \begin{sc}
%         \begin{tabular}{lcccr}
%           \toprule
%           Data set  & Naive         & Flexible      & Better?  \\
%           \midrule
%           Breast    & 95.9$\pm$ 0.2 & 96.7$\pm$ 0.2 & $\surd$  \\
%           Cleveland & 83.3$\pm$ 0.6 & 80.0$\pm$ 0.6 & $\times$ \\
%           Glass2    & 61.9$\pm$ 1.4 & 83.8$\pm$ 0.7 & $\surd$  \\
%           Credit    & 74.8$\pm$ 0.5 & 78.3$\pm$ 0.6 &          \\
%           Horse     & 73.3$\pm$ 0.9 & 69.7$\pm$ 1.0 & $\times$ \\
%           Meta      & 67.1$\pm$ 0.6 & 76.5$\pm$ 0.5 & $\surd$  \\
%           Pima      & 75.1$\pm$ 0.6 & 73.9$\pm$ 0.5 &          \\
%           Vehicle   & 44.9$\pm$ 0.6 & 61.5$\pm$ 0.4 & $\surd$  \\
%           \bottomrule
%         \end{tabular}
%       \end{sc}
%     \end{small}
%   \end{center}
%   \vskip -0.1in
% \end{table}

% Tables contain textual material, whereas figures contain graphical material.
% Specify the contents of each row and column in the table's topmost row. Again,
% you may float tables to a column's top or bottom, and set wide tables across
% both columns. Place two-column tables at the top or bottom of the page.

% \subsection{Theorems and Such}
% The preferred way is to number definitions, propositions, lemmas, etc.
% consecutively, within sections, as shown below.
% \begin{definition}
%   \label{def:inj}
%   A function $f:X \to Y$ is injective if for any $x,y\in X$ different, $f(x)\ne
%     f(y)$.
% \end{definition}
% Using \cref{def:inj} we immediate get the following result:
% \begin{proposition}
%   If $f$ is injective mapping a set $X$ to another set $Y$,
%   the cardinality of $Y$ is at least as large as that of $X$
% \end{proposition}
% \begin{proof}
%   Left as an exercise to the reader.
% \end{proof}
% \cref{lem:usefullemma} stated next will prove to be useful.
% \begin{lemma}
%   \label{lem:usefullemma}
%   For any $f:X \to Y$ and $g:Y\to Z$ injective functions, $f \circ g$ is
%   injective.
% \end{lemma}
% \begin{theorem}
%   \label{thm:bigtheorem}
%   If $f:X\to Y$ is bijective, the cardinality of $X$ and $Y$ are the same.
% \end{theorem}
% An easy corollary of \cref{thm:bigtheorem} is the following:
% \begin{corollary}
%   If $f:X\to Y$ is bijective,
%   the cardinality of $X$ is at least as large as that of $Y$.
% \end{corollary}
% \begin{assumption}
%   The set $X$ is finite.
%   \label{ass:xfinite}
% \end{assumption}
% \begin{remark}
%   According to some, it is only the finite case (cf. \cref{ass:xfinite}) that
%   is interesting.
% \end{remark}
% %restatable

% \subsection{Citations and References}

% Please use APA reference format regardless of your formatter or word processor.
% If you rely on the \LaTeX\/ bibliographic facility, use \texttt{natbib.sty} and
% \texttt{icml2026.bst} included in the style-file package to obtain this format.

% Citations within the text should include the authors' last names and year. If
% the authors' names are included in the sentence, place only the year in
% parentheses, for example when referencing Arthur Samuel's pioneering work
% \yrcite{Samuel59}. Otherwise place the entire reference in parentheses with the
% authors and year separated by a comma \cite{Samuel59}. List multiple references
% separated by semicolons \cite{kearns89,Samuel59,mitchell80}. Use the `et~al.'
% construct only for citations with three or more authors or after listing all
% authors to a publication in an earlier reference \cite{MachineLearningI}.

% Authors should cite their own work in the third person in the initial version
% of their paper submitted for blind review. Please refer to \cref{author info}
% for detailed instructions on how to cite your own papers.

% Use an unnumbered first-level section heading for the references, and use a
% hanging indent style, with the first line of the reference flush against the
% left margin and subsequent lines indented by 10 points. The references at the
% end of this document give examples for journal articles \cite{Samuel59},
% conference publications \cite{langley00}, book chapters \cite{Newell81}, books
% \cite{DudaHart2nd}, edited volumes \cite{MachineLearningI}, technical reports
% \cite{mitchell80}, and dissertations \cite{kearns89}.

% Alphabetize references by the surnames of the first authors, with single author
% entries preceding multiple author entries. Order references for the same
% authors by year of publication, with the earliest first. Make sure that each
% reference includes all relevant information (e.g., page numbers).

% Please put some effort into making references complete, presentable, and
% consistent, e.g. use the actual current name of authors. If using bibtex,
% please protect capital letters of names and abbreviations in titles, for
% example, use \{B\}ayesian or \{L\}ipschitz in your .bib file.

% \section*{Accessibility}

% Authors are kindly asked to make their submissions as accessible as possible
% for everyone including people with disabilities and sensory or neurological
% differences. Tips of how to achieve this and what to pay attention to will be
% provided on the conference website \url{http://icml.cc/}.

% \section*{Software and Data}

% If a paper is accepted, we strongly encourage the publication of software and
% data with the camera-ready version of the paper whenever appropriate. This can
% be done by including a URL in the camera-ready copy. However, \textbf{do not}
% include URLs that reveal your institution or identity in your submission for
% review. Instead, provide an anonymous URL or upload the material as
% ``Supplementary Material'' into the OpenReview reviewing system. Note that
% reviewers are not required to look at this material when writing their review.

% % Acknowledgements should only appear in the accepted version.
% \section*{Acknowledgements}

% \textbf{Do not} include acknowledgements in the initial version of the paper
% submitted for blind review.

% If a paper is accepted, the final camera-ready version can (and usually should)
% include acknowledgements.  Such acknowledgements should be placed at the end of
% the section, in an unnumbered section that does not count towards the paper
% page limit. Typically, this will include thanks to reviewers who gave useful
% comments, to colleagues who contributed to the ideas, and to funding agencies
% and corporate sponsors that provided financial support.

\section{Introduction}





% Adapting large-scale pre-trained models to new tasks is now routine in modern machine learning~\cite{HAN2021225}.
% Continual learning (CL) studies sequential task learning while maintaining generalization on previously learned tasks~\cite{wangComprehensiveSurveyContinual2024}. Catastrophic forgetting is the deterioration of this past-task generalization, measured as increased test risk on earlier task distributions after learning new tasks~\cite{pmlr-v202-lin23f}.
% To reduce the computational and storage cost of full-model updates, parameter-efficient methods such as Low-Rank Adaptation (LoRA)\cite{hu2022lora} introduce lightweight task-specific modules while freezing the backbone weights~\cite{dingParameterefficientFinetuningLargescale2023,hanParameterEfficientFineTuningLarge2024}. Combining parameter-efficient adaptation with continual learning objectives yields parameter-efficient continual learning (PECL), which enables scalable sequential adaptation under restricted update subspaces but can still suffer from catastrophic forgetting~\cite{gaoUnifiedContinualLearning2023,qiaoLearnMoreBother2024}. 
% Flat minima are widely associated with improved generalization and robustness, motivating Sharpness-aware minimization (SAM)~\cite{foretSharpnessAwareMinimizationEfficiently2021} and related methods to control sharpness under local perturbations in the full parameter space. Beyond single task setting, such flatness control has been argued to mitigate forgetting in CL~\cite{NEURIPS2020_518a38cc,JMLRAnEmpiricalInvestigatio,shiOvercomingCatastrophicForgetting2021,liRevisitingCatastrophicForgetting2024}.

Parameter-efficient continual learning (PECL)~\cite{Gao_2023_ICCV,qiaoLearnMoreBother2024} adapts a frozen pre-trained backbone to a task sequence by updating only lightweight modules such as LoRA~\cite{hu_lora_2021}, reducing training and storage costs while preserving and reusing the backbone’s pre-trained knowledge. While this design makes continual adaptation scalable, it constrains learning to a low-dimensional, \emph{task-permissible subspace}, reshaping the local loss geometry along feasible directions and potentially exacerbating generalization fragility and forgetting~\cite{NEURIPS2024_4eb2c0ad,li2024flatlora}.
% While this design makes continual adaptation scalable, it constrains learning to a low-dimensional, task-permissible subspace, reshaping the local loss geometry along feasible directions and potentially exacerbate generalization fragility and forgetting~\citep{NEURIPS2024_4eb2c0ad,li2024flatlora}.
Flat minima in the loss landscape have long been associated with improved generalization and robustness~\citep{FlatMinima1997,wu2017understandinggeneralizationdeeplearning}. 
% A natural response to such sharp local geometry is to bias optimization toward flat minima by injecting perturbations during training.  operationalize this principle
Building on this connection, Sharpness-aware minimization (SAM)~\cite{foretSharpnessAwareMinimizationEfficiently2021} and related methods~\citep{pmlr-v151-bisla22a,Zhang_GAM,NEURIPS2024_C-Flat} incorporate local perturbations into the training objective to penalize sharp solutions and bias optimization toward flat minima.
% Beyond improved generalization, 
Prior work has further suggested that improved flatness can mitigate forgetting in continual learning~\citep{NEURIPS2020_518a38cc,JMLRAnEmpiricalInvestigatio,liRevisitingCatastrophicForgetting2024}.

% optimize the loss with sharpness by injecting perturbations to bais soultion to flat minima.

% by seeking parameters that maintain low loss under small perturbations. 



% as tasks accumulate




% This makes sequential adaptation scalable,
% % This yields scalable sequential adaptation, 
% but it also restricts learning to task permissible update directions induced by the adapters, which can aggravate generalization under sharp local geometry~\cite{NEURIPS2024_4eb2c0ad,li2024flatlora}.
% %which can make generalization more fragile in the presence of sharp local geometry
% A common strategy to improve generalization is to inject parameter perturbations during training and optimize a sharpness-related objective, encouraging convergence toward flatter solutions, including sharpness-aware minimization (SAM)~\cite{foretSharpnessAwareMinimizationEfficiently2021} and related methods~\cite{pmlr-v151-bisla22a,Zhang_GAM,NEURIPS2024_C-Flat,yangRevisitingFlatnessAwareOptimization2025,chen2026sharpnessflatnessdecompositionframework}.



{
\begin{figure}[tbp]
    \centering
\includegraphics[width=0.93\linewidth]{figures_TRML/liuxingshiyitu5.pdf}
    % \vspace{-6pt}
    \caption{\textbf{Full-parameter perturbations versus restricted subspace perturbations.} PECL updates are confined to task-permissible spaces (green), but standard sharpness perturbs in the full parameter space (blue).
    % creating a conflict between possible perturbation geometry and feasible update dynamics. 
    Subspace perturbations within the task-permissible scope (red) align the possible perturbation geometry with the constrained update dynamics.
    % Subspace perturbations within the task-permissible scope (red) steer the optimization by emphasizing curvature along the actually reachable directions, so the sharpness signal directly influences the subsequent task-admissible updates.
    % The green ribbon denotes the task-admissible update subspace and black dots and trajectory illustrate  optimization path within this subspace. 
    % Around the current solution, the red dashed ellipse represents a subspace-restricted perturbation, the blue dashed ellipse represents a full-space perturbation.
    % PECL updates are confined to task-admissible directions (green), but standard sharpness perturbs in full-parameter neighborhoods (blue), creating a scope mismatch between perturbation geometry and feasible update dynamics. However subspace-restricted perturbations within this scope (red) align the perturbation geometry with the constrained update dynamics.
    }
    \vspace{-15pt}
    \label{fig:scope_type1}
\end{figure}
}


% Adapting large-scale pre-trained models to new data is now routine in modern machine learning. Due to the scale of contemporary models, this adaptation is typically performed using parameter-efficient transfer learning methods~\cite{han2024parameter} such as low-rank adaptation (LoRA)~\cite{hu2022lora}, which restrict task-specific updates to low-rank adapter subspaces. 

% Adapting large-scale pre-trained models to new data is typically performed using parameter-efficient transfer learning methods~\cite{han2024parameter} such as low-rank adaptation (LoRA)~\cite{hu2022lora}, which restrict task-specific updates to low-rank adapter subspaces. 
% % commonly achieved through parameter-efficient transfer learning methods, which restrict task-specific updates to low-rank adapter subspaces, such as LoRA~\cite{hu2022lora,han2024parameter}.
% % However, parameter-efficient updates remain susceptible to catastrophic forgetting in sequential training~\cite{mccloskey1989catastrophic}, phenomenonthe deterioration for past-task  measured as increased test risk on earlier task distributions after learning new tasks.
% However, parameter-efficient updates remain susceptible to catastrophic forgetting in sequential training~\cite{mccloskey1989catastrophic}, where learning new tasks increases the test risk on earlier task.
% % as updates for new data can interfere with previously acquired knowledge.
% Continual learning (CL) provides a principled framework for mitigating forgetting~\cite{doi:10.1073/pnas.1611835114,rebuffi2017icarl,de2021continual}, and its combination with parameter-efficient adaptation yields parameter-efficient continual learning (PECL), which enables continual adaptation of pre-trained models under constrained parameter budgets~\cite{dingParameterefficientFinetuningLargescale2023,Gao_2023_ICCV}.
% Continual learning (CL) provides a principled framework for mitigating this interference~\cite{doi:10.1073/pnas.1611835114,rebuffi2017icarl, de2021continual}. Combining parameter-efficient updates with continual learning objectives leads to parameter-efficient continual learning (PECL), which enables continual adaptation of pre-trained models while mitigating forgetting under constrained parameter budgets~\cite{dingParameterefficientFinetuningLargescale2023,Gao_2023_ICCV}.


% Better to move to low-rank already here
% Another parallel strategy for mitigating forgetting is flat-minima optimization. Originally, flat minima in the loss landscape were shown to promote generalization and robustness~\citep{FlatMinima1997,wu2017understandinggeneralizationdeeplearning}. Sharpness-aware minimization (SAM))~\cite{foretSharpnessAwareMinimizationEfficiently2021} and related methods~\citep{pmlr-v151-bisla22a,Zhang_GAM,NEURIPS2024_C-Flat} operationalize this principle by seeking parameters that maintain low loss under small perturbations. 

% Beyond improved generalization, flat minima have also been argued to reduce catastrophic forgetting as flatter solutions are less sensitive to parameter perturbations. Prior work derives that loss landscape flatness helps in reducing forgetting~\cite{NEURIPS2020_518a38cc,JMLRAnEmpiricalInvestigatio,shiOvercomingCatastrophicForgetting2021,liRevisitingCatastrophicForgetting2024}.
% flat-minima
% , adapter-induced
While sharpness-aware optimization has proven effective in the context of full fine-tuning, its behavior under parameter-efficient updates, as considered in PECL, remains underexplored. PECL constrains adaptation to a task-permissible subspace, whereas standard sharpness-aware objectives typically rely on perturbations in the full parameter space. This creates a fundamental tension in perturbation design as shown in Fig.~\ref{fig:scope_type1}: full-parameter perturbations penalize sensitivity along directions that the subsequent LoRA updates cannot realize, while subspace-restricted perturbations can only regularize the landscape within the feasible update directions. 
% full-parameter perturbations perturb directions that subsequent LoRA updates cannot compensate for, 
% while subspace perturbations enforce flatness only within the permissible subspace. 
The key question is whether reducing sharpness
% suppressing curvature 
within the task-permissible subspace
% this task permissible subspace 
is sufficient to secure robust generalization across tasks, or indeed even if it induces any meaningful control over 
the sharpness of the full model.
% curvature
% In this work we investigate the relationship between curvature in task-permissible subspace and curvature in the full parameter space. 
In this work, we investigate how sharpness restricted to the task-permissible subspace relates to sharpness in the full parameter space, and what this implies for continual learning performance.


% full space sharpness need inject perturbation into frozen directions that PECL is designed not to modify, while adapter-restricted perturbations may only enforce
% flatness locally within the permissible subspace. However, it is unclear whether minimizing loss-landscape curvature within a restricted task-permissible subspace is sufficient to ensure robust generalization across tasks or if it induces any meaningful control over the curvature of the full model. In this work, we investigate the relationship between curvature in LoRA update subspaces and curvature in the full parameter space. 



% This raises the question of whether low-rank parameter updates, by severely restricting the optimization subspace, inherently preclude sharp minima, and thus render flat-minima optimization ineffective in parameter-efficient regimes. In this work, we investigate the relationship between curvature in LoRA update subspaces and curvature in the full parameter space. 


% This creates a fundamental tension in perturbation design shown in Fig.~\ref{fig:scope_type1}: full space sharpness need inject perturbation into frozen directions that PECL is designed not to modify. 



% while adapter-restricted perturbations may only enforce flatness locally within the permissible subspace. This mismatch motivates a PECL-specific investigation of perturbation scope and geometry, and their implications for continual generalization.




% In particular, it is unclear whether minimizing loss-landscape curvature within a restricted task-permissible subspace as shown in Fig.~\ref{fig:scope_type1} is sufficient to ensure robust generalization across tasks or if it induces any meaningful control over the curvature of the full model.




% Several works have suggested that over-parameterization is a key factor of sharp minima in deep networks~\cite{shin2025criticalinfluenceoverparameterizationsharpnessaware}. This raises the question of whether low-rank parameter updates, by severely restricting the optimization subspace, inherently preclude sharp minima, and thus render flat-minima optimization ineffective in parameter-efficient regimes. In this work, we investigate the relationship between curvature in LoRA update subspaces and curvature in the full parameter space. 

Despite its practical relevance, the connection between PECL-specific constraints and generalization has not been theoretically established. The PAC-Bayesian framework~\cite{mcallesterPACBayesianStochasticModel2003,pentinaPACBayesianBoundLifelong2014,germainPACBayesianTheoryMeets2016,NEURIPS2024_xinrui} provides a principled route to generalization guarantees via posterior complexity. However, existing studies rarely address continual learning and typically do not incorporate the curvature of the model. 
% In particular, the interaction between adapter subspace sharpness and generalization across tasks remains poorly understood.
In this paper, we develop a sequential hierarchical PAC-Bayes analysis tailored to PECL. Unlike classical hierarchical PAC-Bayes bounds that assume a fixed hyperprior~\cite{pentinaPACBayesianBoundLifelong2014}, our formulation allows the task-level hyperposterior to evolve across tasks and introduces an explicit drift penalty that quantifies cross-task interference. 

We further adapt our framework to LoRA-based PECL and show that both the KL complexity term and the sharpness-dependent empirical term reduce to quantities defined on the task-permissible subspace. This result provides a principled theoretical criterion: controlling sharpness within the subspace is sufficient when the dominant task-conditioned curvature is concentrated in the task-permissible subspace.
% for when flatness control restricted to adapter parameters is sufficient.
% We look at the effectiveness of two flat minima optimization methods 
Motivated by this criterion, we compare representative sharpness-aware objectives and show that first-order sharpness minimization more effectively suppresses dominant curvature directions within the task-permissible subspace.
% algorithms reduce dominant curvature directions more significantly for LoRA updates. 
% than for full model updates
% We hypothesize that this is due ....  
In extensive experiments we demonstrate that the reduced subspace sharpness in task-permissible subspace translates to considerable performance gains for continual learning.

% improved flatness 

% \JW{Surprisingly, we find that the dominant curvature directions of the full model substantially align with the sparse update directions, indicating that flat-minima optimization can remain effective even under highly constrained parameter updates. 
% In addition, 






% tttttttttttttttt

%and Adapter\cite{pmlr-v97-houlsby19a} 
% ,He_2025_CVPR,
% wang2023orthogonal,yang2024parametercollisionhinderingcontinual
%,tsuzukuNormalizedFlatMinima2020,NEURIPS2021_995f5e03
% ~\cite{NIPS1994_01882513,wu2017understandinggeneralizationdeeplearning,jiangFantasticGeneralizationMeasures2019
% Flat minima are widely associated with improved generalization and robustness, motivating Sharpness-aware minimization (SAM)~\cite{foretSharpnessAwareMinimizationEfficiently2021} and related methods to control sharpness under local perturbations in the full parameter space. Beyond single task setting, such flatness control has been argued to mitigate forgetting in CL~\cite{NEURIPS2020_518a38cc,JMLRAnEmpiricalInvestigatio,shiOvercomingCatastrophicForgetting2021,liRevisitingCatastrophicForgetting2024}.


% In PECL, however, feasible adaptation is restricted to task-admissible update directions induced by the trainable adapters, 
% In PECL, however, the update trajectory is confined to the task-permissible subspace, while standard sharpness objectives are typically defined using perturbations in the full parameter space.
% % definitions are stated over full-parameter neighborhoods. 
% This creates a scope mismatch between the perturbation geometry and the feasible update dynamics as shown in Fig.~\ref{fig:scope_type1}.
% Consequently, a central question in PECL is whether controlling sharpness only within the task-admissible update subspace is sufficient to ensure robust generalization across tasks.
% The scope mismatch naturally motivates two PECL-specific design questions: \emph{where} perturbations should be applied and \emph{how} they should be instantiated.
% Even in the single-task setting, prior work remains divided on perturbation scope, advocating either full-parameter perturbations or adapter-restricted perturbations~\citep{NEURIPS2024_4eb2c0ad,li2024flatlora}.

%, as illustrated in Figure~\ref{fig:scope_type1}
% {
% \begin{figure*}[htbp]
%     \centering
%     \makebox[\linewidth]{%
%     \begin{subfigure}[t]{0.42\linewidth}
%         \raggedright
% \includegraphics[height=3.8cm, keepaspectratio]{figures_TRML/shiyi_yanse.pdf}
%         \subcaption{Scope mismatch: full vs.\ adapter subspace}
%     \end{subfigure}
%     \hfill
%   % \hspace{4pt}
%     \begin{subfigure}[t]{0.55\linewidth}
%         \raggedleft
% \includegraphics[height=3.8cm, keepaspectratio]{figures_TRML/lossland_curv1d_combine_loss_title2.pdf}
%        \subcaption{Type effects: Different subspace perturbations shape full-model flatness}
%         \label{fig:perturb_type_landscape}
%     \end{subfigure}
%   }
%     % \vspace{-6pt}
%     \caption{\textbf{Where and How to Perturb in LoRA-Based CL}.
%     \textbf{(a) Where (scope)}:  Parameter-efficient Continual Learning continual learning restricts adaptation to a task-admissible set of update directions.
% Full-space perturbations explore ambient directions that may lie outside the feasible update dynamics, whereas adapter subspace perturbations confine the scope to update consistent directions. \textbf{(b) How (type)}: Different perturbation types shape the geometry of updates within the subspace, leading to different solutions and different full-parameter loss landscapes.}
%     \label{fig:scope_type1}
% \end{figure*}
% }





%     \begin{subfigure}[t]{0.55\linewidth}
%         \raggedleft
% \includegraphics[height=3.8cm, keepaspectratio]{figures_TRML/lossland_curv1d_combine_loss_title2.pdf}
%        \subcaption{Type effects: Different subspace perturbations shape full-model flatness.\textbf{(b) How (type)}: Different perturbation types shape the geometry of updates within the subspace, leading to different solutions and different full-parameter loss landscapes.}
%         \label{fig:perturb_type_landscape}
%     \end{subfigure}
  


%PAC-Bayesiansupervisedclassification2007,,lotfiPACBayesCompressionBounds2022,



% {
% \begin{figure}[tbp]
%     \centering
% \includegraphics[width=\linewidth]{figures_TRML/duibi2.pdf}
%     % \vspace{-6pt}
%     \caption{\textbf{Where and How to Perturb in LoRA-Based PECL}: \emph{where} to apply perturbations (scope: full- or low-rank) and \emph{how} to perturb (type: random perturbation $\mathrm{RS}^{(0)}$, adversarial zeroth-order perturbation $\mathrm{AS}^{(0)}$ and adversarial first-order perturbation $\mathrm{AS}^{(1)}$).}
%     \label{fig:scope_type1}
% \end{figure}
% }



% In this work, we revisit how sharpness could be defined and controlled in LoRA-based PECL.
% % To establishing a sequential hierarchical PAC-Bayesian framework for continual learning and specializing it to PECL. 
% To provide a theoretical explanation for this scope mismatch and its impact on continual generalization, we establish a sequential hierarchical PAC-Bayesian framework for continual learning and specialize it to PECL.
% The bound reveals that task-level priors and posteriors are supported on the task-permissible adapter subspace, so both the sharpness-dependent empirical term and the cumulative complexity term are determined by quantities within this permissible update geometry. 


% This yields a principled criterion for perturbation design in PECL: where to perturb and how to perturb.  We demonstrate that perturbations should be confined to the task-permissible adapter subspace, rather than the full parameter space. Within this permissible scope, the perturbation geometry should be chosen to emphasize the task-conditioned curvature directions that remain reachable by subsequent adapter updates, since different geometries induce systematically different continual outcomes. The experiment  study perturbation scope and type across representative LoRA organizations, with matching trends in multiple curvature diagnostics derived from both Hessian- and Fisher-based estimates.


% : robust generalization across tasks should be controlled within the same update-consistent directions, while imposing robustness constraints in the full parameter space is neither required by the bound nor aligned with the feasible adaptation dynamics. Guided by this criterion, we study perturbation scope and type across representative LoRA organizations, and find that methods that more effectively suppress sharpness within the permissible adapter directions consistently achieve better continual accuracy and stronger robustness under distribution shifts, with matching trends in multiple curvature diagnostics derived from both Hessian- and Fisher-based estimates.





% Under this specialization, task-level priors and posteriors are supported on the task-admissible adapter update subspace. As a result, both the sharpness-dependent empirical contribution and the continual complexity penalty are governed solely by quantities computed within this admissible update geometry.
% % $\mathcal W_{\Delta,t}$, which forces both the sharpness-dependent empirical contribution and the continual complexity penalty to depend only on subspace-defined quantities.
% This yields a principled criterion for perturbation design in PECL: robust generalization across tasks should be controlled in the update space , while enforcing robustness  in the full parameter is neither required by the bound nor compatible with the feasible update dynamics. Guided by this theory, we systematically study perturbation scope (where to perturb) and type (how to perturb) across representative LoRA organizations, and empirically show that sharpness minimization within update subspace translates into improved continual accuracy and robustness.

% We answer this question through a sequential hierarchical PAC-Bayesian framework for continual learning, and then specialize it to frozen-backbone PECL. Under this specialization, task-level priors and posteriors are supported on the task-admissible adapter update subspace, which forces both the sharpness-dependent empirical term and the continual complexity term to reduce to subspace-defined quantities. In this work, we develop a PAC-Bayesian framework for continual learning. We derive a sequential hierarchical PAC-Bayes framework and specialize it to frozen-backbone PECL.
% Under this specialization, task-level priors and posteriors are supported on the task-admissible adapter update subspace, and both the sharpness-dependent empirical term and the continual complexity term reduce to quantities defined on this subspace. We further conduct a systematic empirical study across perturbation strategies and LoRA organizations, providing consistent evidence that adapter-restricted sharpness regularization improves continual performance and robustnes.



Our main contributions can be summarized as:
% Our analysis and experiments support three conclusions.
\begin{itemize}
\vspace{-8pt}
\item
%study the impact of flat-minima optimization for PECL
We revisit sharpness-aware learning for PECL and analyze how sharpness in the restricted task-permissible LoRA subspace affects the curvature of the full model and the generalization over tasks.
% We derive a sequential hierarchical PAC-Bayesian bound for continual learning, where a transition-induced complexity term captures cross-task interference. When specialized to frozen-backbone PECL, both the sharpness-dependent empirical contribution and the cumulative complexity penalty depend only on quantities defined within the task-admissible adapter update directions. This formalizes why subspace sharpness, rather than full-parameter sharpness, is the bound-relevant notion for PECL generalization across tasks.
% We investigate sharpness and derive a sequential hierarchical PAC-Bayesian bound for continual learning. Under PECL, the bound depends only on sharpness and complexity quantities defined within task-permissible adapter update directions, identifying subspace sharpness as the bound-relevant notion for PECL generalization across tasks.
% The sequential hierarchical PAC-Bayesian bound for PECL shows that both the sharpness-dependent empirical term and the cumulative complexity term are governed by quantities supported on the task-permissible adapter update directions.
% We derive a sequential PAC-Bayesian bound for continual learning with a hierarchical prior, yielding a transition-induced KL drift term that quantifies cross-task interference.
% We specialize the bound to frozen-backbone PECL in which both the perturbation term and the complexity term are supported on the trainable adapter subspace.
\vspace{-2pt}
% \textbf{Principle}:
\item
% We derive a sequential hierarchical PAC-Bayesian bound for continual learning yielding a transition-induced KL drift term that quantifies cross-task interference. We show that in PECL, the bound depends on sharpness and complexity quantities defined within the task-permissible subspace.
We derive a sequential hierarchical PAC-Bayesian bound and show that under PECL both the bound-relevant sharpness penalty and the cumulative complexity are supported on the task-permissible subspace.
% identifying subspace sharpness as the bound-relevant notion for PECL generalization across tasks.
% The bound yields a task-conditional perturbation rule: \emph{where} to perturb is within the task-permissible adapter update directions, and \emph{how} to perturb is to use increasingly curvature-aligned perturbation types within that permissible geometry.
% Across LoRA organizations, applying perturbations within the task-permissible adapter scope consistently improves continual accuracy and robustness compared with full-parameter perturbations.
% The bound makes perturbation design in PECL explicitly {task-conditional} and decomposes it into two coupled choices that directly address the central question. Where: perturbations should be restricted to $\mathcal W_{\Delta,t}$, since perturbing outside this subspace enforces robustness along unreachable directions and can introduce unnecessary loss elevation. How: within $\mathcal W_{\Delta,t}$, perturbation geometries that concentrate variance along dominant high-curvature directions are expected to bias training toward flatter task-conditioned solutions, improving continual generalization.


% : \emph{where} the bound suggests restricting to the current task's admissible update subspace rather than the full space, and \emph{how} to perturb, where perturbation that align with dominant high-curvature directions are expected to bias the training trajectory toward flatter regions.
\vspace{-4pt}
%\textbf{Evidence}:
\item
% We demonstrate using different flat-minima optimization methods and LoRA-based CL methods that minimizing sharpness in the restricted subspace achieves better performance over conperturbation in full model space.
% We demonstrate that sharpness control restricted to the task-permissible LoRA update subspace is consistently associated with stronger continual performance and reduced forgetting across LoRA-based CL methods.
We demonstrate that continual performance differences across perturbation designs are consistent with curvature localization: methods that more strongly suppress subspace sharpness along dominant curvature directions exhibit higher stability and less forgetting across tasks.
% study perturbation design under the PECL constraint by contrasting subspace-restricted versus full-parameter perturbation scopes and evaluating multiple sharpness-aware perturbation types within the LoRA update subspace. Curvature diagnostics further show that lower subspace sharpness is associated with higher continual accuracy and less forgetting across tasks.

% demonstrate using different flat-minima optimization methods and LoRA-based CL methods that minimizing sharpness in the restricted subspace achieves better performance over conperturbation in full model space.

% Under a fixed task-permissible scope, different perturbation geometries lead to differences in perturbation-aware loss and curvature diagnostics, and these differences track the observed performance gaps along the task sequence.
% Across LoRA organizations and perturbation types, adapter-restricted sharpness control consistently improves continual accuracy and robustness under distribution shifts.
% Across LoRA organizations and perturbation types, we observe a consistent ordering: methods that more strongly suppress adapter-subspace sharpness yield better continual accuracy and robustness. 


% We systematically evaluate perturbation scope/type across LoRA organizations and robustness benchmarks, confirming that adapter-restricted sharpness regularization improves continual performance and robustness.
\end{itemize}

\section{Related Work}

\minisection{LoRA-Based CL.}
PECL has emerged as a practical solution for incrementally adapting large pre-trained models
% , with LoRA becoming a representative approach due to its efficiency and compatibility with frozen backbones
~\cite{pmlr-v97-houlsby19a,Wang_2022_CVPR,dingParameterefficientFinetuningLargescale2023,xu2023parameterefficientfinetuningmethodspretrained}. By introducing trainable low-rank matrices, LoRA enables continual adaptation with minimal overhead, and has demonstrated strong performance~\cite{hu_lora_2021}.
However, LoRA-based methods remain vulnerable to forgetting and task interferences~\cite{Gao_2023_ICCV,10.1145/3730436.3730456,yang2024parametercollisionhinderingcontinual}. 
 % Gated variants further select or combine adapters with task-dependent weights~\citep{Gao_2023_ICCV,10.1145/3730436.3730456,zhao-etal-2024-sapt,cheng2025exploiting,wu2025sdlorascalabledecoupledlowrank}.
Recent variants such as SDLoRA~\cite{wu2025sdlorascalabledecoupledlowrank} and OLoRA~\cite{NEURIPS2020_70d85f35,wang2023orthogonal} aim to alleviate these issues via adapter modularity and orthogonal constraints. 
Although such designs yield empirical improvements, they lack a principled understanding of how adapter structure affects forgetting and generalization.
% In particular, the interaction between adapter subspace design and solution stability remains insufficiently explored.

\minisection{Flat Minima.}
Prior work has explored flat minima in CL through sharpness-aware optimization~\cite{pmlr-v151-bisla22a,Zhang_GAM,NEURIPS2024_C-Flat,yangRevisitingFlatnessAwareOptimization2025,chen2026sharpnessflatnessdecompositionframework}, mode connectivity~\cite{izmailov2019averagingweightsleadswider,mirzadeh2020linearmodeconnectivitymultitask}, and improved representation learning~\cite{hu2022doesselfsupervisedpretrainingperform,ramasesh2022effect}.
These methods proposed for full model training do not account for continual training dynamics in the low-rank subspace~\cite{yangRevisitingFlatnessAwareOptimization2025}. Whether flatness in the low-rank subspace alone can ensure generalization in CL remains unclear, motivating further investigation.


\minisection{PAC-Bayesian Theory.}
The PAC-Bayesian framework provides a principled foundation for analyzing generalization by balancing empirical risk and posterior complexity~\cite{mcallesterPACBayesianModelAveraging1999,mcallesterPACBayesianStochasticModel2003,pentinaPACBayesianBoundLifelong2014,germainPACBayesianTheoryMeets2016,lotfiPACBayesCompressionBounds2022}.
Online PAC-Bayes~\citep{NEURIPS2022_onlineBayes,NEURIPS2024_xinrui} allows time-varying, data-dependent priors, but it is formulated at the level of individual examples rather than tasks and does not explicitly model task-level inheritance in continual learning.
Hierarchical PAC-Bayes bounds for meta learning~\citep{pentinaPACBayesianBoundLifelong2014,JMLR:v24:22-1254} introduce hyperposteriors over task priors, but they assume a single stationary hyperposterior and therefore do not model the sequential, history-dependent evolution of the inductive bias that underlies forgetting in continual learning.
% In continual learning, the inductive bias must adapt to the observed task history, and this non-stationarity is particularly pronounced in LoRA-based PECL, where the update scope and the initialization of incremnental adapters could change with training.

\section{Preliminaries}

\subsection{LoRA-Based CL}
\label{sec:pecl-setup}

Let $\mathcal X$ be the input space, $\mathcal Y$ the label space, and $\mathcal Z = \mathcal X \times \mathcal Y$ the data space.
A model with parameters $W \in \mathbb R^d$ induces a predictor $f_W : \mathcal X \to \mathcal Y$.
We consider a bounded loss $\ell : \mathcal H \times \mathcal Z \to [0,1]$ and write $\ell(W,z)$ for $\ell(f_W,z)$. 
Given a distribution $\mathcal D$ over $\mathcal Z$ and a training set $S=\{z_j\}_{j=1}^m$ drawn i.i.d.\ from $\mathcal D$, the population and empirical risks are
$
L_{\mathcal D}(W)
=
\mathbb E_{z\sim\mathcal D}\ell(W,z),
\hat L_{S}(W)
=
\frac{1}{m}\sum_{j=1}^{m}\ell(W,z_{j}).
$



% LoRA-based PECL adapts a frozen backbone by learning low-rank updates across tasks. A key design axis is how these updates are organized in parameter space, since the resulting sharing pattern determines the degree of knowledge transfer and cross-task interference. We focus on \emph{shared-subspace} PECL, where tasks interact through a common adapter mechanism, and summarize three representative organizations.

% \textbf{Fully shared.}
% SeqLoRA reuses a single LoRA pair $(A,B)$ for all tasks and updates it sequentially on the fixed backbone $W_0$. This maximizes parameter sharing but also couples tasks most strongly, making interference accumulate within the same low-rank directions.

% \textbf{Incremental shared.}
% Incremental designs expand the adapter set over time by introducing a new pair $(A_t,B_t)$ at task $t$ while keeping $W_0$ and earlier adapters fixed, and then composing multiple adapters at inference.

% \textbf{Constraint incremental shared.}
% Constraint-based methods regularize the geometry of the adapter space to reduce conflicts based incremental designs. OLoRA~\citep{wang2023orthogonal} encourages orthogonality between task-specific directions to reduce parameter overlap, while gradient-based approaches align or project updates to avoid previously occupied directions~\citep{yang2024parametercollisionhinderingcontinual}.

In continual learning with \(T\) tasks indexed by \(t=1,\dots,T\), each task provides a sample \(S_t=\{z_{tj}\}_{j=1}^{m_t}\) drawn i.i.d. from a task distribution \(\mathcal{D}_t\). LoRA-based PECL adapts a frozen backbone by learning task-wise low-rank weight increments. Throughout, we treat the LoRA {weight increment} as the central object and use \emph{task-permissible subspace} to refer to the feasible update directions of this increment. A key design axis is how task-permissible subspaces are organized across tasks within a shared parameter space. We focus on such methods and distinguish three categories.

% In continual learning with \(T\) tasks indexed by \(t=1,\dots,T\), each task provides a sample \(S_t =\{z_{tj}\}_{j=1}^{m_t}\) drawn i.i.d. from a task distribution \(\mathcal{D}_t\). LoRA-based PECL adapts a frozen backbone by learning task-wise low-rank updates. Throughout, we treat the LoRA \emph{weight increment} as the central object and use \emph{task-permissible subspace} to refer to the feasible update directions of this increment. A key design axis is how these task-permissible subspace are organized. 
%  We focus on methods that share a common parameter space and distinguish three categories.
% LoRA-based PECL adapts a frozen backbone by learning low-rank updates across tasks and A key design axis is how these updates are organized in parameter space.  We focus on methods with a shared parameter space and
% distinguish three categories.

% In the context of CL, we focus on three representative LoRA-
% based training strategies:
\begin{itemize}
\vspace{-1em}
\item \textbf{Fully shared (SeqLoRA)} : A single LoRA pair $(A, B)$ is trained sequentially and shared across tasks while keeping the backbone weights $W_0$ frozen. The final model after $N$ tasks is given by $W_0 + AB$.
\vspace{-0.5em}
\item \textbf{Incremental shared(IncLoRA)} : A new LoRA pair $(A_t, B_t)$ is incrementally added for each task $t$, and only the current pair is updated while all previous modules and $W_0$ are fixed. The resulting model is $W_0 + \sum_{t=1}^N A_t B_t$.
\vspace{-0.25em}
\item \textbf{Constraint shared(OLoRA)} : Extending IncLoRA, OLoRA~\cite{wang2023orthogonal} introduces an orthogonality constraint across tasks to encourage task-permissible subspace disentanglement. Specifically, during task $t$, a penalty $\mathcal{L}_{\text{orth}}^{(t)} = \sum_{i=1}^{t-1} \|A^{(i)\top} A^{(t)}\|_F^2$ is applied, where $A^{(i)}$ denotes the LoRA-A matrix from task $i$.
\end{itemize}

% For completeness, some methods allocate fully disjoint adapters per task and activate them conditionally by task identification at inference. Such isolated designs are not the focus of our analysis.

\subsection{Classical PAC-Bayesian Bound}



The PAC-Bayesian viewpoint~\cite{mcallesterPACBayesianModelAveraging1999,mcallesterPACBayesianStochasticModel2003} place probability measures over the hypothesis space $\mathcal{H}$ rather than focusing on a single function: starting from a prior $P\in\mathcal M(\mathcal H)$, the learner observes data $S$ and outputs a posterior $Q(\cdot\mid P,S)\in\mathcal M(\mathcal H)$. The corresponding posterior-averaged population and empirical risks  are $
L_{\mathcal D}(Q)
=
\mathbb E_{W \sim Q(\cdot \mid P,S)}\mathbb E_{z\sim\mathcal D}\ell(W,z)
$ and
$
\hat L_{S}(Q_t)
=
\mathbb E_{W \sim Q(\cdot \mid P,S)}\frac{1}{m}\sum_{j=1}^{m_t}\ell(W,z_{j}).
$
A classical PAC--Bayes bound~\cite{mcallesterPACBayesianStochasticModel2003} for a single task under a fixed distribution is stated as follows.
\begin{theorem}[PAC-Bayes Generalization Bound for single task]
\label{theorem:classical}
For any loss $\ell: \mathcal{F} \times \mathcal{X} \times \mathcal{Y} \to [0,1]$, with probability at least $1-\delta$ over the choice of training set $S$:
{%
% \setlength{\abovedisplayskip}{2pt}
\setlength{\belowdisplayskip}{3pt}
% \setlength{\abovedisplayshortskip}{3pt}
\setlength{\belowdisplayshortskip}{3pt}
    \[\mathbb{E}_{f \sim Q}[L_{\mathcal{D}}^{\ell}(f)] \leq \mathbb{E}_{f \sim Q}[\hat{L}_S^{\ell}(f)] + \sqrt{\frac{KL(Q\|P) + \log \frac{m}{\delta}}{2(m-1)}}\]
}
\end{theorem}

\subsection{Sharpness}
\label{section:sahrpness}
% Flat minima are often associated with improved generalization and robustness in standard supervised learning. 
Motivated by PAC-Bayes analysis, SAM~\citep{foretSharpnessAwareMinimizationEfficiently2021} seeks to reduce the empirical contribution in generalization bounds by controlling the loss elevation in a local neighborhood. In particular, SAM can be interpreted as minimizing an \textbf{adversarial zeroth-order sharpness}~\citep{AWP2020,foretSharpnessAwareMinimizationEfficiently2021} on the training set $S$ within a radius $\rho$:
{%
\[\mathrm{AS}^{(0)}_S(W,\rho) =\max_{ \| \epsilon \|\leq \rho } \Big[ \hat{L}_{S}(W+{\epsilon})-\hat{L}_{S}(W)\ \Big].\]
}
% \vspace{-5mm}
A complementary formulation considers randomized perturbations and minimizes the expected loss elevation under Gaussian noise, referred to as \textbf{random zeroth-order sharpness}~\citep{liRevisitingRandomWeight2024,jiang2019fantasticgeneralizationmeasures}:
{%
\[
\mathrm{RS}_S^{(0)}({W}, \Sigma)
 =
 \mathbb{E}_{\varepsilon \sim \mathcal{N}(0,\Sigma)}
 \big[
 \hat{L}_{S}({W}+\varepsilon) - \hat{L}_{S}({W})
 \big],
 \]}Typical choices for $\mathcal{N}(0,\Sigma)$ include isotropic Gaussians and variants such as filter-wise Gaussians~\citep{pmlr-v151-bisla22a}. Moreover, SAM can be interpreted as optimizing the mean loss under a fixed perturbation variance: for a suitable correspondence between radius $\rho$ and covariance $\Sigma$, 
 % $\mathrm{AS}^{(0)}_S(W,\rho)$ provides an upper bound on $\mathrm{RS}_S^{(0)}({W}, \Sigma)$~\cite{mollenhoffSAMOptimalRelaxation2022}.
 the adversarial zeroth-order sharpness provides an upper bound on the random zeroth-order sharpness~\cite{mollenhoffSAMOptimalRelaxation2022}.

% Beyond loss elevation, sharpness can be measured through the neighborhood gradient norm.
The \textbf{adversarial first-order sharpness} proposed in GAM~\citep{Zhang_GAM} is defined as the maximum gradient norm within a radius $\rho$ on the training set:
{%
\[\mathrm{AS}^{(1)}_S(W,\rho) = \max_{ \| \epsilon \|\leq \rho } \| \nabla_{W} \hat{L}_S(W+\epsilon)\|.\]}Compared with zeroth-order sharpness which measures  loss increase, the first-order criterion emphasizes the sharpness of the gradient itself, and is closely tied to dominant curvature directions. For any $W$ and fixed $\rho$, the adversarial first-order sharpness dominates the adversarial zeroth-order sharpness, hence minimizing $\mathrm{AS}^{(1)}_S$ also reduces $\mathrm{AS}^{(0)}_S$, making it a stronger flatness measure~\citep{Zhang_GAM}.








\section{Revisiting Sharpness for PECL}
\label{sec:pac-pecl}

We revisit sharpness from the perspective of PAC-Bayesian Analysis~\cite{mcallesterPACBayesianModelAveraging1999}. Theorem~\ref{theorem:classical} provides a generalization guarantee for a single task under a fixed data distribution and a prior that is fixed before observing the corresponding training set. Continual learning instead proceeds over a task sequence and the learner updates its hypothesis after observing task data. This sequential setting implies that the inductive bias should adapt with the task history, so the task prior at time $t$ must reflect information from tasks $1,\dots,t-1$, which cannot be captured by treating $\{P_t\}$ as externally specified constants. We now discuss how to apply PAC-Bayesian bounds for the continual setting.

\subsection{PAC-Bayesian Bound for Continual Learning}
\label{sec:single_to_cl}

We model transferable inductive bias by treating the task prior as a random object and introducing a hyperposterior $\mathcal P\in\mathcal M(\mathcal M(\mathcal H))$ over task-level priors. The key motivation is that, in continual learning, the prior used at step $t$ is typically produced by a history-dependent rule that progressively incorporates information from earlier tasks, rather than being generated from a single time-invariant hyperdistribution (Multi-task Learning) as shown in Fig.~\ref{fig:showPAC3}. Concretely, common CL mechanisms make the prior sequence nonstationary, including curvature-based regularization that updates the effective prior through Fisher or Hessian summaries~\citep{EWC2017,SI2017} and scope-allocation strategies that modify which parameters remain trainable over time~\citep{progressive2016,HAT2018,Supermasks2020,pmlr-v97-houlsby19a}, which naturally introduces a notion of drift between priors.

{%
\setlength{\intextsep}{8pt}        % [h] 浮动体与正文的上下间距（最关键）
\setlength{\textfloatsep}{6pt}     % [t]/[b] 时浮动体与正文间距（备用）
\setlength{\abovecaptionskip}{2pt} % 图与caption之间（上方）间距
\setlength{\belowcaptionskip}{0pt} % caption下方到正文的间距
\captionsetup{skip=4pt}            % 你已在用：图与caption的距离（更细粒度）

\begin{figure}[t]
\centering
\includegraphics[width=0.95\linewidth]{figures_TRML/multitask5.pdf}
\caption{
Time-varying Hyperposterior in Continual Learning. \textbf{(Left)} In multi-task learning, task priors are drawn from a single time-invariant hyperdistribution, so there is no sequential drift. \textbf{(Right)} In continual learning, the hyperposterior is updated over time, inducing a history-dependent cumulative drift over tasks.
% Time-varying Hyperposterior in Continual Learning. \textbf{Left:} Multi-task uses a single hyperdistribution drift term to sample task priors.
% \textbf{Right:} Continual learning updates the hyperposterior over time, yielding history-dependent cumulate hyperdistribution drift.
}
\vspace{-10pt}
\label{fig:showPAC3}
\end{figure}
}% 结束局部分组

Let $\mathcal F_{t-1}=\sigma(S_1,\dots,S_{t-1})$ denote the information available before observing task $t$. We model continual learning as a hyperposterior process $\{\mathcal P_t\}_{t=0}^T$ where $\mathcal P_{t-1}$ is $\mathcal F_{t-1}$ measurable, which yields a pathwise bound with an explicit hyperposterior drift penalty that quantifies the information budget of sequential belief updates. We define the hierarchical test and empirical risks
$
\mathcal L_t
=
\mathbb E_{P_t \sim \mathcal{P}_{t-1}}
\,
\mathbb E_{W_t \sim Q_t(\cdot \mid P_t,S_t)}
\big[
L_{\mathcal D_t}(W_t)
\big]
$ and
$
\hat{\mathcal L}_t
=
\mathbb E_{P_t \sim \mathcal{P}_{t-1}}
\,
\mathbb E_{W_t \sim Q_t(\cdot \mid P_t,S_t)}
\big[
\hat L_{S_t}(W_t)
\big].
$ We propose the following PAC-Bayes bound for continual learning:

% \vspace{0.5em}
\begin{theorem}[PAC Bayes with time varying hyperposterior]
\label{thm:pathwise_pac}
For any $\delta\in(0,1]$, with probability at least $1-\delta$ over $S_{1:T}$,
{%
% \setlength{\abovedisplayskip}{2pt}
\setlength{\belowdisplayskip}{3pt}
% \setlength{\abovedisplayshortskip}{3pt}
\setlength{\belowdisplayshortskip}{3pt}
\begin{equation*}
\resizebox{\linewidth}{!}{$
\begin{aligned} 
&\quad\frac{1}{T} \sum_{t=1}^T \mathcal L_t
\le  
\frac{1}{T} \sum_{t=1}^T \hat{\mathcal L}_t  + \\
&
\sqrt{
\frac{
\displaystyle
\sum_{t=1}^T
\mathbb E_{P_t\sim\mathcal P_{t-1}}
\mathrm{KL}\bigl(Q_t \,\Vert\, P_t\bigr)
+
\sum_{t=1}^T
\mathrm{KL}\bigl(\mathcal P_t\Vert\mathcal P_{t-1}\bigr)
+
\log\frac{1}{\delta}}
{2\displaystyle\sum_{t=1}^T (m_t - 1)}
}.
\end{aligned}
$}
\end{equation*}
}
\end{theorem}

Theorem~\ref{thm:pathwise_pac} bounds the trajectory-averaged hierarchical test risk by the corresponding trajectory-averaged empirical risk plus a complexity penalty, aligning with standard continual learning evaluation over a task sequence.
The complexity decomposes into a within task adaptation term $\mathbb E_{P_t\sim\mathcal P_{t-1}}\mathrm{KL}(Q_t\Vert P_t)$ and a hyperdistribution drift term $\mathrm{KL}(\mathcal P_t\Vert\mathcal P_{t-1})$. The former quantifies how much the learner must deviate from the task prior to fit $S_t$, while the latter quantifies how the task prior generation mechanism evolves as new tasks are incorporated.
This drift is intrinsic to continual learning because the prior used at step $t$ is typically produced by a history-dependent rule rather than fixed externally,
for example by specifying the permissible update scope~\citep{HAT2018,Supermasks2020,hu_lora_2021}, or by shaping the prior geometry using curvature summaries as regularization~\citep{EWC2017,SI2017}.
Theorem~\ref{thm:pathwise_pac} also contains the classical exchangeable-task hierarchical PAC-Bayes bound~\citep{pentinaPACBayesianBoundLifelong2014} as a special case: if the hyperposterior is not a sequence change then the result reduces to a single shared hyperdistribution over task priors. We refer to Appendix~\ref{proof:theorem all} for the detailed proof of Theorem~\ref{thm:pathwise_pac}.

\subsection{Moving from Full-Finetuning to LoRA}
\label{sec:kl-decomp-pecl}

% In PECL, we represent the task-$t$ parameters in an affine form
% \[
% W_t = W^{\mathrm{froz}}_t + \Delta W_t ,
% \]
% where $W^{\mathrm{froz}}_t\in\mathbb R^d$ collects all parameters that are frozen while learning task $t$
% (with $d=\dim(\mathrm{vec}(W))$ under the chosen vectorization), and $\Delta W_t$ denotes the trainable
% increment induced by the current-task LoRA coordinates.
% Accordingly, the feasible learning dynamics are restricted to a task-dependent \emph{increment-direction subspace}:
% \[
% \Delta W_t \in \mathcal W_{\Delta,t}\subseteq \mathbb R^d .
% \]
% Here $\mathcal W_{\Delta,t}$ denotes the task-admissible linear space of \emph{increment directions} in a
% neighborhood of the learned solution, which provides the geometric support for the PAC--Bayesian
% distributions and is instantiated from the LoRA implementation via the local Jacobian construction
% in Appendix~\ref{app:lora_subspace_bridge}.
\begin{figure}[t]
\centering
% \begin{minipage}{1\linewidth}
\centering
\includegraphics[width=0.9\linewidth]{figures_TRML/KLshiyi1.pdf}
% {figures_TRML/KLshiyitu1.pdf}
% \caption{Illustration of the KL divergence decomposition and coordinate equivalence. 
%     (1) \textbf{KL Decomposition}: In PECL, the overall model complexity, measured by $\mathrm{KL}(Q^{\mathrm{full}} \| P^{\mathrm{full}})$, is solely determined by the distribution of the updated parameters in the adapter subspace $\mathcal{W}_\Delta$, formalized as $\mathrm{KL}(Q^{\Delta} \| P^{\Delta})$, while the frozen backbone contributes no divergence. 
%     (2) \textbf{Two equivalent views}: The optimization in the low-rank factor coordinates $(A, B)$ is theoretically isomorphic to operating directly in the update subspace $\mathcal{W}_\Delta$. This equivalence preserves geometric properties such as flatness and curvature, which allows sharpness-aware optimization to be confined to the factorized LoRA parameters.}
% \vspace{-4pt}
%$T_t(\Delta)=W_t^{\mathrm{froz}}+\Delta$
%: Full-Model KL vs Adapter-Space KL.
\caption{\textbf{KL Decomposition}
 The posterior/prior defined on the task-permissible subspaces $(Q_t^\Delta,P_t^\Delta)$ induce full-model distributions $(Q_t,P_t)$ with $\mathrm{KL}(Q_t\Vert P_t)=\mathrm{KL}(Q_t^\Delta\Vert P_t^\Delta)$, so the complexity term depends only on task-permissible update subspace.}
\vspace{-10pt}
% \caption{ Moving from Full-Finetuning to LoRA: Full-Model KL vs Adapter-Space KL.
% $\mathrm{KL}(Q^{\mathrm{full}}\|P^{\mathrm{full}})=\mathrm{KL}(Q^{\Delta}\|P^{\Delta})$, so the complexity depends only on the adapter subspace $\mathcal W_{\Delta}$.
% LoRA factors $(A,B)$ parameterize the same update subspace, allowing sharpness control to be applied in factor space.
% }
\label{fig: KL_decompose}
% \end{minipage}
\end{figure}
% In PECL, the parameters after learning task $t$ can be written as
% \[
% W_t = W^{\mathrm{froz}}_t + \Delta W_t,
% \]
% where $W^{\mathrm{froz}}_t$ collects all frozen parameters at task $t$, and $\Delta W_t$ denotes the trainable LoRA weight increment for the current task.
In PECL, the task-$t$ parameters can be written as
$
W_t = W^{\mathrm{froz}}_t + \Delta W_t,
$
where $W^{\mathrm{froz}}_t$ refers to
% $W^{\mathrm{froz}}_t\in\mathbb R^d$ collects
%from previous task while
all frozen parameters when learning task $t$, and $\Delta W_t$ are the trainable LoRA parameters for task $t$. 
% The permissible update dynamics are thus restricted to a task-dependent incremental subspace
The PECL constraint restricts the increment to a task-permissible subspace,
$
\Delta W_t \in \mathcal W_{\Delta,t},
$
where $\mathcal W_{\Delta,t}$ is a linear subspace of the ambient weight-increment space that locally models the feasible increment directions around the task-$t$ solution (see Appendix~\ref{app:lora_subspace_bridge}). Consequently, the hypothesis class at task $t$ is the affine set $W^{\mathrm{froz}}_t + \mathcal W_{\Delta,t}$, and all task-dependent randomness in our analysis is supported on $\mathcal W_{\Delta,t}$ and then translated to the affine class.


% The update dynamics are restricted to a \emph{task-permissible subspace}
% $
% \Delta W_t \in \mathcal W_{\Delta,t}
% $
% where $\mathcal W_{\Delta,t}$ denotes the set of admissible increment directions around the learned solution  
% % the task-specific space of possible directions around the learned
% solution (formalized in Appendix \ref{app:lora_subspace_bridge}). Consequently, the hypothesis class at task $t$ is the affine
% set $W^{\mathrm{froz}}_t + \mathcal W_{\Delta,t}$ and all task-dependent randomness is supported
% on $\mathcal W_{\Delta,t}$ and then pushed forward to the affine class by translation. 

% in the PAC--Bayes analysis 

We place a Gaussian posterior over increments on the task-permissible subspace,
$
Q_t^\Delta = \mathcal N(\widehat{\Delta W}_t,\Sigma_t),
$
where $\widehat{\Delta W}_t$ is the learned LoRA weight increment and $\Sigma_t$ is a covariance operator supported on $\mathcal W_{\Delta,t}$. Equivalently, sampling $W_t \sim Q_t$ on the affine class $W^{\mathrm{froz}}_t+\mathcal W_{\Delta,t}$ amounts to setting $W_t = W^{\mathrm{froz}}_t + \Delta W_t$ with
% \[
% \Delta W_t = \widehat{\Delta W}_t + \varepsilon,\quad \varepsilon \sim \mathcal N(0,\Sigma_t),\ \ \varepsilon \in \mathcal W_{\Delta,t}.
% \]
$\Delta W_t = \widehat{\Delta W}_t + \varepsilon,\quad \varepsilon \sim \mathcal N(0,\Sigma_t),\ \ \varepsilon \in \mathcal W_{\Delta,t}.$
% {%
% \setlength{\abovedisplayskip}{0.3pt}
% \setlength{\belowdisplayskip}{0.3pt}
% \setlength{\abovedisplayshortskip}{0.3pt}
% \setlength{\belowdisplayshortskip}{0.3pt}
% \begin{equation*}
% \Delta W_t = \widehat{\Delta W}_t + \varepsilon,\quad \varepsilon \sim \mathcal N(0,\Sigma_t),\ \ \varepsilon \in \mathcal W_{\Delta,t}.
% \end{equation*}
% }% \in \mathcal W_{\Delta,t}
% We formalize the LoRA adapter posterior on a Gaussian distribution in the permissible space $\mathcal W_{\Delta,t}$ as
% $
% Q_t^\Delta = \mathcal N(\widehat{\Delta W}_t,\Sigma_t)$ where
% $\widehat{\Delta W}_t$ refers to the learned weight update 
% and $\Sigma_t$ is a covariance around $\widehat{\Delta W}_t$ on $\mathcal W_{\Delta,t}$ .
% Equivalently, sampling $W_t\sim Q_t$ on $W^{\mathrm{froz}}_t+\mathcal W_{\Delta,t}$ amounts to 
% % sampling $ \Delta W_t \sim Q_t^{\Delta}$ \text{ on } $W_{\Delta,t}$ and 
% setting $W_t=W^{\mathrm{froz}}_t+\Delta W_t$ and sampling $ \Delta W_t \sim Q_t^{\Delta}$ where $\Delta W_t = \widehat{\Delta W}_t + \varepsilon$
% and
% $
% \varepsilon \sim \mathcal N(0,\Sigma_t).$
% \text{ on }  \mathcal W_{\Delta,t}.$
% and setting $W=W_{\mathrm{frozen}}+\Delta W$.
We propose the following lemma to formalize the resulting translation-invariance of the KL divergence:
% , and its proof is included in Appendix~\ref{proof:Lemma1}.
\begin{lemma}
\label{lem:kl-decomp-pecl}
Assume $Q_t^\Delta \ll P_t^\Delta$. Then $Q_t \ll P_t$ and
\[
\mathrm{KL}\bigl(Q_t \,\Vert\, P_t\bigr)
=
\mathrm{KL}\bigl(Q_t^\Delta \,\Vert\, P_t^\Delta\bigr).
\]
\end{lemma}
Lemma~\ref{lem:kl-decomp-pecl} shows that, under a frozen backbone, the  KL divergence in the full parameter space is exactly the KL divergence on the task-permissible update subspace. We refer to Appendix~\ref{proof:Lemma1} for the proof.

 
Now, we define the subspace sharpness as the expected loss elevation under perturbations supported on the task-permissible subspace $\mathcal W_{\Delta,t}$ for task $t$. The following Lemma~\ref{lem:subspace-loss-pecl} isolates an additive loss elevation caused by perturbations supported on $\mathcal W_{\Delta,t}$.

\begin{lemma}[subspace sharpness]
\label{lem:subspace-loss-pecl}
For each task $t$,
{%
\setlength{\abovedisplayskip}{1pt}
\setlength{\belowdisplayskip}{2pt}
\setlength{\abovedisplayshortskip}{1pt}
\setlength{\belowdisplayshortskip}{2pt}
\begin{equation*}
\label{eq:subspace-loss-pecl}
\begin{aligned}
&\mathbb E_{W_t\sim Q_t}\bigl[\hat L_{S_t}(W_t)\bigr] \\
&= \mathbb{E}_{\varepsilon\sim\mathcal N(0,\Sigma_t)}
    \big[\hat{L}_{S_t}(W^{\mathrm{froz}}_t+\widehat{\Delta W}_t+\varepsilon)\big] \\
& =\hat L_{S_t}(W^{\mathrm{froz}}_t + \widehat{\Delta W}_t)
+
\mathrm{TS}_t^{\Delta}(\widehat{\Delta W}_t,\Sigma_t).
\end{aligned}
\end{equation*}
}where the subspace sharpness is $
\mathrm{TS}_t^{\Delta}(\widehat{\Delta W}_t,\Sigma_t)
= \mathbb{E}_{\varepsilon\sim\mathcal N(0,\Sigma_t)}
   \big[\hat{L}_{S_t}(W^{\mathrm{froz}}_t+\widehat{\Delta W}_t+\varepsilon)
   - \hat{L}_{S_t}(W^{\mathrm{froz}}_t+\widehat{\Delta W}_t)\big] \text{\,and\, } 
     \varepsilon\sim\mathcal N(0,\Sigma_t) \in  \mathcal W_{\Delta,t}$.
     % \begin{equation*}
     % \begin{aligned}
     % &\mathrm{RS}_t^{\Delta}(\hat{\Delta W}_t,\Sigma_t):= \\
     % &\qquad \mathbb{E}_{\varepsilon\sim\mathcal N(0,\Sigma_t)}
     %    \big[\hat{L}_{S_t}(W_\mathrm {frozen }+\hat{\Delta W}_t+\varepsilon) \\
     % & \qquad- \hat{L}_{S_t}(W_\mathrm {frozen }+\hat{\Delta W}_t)\big] \\
     % &\qquad \varepsilon\sim\mathcal N(0,\Sigma_t) \in \mathcal W_\Delta
     % \end{aligned}
     % \end{equation*}
     \end{lemma}

We refer to Appendix~\ref{proof:lemma2} for the proof.



\subsection{PAC-Bayesian Bound for PECL}

In PECL, Lemma~\ref{lem:kl-decomp-pecl} and Lemma~\ref{lem:subspace-loss-pecl} reformulates the complexity and empirical terms entirely in the task-permissible subspace $\mathcal W_{\Delta,t}$ through the task KL divergence $\mathrm{KL}\bigl(Q_t^\Delta \,\Vert\, P_t^\Delta\bigr)$ and the subspace sharpness $\mathrm{TS}_t^{\Delta}(\widehat{\Delta W}_t,\Sigma_t)$ respectively. 
% See Fig.~\ref{fig: showPAC3} and Fig.~\ref{fig: KL_decompose} for an overview of the subspace reduction and the hierarchical construction.
Substituting these reductions into the bound from Theorem~\ref{thm:pathwise_pac}, we obtain the PAC-Bayesian bound for PECL:
\begin{theorem}
\label{thm:pecl-bound}
Consider a PECL model with frozen backbone $W^{\mathrm{froz}}_t$ and Gaussian posteriors 
$Q_t^\Delta = \mathcal N(\widehat{\Delta W}_t,\Sigma_t)$ supported on $\mathcal W_{\Delta,t}$. 
Let $P_t^\Delta$ be task priors sampled from a hyper posterior $\mathcal P_{t-1}$ over adapter 
distributions. For any $\delta\in(0,1]$, with probability at least $1-\delta$ over $S_{1:T}$,
{%
\setlength{\abovedisplayskip}{2pt}
\setlength{\belowdisplayskip}{3pt}
\setlength{\abovedisplayshortskip}{3pt}
\setlength{\belowdisplayshortskip}{3pt}
\begin{equation*}
\label{eq:pecl-main-bound}
\resizebox{\linewidth}{!}{$
\begin{aligned}
&\frac{1}{T}\sum_{t=1}^T
\mathbb E_{P_t^\Delta\sim\mathcal P_{t-1}}
\mathbb E_{W_t\sim Q_t}
L_{\mathcal D_t}(W_t) \\
&\le
\frac{1}{T}\sum_{t=1}^T
\mathbb E_{P_t^\Delta\sim\mathcal P_{t-1}}
\bigl[
\hat L_{S_t}(W^{\mathrm{froz}}_t + \widehat{\Delta W}_t)
+
\mathrm{TS}_t^{\Delta}(\widehat{\Delta W}_t,\Sigma_t)
\bigr] \\
&+
\sqrt{
\frac{
\displaystyle
\sum_{t=1}^T
\mathbb E_{P_t^\Delta\sim\mathcal P_{t-1}}
\mathrm{KL}\bigl(Q_t^\Delta \,\Vert\, P_t^\Delta\bigr)
+
\sum_{t=1}^T
\mathrm{KL}\bigl(\mathcal P_t\Vert\mathcal P_{t-1}\bigr)
+
\log\frac{1}{\delta}}
{2\displaystyle\sum_{t=1}^T (m_t - 1)}
}.
\end{aligned}
$}
\end{equation*}
}
\end{theorem}
Theorem~\ref{thm:pecl-bound} specializes the Theorem~\ref{thm:pathwise_pac} to PECL by exploiting the fact that all task dependent updates are confined to the permissible task subspace $\mathcal W_{\Delta,t}$. 
% This restriction is not merely technical.
Under a frozen backbone $\mathcal W_{\Delta,t}$ is exactly the set of directions along which the model can be modified by future tasks.
Therefore, in PECL the relevant notion of flatness for generalization is the effective geometry induced on $\mathcal W_{\Delta,t}$, rather than a global property of the full parameter space. We provide the proof in Appendix~\ref{proof:theorem}.

% Since subsequent adaptation is restricted to $\mathcal W_{\Delta,t}$, this term quantifies the sensitivity of the task $t$ solution to admissible variations, and thus how robust the learned solution is with respect to the only perturbations that future tasks can effectively realize.
% The complexity penalty is expressed through the adapter level divergence $\mathrm{KL}(Q_t^\Delta\Vert P_t^\Delta)$ together with the hyperposterior drift, yielding a structurally aligned form of complexity control compared with full finetuning.

\subsection{Perturbation Scope and Type}

Having established the PAC-Bayesian bound for PECL, we now explore where and how to apply perturbations to obtain a flatter minima and improve the performance in PECL.

% \vspace{-1pt}
%the backbone $W^{\mathrm{froz}}_t$ is frozen and 
%the backbone $W^{\mathrm{froz}}_t$ is frozen and 
\minisection{Scope - Why permissible perturbations should lie in $\mathcal W_{\Delta,t}$?}
In PECL, learning is confined to the LoRA update subspace $\mathcal W_{\Delta,t}$. Consequently, only directions in $\mathcal W_{\Delta,t}$ are reachable by future task updates and can affect learning along the training trajectory.
Theorem \ref{thm:pecl-bound} makes this restriction explicit: both the task-wise KL term and the sharpness-dependent empirical term depend only on distributions supported on $\mathcal W_{\Delta,t}$.
Therefore, perturbing the parameter outside $\mathcal W_{\Delta,t}$ imposes impact along unreachable directions, which cannot be compensated by permissible LoRA updates under PECL constraint and can introduce unnecessary loss elevation and unstable optimization behavior.
 

\begin{figure}[tbp]
    \centering
\includegraphics[width=\linewidth]{figures_TRML/lossland_curv1d_combine_loss_title2.pdf}
    % \vspace{-6pt}
    \caption{
    % {How (type)}: 
    We visualize the loss landscape of the models by a 1D probe along the leading eigen-direction of the full curvature.
\textbf{Top}: The probe is restricted to the trainable LoRA parameters, using the leading eigen-direction of the restricted curvature.
\textbf{Bottom}: The probe is applied in the full parameter space, using the leading eigen-direction of the full curvature. See Appendix~\ref{app:curv1d_probe} for implementation details.
%     We visualize the loss of the same trained solutions by a 1D probe along the leading eigen-direction under two evaluation scopes.
% Top: The probe is restricted to the trainable LoRA parameters and follows the leading eigen-direction of the restricted curvature $C_{\Delta,t}$.
% Bottom: The probe is applied in the full parameter space and follows the leading eigen-direction of the full curvature $C_t$.
% Models trained with different flatness objectives show similar profiles in the LoRA-restricted probe but clearly different sensitivity under full-parameter perturbations.
     % Models trained with different flatness objectives yet different sensitivity under full-parameter perturbations. 
    % Different perturbation types shape the geometry of updates within the subspace, leading to different full-parameter loss landscapes. 
    }
    \label{fig:scope_type2}
    \vspace{-10pt}
\end{figure}

\minisection{Type: How Different Perturbations Shape Flatness.}
% Our sharpness measures quantify how much the loss increases under small perturbations around the current parameters. In PECL, perturbations that fall outside the task-permissible subspace are not actionable by subsequent updates, so we focus on subspace-supported covariances $\Sigma_{\Delta,t}$ and analyze how different choices of $\Sigma_{\Delta,t}$ shape the effective curvature seen by training.
Here \emph{type} refers to the different sharpness-aware objectives defined in the section~\ref{section:sahrpness}, each inducing a distinct perturbation rule. 
% , which induce distinct perturbation rules. 
However these objectives are typically formulated in the full parameter space, under PECL they should be interpreted within the task-permissible LoRA update subspace $\mathcal W_{\Delta,t}$, where different types correspond to different ways of allocating perturbation mass and aligning it with the task-conditioned curvature.
% they must be interpreted under the task-permissible constraint under PECL: both updates and perturbations are restricted to the current LoRA update subspace $\mathcal W_{\Delta,t}$, so different types correspond to different ways of distributing perturbation mass and aligning it with the task-conditioned curvature \emph{within} $\mathcal W_{\Delta,t}$.


Let $\Pi_{\Delta,t}$ be the projector onto $\mathcal W_{\Delta,t}$ and let $C_t=C(\hat W_t;S_t)$ be a task-dependent curvature operator of $\hat L_{S_t}$ at $\widehat W_t$.
In PECL, task-relevant parameter updates and subspace sharpness are restricted to $\mathcal W_{\Delta,t}$, hence only the restricted curvature $C_{\Delta,t}=\Pi_{\Delta,t}C_t\Pi_{\Delta,t}$ and perturbations supported on $\mathcal W_{\Delta,t}$ can contribute.
Under a local second-order approximation and $\varepsilon\sim\mathcal N(0,\Sigma_t)$ supported on $\mathcal W_{\Delta,t}$,
{%
% \setlength{\abovedisplayskip}{4pt}
% \setlength{\belowdisplayskip}{3pt}
% \setlength{\abovedisplayshortskip}{4pt}
% \setlength{\belowdisplayshortskip}{3pt}
\[
\mathrm{RS}_t^{\Delta}(\widehat{\Delta W}_t,\Sigma_t)
% =
% \mathbb E_{\varepsilon}\big[\hat L_{S_t}(\hat W_t+\varepsilon)-\hat L_{S_t}(\hat W_t)\big]
\approx
\frac12\operatorname{tr}\!\big(C_{\Delta,t}\Sigma_t\big),
\]
}where $\mathcal W_{\Delta,t}$ fixes the permissible directions, $C_{\Delta,t}$ specifies the restricted curvature on them, $\Sigma_t$ allocates perturbation variance within them, and $\operatorname{tr}(C_{\Delta,t}\Sigma_t)$ measures curvature–noise alignment inside the task-permissible geometry.

{%
% 1) table* 与上下正文的距离（页顶/页底浮动体）
\setlength{\textfloatsep}{2pt}      % 可在 4pt--10pt 之间调
\setlength{\floatsep}{2pt}          % 多个浮动体之间（备用）
\setlength{\dbltextfloatsep}{2pt}   % 某些模板对 table* 用这个（双栏更关键）
\begin{table*}[ht]
\centering
\caption{
Class incremental performance and weight space flatness metrics on ImageNet-R ($T=20$, $R=16$).
% We report standard class-incremental metrics together with sharpness and empirical curvature diagnostics (Hessian and Fisher) {evaluated on the trainable LoRA parameters only}, namely restricted to the adapter-induced subspace rather than the full model parameter space.
We report standard class-incremental metrics together with sharpness measures of the full-model loss under full-parameter perturbations, and empirical curvature diagnostics (Hessian/Fisher) computed on the task-permissible subspace.
% The table reports standard class incremental metrics  together with the sharpness and empirical Hessian and Fisher statistics restricted in the weight space.
}
\label{table:q1_ci_with_weight_flatness_reorg}
\setlength{\tabcolsep}{2.5pt}       % 列间距更小
\renewcommand{\arraystretch}{0.92}  % 行高更紧凑（可选）
\resizebox{\linewidth}{!}{
\scriptsize % 整体字号更小（可改为 \footnotesize 或 \tiny）
\begin{tabular}{@{}p{1.15cm} p{1.8cm} c c c c c c c c@{}}
\toprule
\multirow{2}{*}{\textbf{Method}} & \multirow{2}{*}{\textbf{Opt.}} &
\multicolumn{2}{c}{\textbf{CL Performance ($\uparrow$)}} &
\multicolumn{2}{c}{\textbf{Sharpness ($\downarrow$)}} &
\multicolumn{4}{c}{\textbf{Curvature ($\downarrow$)}} 
% \multicolumn{2}{c}{\textbf{Emp Fisher ($\downarrow$)}}
\\
\cmidrule(lr){3-4}\cmidrule(lr){5-6}\cmidrule(lr){7-10}
 & & $\operatorname{Acc}_{20}^{\mathrm{cls}}$  %&$\operatorname{ACC}^{\mathrm{cls}}$ 
&$\operatorname{Acc}_{20}^{\mathrm{nme}}$ 
 & AS$^{(0)}(\rho)$  & AS$^{(1)}(\rho)$
 & $\lambda_{\max}(H)$ & $\mathrm{tr}(H)$  
 & $\lambda_{\max}(F)$ & $\mathrm{tr}(F)$   \\
\midrule
% \multirow{1}{*}{\textbf{PerTaskLP}}
% & SGD   & $55.22^{\pm 0.31}$ & $59.36^{\pm 0.19}$ & $-8.35^{\pm 0.75}$ & 0.072   & 0.072 &     &   &    &   \\
% \midrule
% \multirow{2}{*}{\textbf{SeqFT}}
% & SGD   &$50.72^{\pm 2.20}$ & $59.89^{\pm 1.32}$ & $-38.94^{\pm 1.87}$ &   &  &   &  &   & \\
% & SAM   &$56.07^{\pm 2.60}$ & $66.47^{\pm 1.53}$ & $-33.04^{\pm 4.71}$&   &  &   &  &   &  \\
%  &GAM  &$59.78^{\pm 0.74}$ & $67.31^{\pm 1.18}$ & $-32.70^{\pm 0.26}$ &   &  &   &  &   &   \\
% \midrule
\multirow{4}{*}{{SeqLoRA}}
& SGD    &$63.14^{\pm 2.17}$ & $72.82^{\pm 1.27}$  & 0.142  & 0.114 & 47.37 & 1502.55 & 35.51 & 1401.15 \\\
& + $\mathrm{RS}^{0}_S(W,\rho)$  &$64.25^{\pm 0.87}$ & $73.01^{\pm 1.11}$ &0.121 &0.108 &47.29 &1444.96 &30.35 &1338.16\\
& + $\mathrm{AS}^{0}_S(W,\rho)$    &$67.04^{\pm 0.71}$ & $74.53^{\pm 0.60}$ & 0.120  & 0.093 & 24.41   & 996.59   & 21.39  & 821.45   \\
 & + $\mathrm{AS}^{(1)}_S(W,\rho)$   &$73.31^{\pm 0.52}$ & $75.74^{\pm 0.15}$  & 0.060  &0.058  & 4.81    & 339.49   & 3.72   & 240.91   \\
\midrule
\multirow{4}{*}{{IncLoRA}}
& SGD    &$65.77^{\pm 2.84}$  &$75.41^{\pm 0.26}$ & 0.076    & 0.062 & 15.46 & 969.34  & 8.89  & 836.91 \\
& + $\mathrm{RS}^{0}_S(W,\rho)$ & $67.42^{\pm 1.23}$ & $75.79^{\pm 0.42}$  &0.075 &0.073 &7.54 &917.35 &8.82 &774.56\\
& + $\mathrm{AS}^{(0)}_S(W,\rho)$    &$68.86^{\pm 0.92}$ & $75.95^{\pm 0.07}$  & 0.053   & 0.052 & 3.72    & 712.88   & 4.35   & 526.49   \\
 &+ $\mathrm{AS}^{(1)}_S(W,\rho)$  &$72.86^{\pm 0.23}$ & $76.72^{\pm 0.54}$  & 0.035   & 0.035 & 2.69    & 281.22   & 2.09   & 192.25   \\
\midrule
\multirow{4}{*}{{OLoRA}}
& SGD    &$66.64^{\pm 0.44}$ & $75.10^{\pm 0.41}$  & 0.078   & 0.058 & 6.94    & 938.93   & 9.00   & 831.23   \\
& + $\mathrm{RS}^{0}_S(W,\rho)$  & $67.53^{\pm 1.04}$ & $75.71^{\pm 0.27}$  &0.076 &0.073 &7.18 &918.33 &8.97 &813.62\\
& + $\mathrm{AS}^{(0)}_S(W,\rho)$     &$68.53^{\pm 0.78}$ & $75.73^{\pm 0.32}$  & 0.066   & 0.048 & 4.40    & 775.40   & 7.41   & 628.12   \\
 & + $\mathrm{AS}^{(1)}_S(W,\rho)$  &$72.44^{\pm 0.36}$ & $76.56^{\pm 0.53}$  & 0.036   & 0.036 & 2.22    & 261.47   & 2.08   & 198.30\\
% \midrule
% \multirow{1}{*}{\textbf{l2p}}
% & Adam   & $71.35^{\pm 0.05}$ & $74.84^{\pm 0.73}$ & $-6.34^{\pm 0.50}$  &-  &- &- &- &- &- \\
% \midrule
% \multirow{1}{*}{\textbf{dualprompt}}
% & Adam   & $71.59^{\pm 0.68}$ & $74.80^{\pm 1.18}$ & $-3.38^{\pm 0.45}$ &-  &- &- &- &- &- \\
\bottomrule
\end{tabular}
}
\end{table*}
} 

This clarifies how perturbation \emph{type} acts in PECL: since adapter update direction is constrained to $\mathcal W_{\Delta,t}$, perturbation types induce distinct curvature-aligned covariances $\Sigma_{\Delta,t}$, leading to different training trajectories and flat minima. 
% different { perturbation types} allocate variance $\Sigma_{\Delta,t }$ in different ways relative to the restricted curvature $C_{\Delta,t}$, leading to different trajectories and  different flatness solutions.
% , by allocating variance differently within $\mathcal W_{\Delta,t}$.
% This clarifies how different optimizers act in PECL. Standard sharpness criteria can be re-expressed in the task-permissible geometry by restricting perturbations to $\mathcal W_{\Delta,t}$.
% Since parameter-efficient continual learning restricts adaptation to a task-admissible set of update directions, it is natural to focus on perturbations supported on the adapter-induced subspace $\mathcal W_{\Delta,t}$ when reasoning about training dynamics and flatness.
% Under this fixed scope, different optimizers mainly differ by the \emph{geometry of the perturbation distribution} within $\mathcal W_{\Delta,t}$, namely by how $\Sigma_t$ concentrates variance inside the same subspace.
% This clarifies how perturbation \emph{type} acts in PECL: since parameter-efficient continual learning restricts adaptation to a task-admissible set of update directions, only perturbations supported on the same adapter-induced subspace $\mathcal W_{\Delta,t}$ can materially alter the optimization trajectory and the resulting flatness; under this fixed scope, different optimizers mainly differ by the \emph{geometry of the perturbation distribution} within $\mathcal W_{\Delta,t}$, namely by how $\Sigma_t$ concentrates variance (or, for adversarial criteria, how the neighborhood is searched) inside the same subspace.
SGD induces an implicit stochastic perturbation on the trainable parameters, which can be viewed as an effective noise covariance on $\mathcal W_{\Delta,t}$~\cite{pmlr-v97-zhu19e}. Random zeroth-order sharpness $\mathrm{RS}^{(0),\Delta}(W,\Sigma)$ corresponds to an explicit Gaussian perturbation $\varepsilon\sim\mathcal N(0,\Sigma)$ supported on $\mathcal W_{\Delta,t}$, and measures the average-case loss elevation. In contrast, adversarial criteria concentrate the perturbation on the same task-permissible subspace: $\mathrm{AS}^{(0),\Delta}$ optimize the worst-case loss increase over the $\rho$-ball in $\mathcal W_{\Delta,t}$ ~\cite{foretSharpnessAwareMinimizationEfficiently2021}, while $\mathrm{AS}^{(1),\Delta}$ further considers the neighborhood gradient norm and is more sensitive to dominant modes of the restricted curvature $C_{\Delta,t}$~\cite{Zhang_GAM}. 


% Under a fixed scope, these choices correspond to increasingly concentrated perturbation geometries within $\mathcal W_{\Delta,t}$~\cite{mollenhoffSAMOptimalRelaxation2022,Zhang_GAM},
% $A$
% and therefore yield progressively stronger control of the sharpness-dependent empirical term in Theorem~\ref{thm:pecl-bound}.
% This mechanism is reflected in Fig.~\ref{fig:scope_type2}. under a full-parameter 1D probe along the leading curvature direction, the trained solutions exhibit clearly separated loss profiles in different order and the order is aligned with perturabtion  definiton. 


Under a fixed scope, these choices correspond to increasingly concentrated perturbation geometries within $\mathcal W_{\Delta,t}$, which induces an ordering of strength under a comparable radius (or covariance scale)~\cite{mollenhoffSAMOptimalRelaxation2022,Zhang_GAM}:
$ \mathrm{AS}^{(1)}_S(W,\rho)\ \ge \mathrm{AS}^{(0)}_S(W,\rho)\ \ge \mathrm{RS}^{(0)}_S(W,\Sigma_t).$
% {%
% \setlength{\abovedisplayskip}{1pt}
% \setlength{\belowdisplayskip}{1pt}
% \setlength{\abovedisplayshortskip}{1pt}
% \setlength{\belowdisplayshortskip}{1pt}
% \begin{equation*}
%     \mathrm{AS}^{(1)}_S(W,\rho)\ \ge \mathrm{AS}^{(0)}_S(W,\rho)\ \ge \mathrm{RS}^{(0)}_S(W,\Sigma_t).
% \end{equation*}
% }
Accordingly, more adversarial objectives impose a stronger penalty on sharp neighborhoods and thus exert progressively stronger control of the sharpness-dependent empirical term in Theorem~\ref{thm:pecl-bound}.
This mechanism is reflected in Fig.~\ref{fig:scope_type2}: under the full-parameter probe along the leading eigen-direction of the full curvature, the trained solutions exhibit clearly separated loss profiles, indicating that perturbation type, although imposed within $\mathcal W_{\Delta,t}$ during training, can translate into different global sensitivity of $W^{\mathrm{froz}}_t+\Delta W_t$.

% and therefore yield progressively stronger control of the sharpness-dependent empirical term in Theorem~\ref{thm:pecl-bound}. This mechanism is reflected in Fig.~\ref{fig:scope_type2}. Under the full-parameter probe along the leading eigen-direction of the full curvature, the trained solutions exhibit clearly separated loss profiles, indicating that different types of perturbation, although imposed inside $\mathcal W_{\Delta,t}$ during training, can translate into different global sensitivity of $W_0+\Delta W_t$.


% This mechanism is  reflected in Fig.~\ref{fig:scope_type2}: under the full-parameter probe, the  solutions exhibit clearly separated profiles.  These observations indicate that different subspace perturbation geometries, while acting inside $\mathcal W_{\Delta,t}$ during training, translate into different global sensitivity of $W_0+\Delta W_t$.


% These definitions imply an ordering in strength within the same subspace and radius constraint~\cite{mollenhoffSAMOptimalRelaxation2022,Zhang_GAM}:
% $\mathrm{AS}^{(1)}_S(W,\rho)\ \ge \mathrm{AS}^{(0)}_S(W,\rho)\ \ge \mathrm{RS}^{(0)}_S(W,\Sigma_t).$
% % {%
% % % \setlength{\abovedisplayskip}{3pt}
% % % \setlength{\belowdisplayskip}{3pt}
% % % \setlength{\abovedisplayshortskip}{1pt}
% % % \setlength{\belowdisplayshortskip}{1pt}
% % \begin{equation*}
% %     \mathrm{AS}^{(1)}_S(W,\rho)\ \ge \mathrm{AS}^{(0)}_S(W,\rho)\ \ge \mathrm{RS}^{(0)}_S(W,\Sigma_t),
% % \end{equation*}
% % }
% Accordingly, optimizing increasingly adversarial perturbation objectives typically yields solutions with smaller sharpness within task subspace. And as the perturbation becomes increasingly concentrated and aligned with the dominant high-curvature direction, it more effectively promotes escape from sharp minima and biases the trajectory toward flatter regions~\cite{pmlr-v97-zhu19e,chen2026sharpnessflatnessdecompositionframework}.



% These definitions imply an ordering in strength within the same subspace and radius constraint~\cite{mollenhoffSAMOptimalRelaxation2022,Zhang_GAM}:
% {%
% \setlength{\abovedisplayskip}{3pt}
% \setlength{\belowdisplayskip}{3pt}
% \setlength{\abovedisplayshortskip}{1pt}
% \setlength{\belowdisplayshortskip}{1pt}
% \begin{equation*}
%     \mathrm{AS}^{(1)}_S(W,\rho)\ \ge \mathrm{AS}^{(0)}_S(W,\rho)\ \ge \mathrm{RS}^{(0)}_S(W,\Sigma_t),
% \end{equation*}
% }
% % $$
% Accordingly, optimizing increasingly adversarial perturbation objectives typically yields solutions with smaller sharpness within task subspace. And as the perturbation becomes increasingly concentrated and aligned with the dominant high-curvature direction, it more effectively promotes escape from sharp minima and biases the trajectory toward flatter regions~\cite{pmlr-v97-zhu19e,chen2026sharpnessflatnessdecompositionframework}.
% This mechanism is directly reflected in Fig.~\ref{fig:scope_type2}: under the full-parameter probe, the same trained solutions exhibit clearly separated profiles. A one-dimensional probe within the LoRA subspace can produce seemingly similar loss curves (top of Fig.~\ref{fig:scope_type2}), however, when the resulting models are evaluated under full-parameter perturbations (bottom of Fig.~\ref{fig:scope_type2}), the profiles separate markedly, revealing that different subspace perturbation geometries lead to different full-model sensitivity. These observations indicate that different subspace perturbation geometries, while acting inside $\mathcal W_{\Delta,t}$ during training, translate into different global sensitivity of $W_0+\Delta W_t$.

% A one-dimensional probe within the LoRA subspace can produce seemingly similar loss curves (top of Fig.~\ref{fig:scope_type2}), however, when the resulting models are evaluated under full-parameter perturbations (bottom of Fig.~\ref{fig:scope_type2}), the profiles separate markedly, revealing that different subspace perturbation geometries lead to different full-model sensitivity.


% the same trained solutions separate markedly when evaluated under \emph{full-parameter} perturbations, indicating that different subspace perturbation geometries translate into different global sensitivity of $W_0+\Delta W_t$




% This clarifies how different optimizers act in PECL. More generally, standard sharpness criteria can be re-expressed in the task-admissible geometry by restricting perturbations to $\mathcal W_{\Delta,t}$. Standard SGD injects an implicit stochastic perturbation with an effective covariance $\Sigma_t^{\mathrm{SGD}}$ determined by minibatch noise\cite{pmlr-v97-zhu19e}. Random zeroth-order sharpness $\mathrm{RS}^{(0)}$ corresponds to explicit isotropic Gaussian noise $\Sigma=\sigma^2 I$, and measures the expected loss elevation under average-case perturbations supported on the admissible subspace. In contrast, adversarial criteria replace the average-case elevation with a worst-case concentration over the same adapter subspace. The adversarial zeroth-order sharpness $\mathrm{AS}^{(0),\Delta}$ chooses the worst-case perturbation that maximizes the local loss increase over $\Sigma=\rho\frac{\nabla L}{|\nabla L|_2}$
% % , yielding a gradient-aligned direction\cite{foretSharpnessAwareMinimizationEfficiently2021}.
% The adversarial first-order sharpness $\mathrm{AS}^{(1),\Delta}$ maximizes the neighborhood gradient norm
% $\Sigma=\rho\frac{\nabla|\nabla L|_2}{|\nabla|\nabla L|_2|_2}$,
% which emphasizes directions directions aligned with dominant eigenmodes of $C_{\Delta,t}$~\cite{zhangGradientNormAware2023}. These definitions imply an ordering in strength within the same subspace and radius constraint~\cite{mollenhoffSAMOptimalRelaxation2022,Zhang_GAM}:
% $\mathrm{AS}^{(1)}_S(W,\rho)\ \ge \mathrm{AS}^{(0)}_S(W,\rho)\ \ge \mathrm{RS}^{(0)}_S(W,\Sigma_t).$
% % {%
% % % \setlength{\abovedisplayskip}{3pt}
% % % \setlength{\belowdisplayskip}{3pt}
% % % \setlength{\abovedisplayshortskip}{1pt}
% % % \setlength{\belowdisplayshortskip}{1pt}
% % \begin{equation*}
% %     \mathrm{AS}^{(1)}_S(W,\rho)\ \ge \mathrm{AS}^{(0)}_S(W,\rho)\ \ge \mathrm{RS}^{(0)}_S(W,\Sigma_t),
% % \end{equation*}
% % }
% Accordingly, optimizing increasingly adversarial perturbation objectives typically yields solutions with smaller sharpness within task subspace. And as the perturbation becomes increasingly concentrated and aligned with the dominant high-curvature direction, it more effectively promotes escape from sharp minima and biases the trajectory toward flatter regions~\cite{pmlr-v97-zhu19e,chen2026sharpnessflatnessdecompositionframework}.



\noindent\textbf{LoRA Organization as a Structural Prior.}\quad 
To make the hyper process concrete in PECL, we represent each task-level prior by a minimal \emph{hyper-state}
$
\phi_t = \bigl(\mathcal W_{\Delta,t},\mu_{P,t},\Sigma_{P,t}\bigr),
$
where $\mathcal W_{\Delta,t}$ specifies the permissible adapter update scope and $(\mu_{P,t},\Sigma_{P,t})$ parameterize a task prior on these coordinates.
Sampling $P_t\sim\mathcal P_{t-1}$ can then be interpreted as drawing $\phi_t\sim\mathcal P_{t-1}$ and instantiating the corresponding task prior. Under this view, SeqLoRA, IncLoRA, and constraint-augmented designs correspond to different evolution rules for $\phi_t$, and thus shape both the subspace sharpness channel and the KL complexity in Theorem~\ref{thm:pathwise_pac}.




SeqLoRA keeps an approximately fixed scope $\mathcal W_{\Delta,t}\equiv\mathcal W_\Delta$ and concentrates successive updates in the same geometry, which tends to induce stronger direction competition across tasks.
% under the deterministic choice $P_t^\Delta=Q_{t-1}^\Delta$, the complexity reduces to $\sum_t\mathrm{KL}(Q_t^\Delta\Vert Q_{t-1}^\Delta)$.
% This shared scope limit admissible directions for new task and amplify cross-task interference for previous task.
% A fixed scope concentrates successive updates in the same geometry, which tends to induce stronger direction competition across tasks.
IncLoRA relaxes the fixed-scope assumption by expanding $\mathcal W_{\Delta,t}$ over time, which redistribute adaptation across additional degrees of freedom and reduced subspace overwrite, while the increased flexibility may also admit more directions with projected curvature.
% . At the same time, a larger freedom may admit more directions with projected curvature.
% This provides extra admissible directions and reduce the influence of previous task direction. 
Constraint-augmented designs further control how new degrees of freedom are allocated and initialized by imposing structure on $\mathcal W_{\Delta,t}$, for example via orthogonality regularization~\cite{wang2023orthogonal}, gradient projection~\citep{liangInfLoRAInterferenceFreeLowRank2024a,yang2024parametercollisionhinderingcontinual}.
It reduce repeated overwrite and stabilize the evolution of the permissible geometry.


\section{Experiments and Analysis}

% \subsection{Experimental Setup}



\noindent\textbf{Datasets.}\quad
We evaluate PECL methods on vision benchmarks derived from ImageNet. ImageNet-R~\cite{ImagenetR_Hendrycks_2021_ICCV} serves as our main class-incremental benchmark following prior work~\cite{liangInfLoRAInterferenceFreeLowRank2024a}.
%and is split into $T=20$ tasks over $200$ classes
% To assess robustness under corruptions and perturbations, we use the Tiny ImageNet variants of ImageNet-C and ImageNet-P (200 classes)~\cite{hendrycks2018benchmarking}. 
% For brevity, we refer to these datasets as “ImageNet-C” and “ImageNet-P”.
To assess robustness under corruptions and perturbations, we use the Tiny ImageNet variants of ImageNet-C and ImageNet-P (200 classes)~\cite{hendrycks2018benchmarking}.
%, where the former applies 15 types of corruption types and the latter provides short perturbation sequences per image

% We further adopt Tiny ImageNet-C~\cite{hendrycks2018benchmarking} and Tiny ImageNet-P~\cite{hendrycks2018benchmarking} to evaluate robustness on common corruptions and perturbations. Tiny ImageNet-C is a corrupted variant of Tiny ImageNet with 200 classes containing 15 types of corruption. Tiny ImageNet-P provides short perturbation sequences for each image
%throughout the rest of the paper
% To assess robustness under corruptions and perturbations, we further use ImageNet-C and ImageNet-P~\cite{hendrycks2018benchmarking}. 
% ImageNet-C applies common corruptions to ImageNet validation images, while ImageNet-P provides short perturbation sequences per image.

{%
\setlength{\intextsep}{4pt}
\setlength{\textfloatsep}{4pt}
\setlength{\abovecaptionskip}{6pt}
\setlength{\belowcaptionskip}{4pt}
\captionsetup{skip=10pt}
\begin{figure*}[t]
\centering
\includegraphics[width=1\linewidth]{figures_TRML/weight_overlap_5fig.pdf}
\vspace{-15pt}
\caption{Evolution of the curvature projection ratio $c_t$ (left), amplification factor $\alpha_t$ (middle) and loss with perturbation $\hat{L}(W+\epsilon)$ and the average accuracy $\operatorname{Acc}_t$ (right) in seqLoRA.
% across the task sequence, together with the corresponding average accuracy $\operatorname{Acc}_t$ (dashed) in seqLoRA.
% under SGD, $\mathrm{AS}^{(0)}$, and $\mathrm{AS}^{(1)}$
% Evolution of the amplification factor $\alpha_t$ (top) and curvature projection ratio $c_t$ (bottom) under SGD and apply perturbation $\mathrm{AS}^{(1)}$ and $\mathrm{AS}^{(1)}$ in seqLoRA.
% Compared with SGD, $\mathrm{AS}^{(1)}$ produces consistently larger and increasing $\alpha_t$ and $c_t$ over tasks, indicating progressively stronger adapter-induced amplification and a growing concentration of task-conditioned curvature in the adapter-induced geometry.
$\alpha_t$ and $c_t$ are computed from the last layer of ViT-B/16.
% Evolution of the amplification factor $A_t$ and curvature projection ratio $C_t$ across tasks in SeqLoRA. The top panel shows $A_t$, measuring how strongly the adapter update $\Delta W_t$ amplifies directions relative to the backbone restricted to the same rank-$r$ subspace. The bottom panel shows $C_t$, comparing the concentration of dominant task-conditioned curvature modes in the adapter-induced subspace versus the backbone subspace. The weight matrices are taken from the  11th layer of ViT-B/16.
}
\label{fig:evolute}
\vspace{-15pt}
\end{figure*}
}


\noindent\textbf{Models.}\quad 
All experiments use a ViT-B/16~\cite{dosovitskiy2020image} pre-trained on ImageNet-21K.
We attach LoRA adapters to each transformer block and consider three variants: SeqLoRA, IncLoRA, and OLoRA~\cite{wang2023orthogonal}.

% Unless stated otherwise, PECL updates and perturbations are applied only to
% to the trainable LoRA coordinates of the current task, while the backbone and past adapters remain frozen. 





% To study perturbation {scope}, we consider a Flat-LoRA-style~\cite{li2024flatlora} as a \emph{full-space} baseline, which applies $\mathrm{RS}^{(0)}_S(W,\Sigma)$ perturbations over all model parameters.
To study perturbation \emph{scope}, we adopt {full-parameter} setting~\cite{li2024flatlora} by injecting isotropic Gaussian perturbation $\mathrm{RS}_S^{(0)}({W}, \sigma^2 
I)$ into all model parameters.
% , instead of restricting perturbations to the adapter subspace.
% To study perturbation {type}, we compare SAM~\cite{foretSharpnessAwareMinimizationEfficiently2021} ($\mathrm{AS}^{(0)}_S(W,\rho)$), RWP~\cite{pmlr-v151-bisla22a} ($\mathrm{RS}^{(0)}_S(W,\Sigma)$), and GAM~\cite{zhangGradientNormAware2023} ($\mathrm{AS}^{(1)}_S(W,\rho)$), all acting on the trainable adapters in PECL.
To study perturbation \emph{type}, we compare SAM~\cite{foretSharpnessAwareMinimizationEfficiently2021}, RWP~\cite{pmlr-v151-bisla22a}, and GAM~\cite{Zhang_GAM}, which optimize different sharpness objectives: SAM minimizes $\mathrm{AS}^{(0)}_S(W,\rho)$, RWP targets $\mathrm{RS}^{(0)}_S(W,\Sigma)$, and GAM targets $\mathrm{AS}^{(1)}_S(W,\rho)$. For a controlled comparison in PECL, all methods apply their perturbations only to the trainable LoRA coordinates.
% $\theta_\Delta$; the induced weight perturbation is $\Delta W(\theta_\Delta+\xi)-\Delta W(\theta_\Delta)$.
%all methods apply their perturbations only to the trainable adapters.
Unless otherwise stated, we fix the perturbation radius to $\rho=0.05$, following the standard empirical setting used in SAM, and use a dedicated radius sweep only when diagnosing the scope effect.




% based on the task accuracy matrix $a_{tj}$, where $a_{tj}$ denotes the accuracy on task $j$ after learning task $t$. 
\noindent\textbf{Metrics.}\quad
The average class-incremental accuracy at step $t$ is
$
\operatorname{Acc}_t = \frac{1}{t} \sum_{j=1}^{t} a_{tj}
$,  where $a_{tj}$ denotes the accuracy on task $j$ after learning task $t$. 
We report accuracy $\operatorname{Acc}_t^{\mathrm{cls}}$ using the linear classifier and
% , and also evaluate representation quality using Nearest Class Mean (NCM) classifier, denoted by $\operatorname{Acc}_t^{\mathrm{nme}}$.
compute the corresponding Nearest-Class-Mean Classifier (NCM) results
% Nearest Mean of Exemplars (NME) metrics 
$(\operatorname{Acc}_t^{\mathrm{ncm}})$ using class prototypes in the feature space.
%, which isolates the feature representations from the classifier head
% We report class-incremental performance using the task-wise classification accuracy $\operatorname{Acc}_t^{\mathrm{cls}}$.
% In addition, we evaluate representation quality via Nearest Mean of Exemplars (NME), denoted by $\operatorname{Acc}_t^{\mathrm{nme}}$.
% This metric reduces the influence of the classifier head and provides a complementary view of how stable and transferable the learned representations are across tasks.
To quantify flatness, we estimate the largest eigenvalue $\lambda_{\max}$ and the trace $\mathrm{tr}(\cdot)$ of both the Hessian and the Fisher information matrix using Lanczos~\cite{ghorbaniInvestigationNeuralNet2019}.

\subsection{Result Analysis}





\noindent\textbf{Scope of perturbations.}\quad
To study perturbation scope in PECL, we inject isotropic Gaussian perturbation $\mathrm{RS}_S^{(0)}({W}, \sigma^2 
I)$ under two scopes: perturbations restricted to the task-permissible adapter coordinates versus perturbations defined over the full parameter space.
As shown in Table~\ref{tab:lora_or_full_cls—show}, restricting perturbations to the trainable adapters consistently improves average accuracy over the SGD baseline, whereas applying perturbations that include the frozen backbone substantially degrades performance, with the strongest drop observed for SeqLoRA.



% Let $U_{\Delta,t}$ and $ V_{\Delta,t}$ are the left- and right-singular matrices of the SVD decomposition of $\Delta W$, and $U_{\Delta,t}U_{\Delta,t}^\top W  V_{\Delta,t}^\top  V_{\Delta,t}$ gives the “projection” of $W $onto the subspace spanned by $\Delta W$. Following~\cite{hu2022lora}, we adopt the feature amplification factor
% $
% A_t =
% \frac{\left\lVert \Delta W_{t}\right\rVert_F}
% {\left\lVert U_{\Delta,t}^\top W  V_{\Delta,t}^\top\right\rVert_F}
% $
% to measure how much of task-specific directions are amplified by $\Delta W$. 
% Here $\|U_{\Delta,t}^\top W V_{\Delta,t}^\top\|_F$ measures the energy of $W$ restricted to the same bilinear subspace induced by $\Delta W_t$ and equals the Frobenius norm of the projected of $W$ onto the subspace spanned by $\Delta W$ $U_{\Delta,t}(U_{\Delta,t}^\top W V_{\Delta,t}^\top) V_{\Delta,t}$.
% where $|U_{\Delta,t}^\top W V_{\Delta,t}|_F$ measures the energy of $W$ restricted to the same bilinear subspace induced by $\Delta W_t$. Equivalently, since $U{\Delta,t}$ and $V_{\Delta,t}$ are orthonormal, this quantity equals the Frobenius norm of the projected weight $\Pi_{\Delta,t}(W)=U_{\Delta,t}(U_{\Delta,t}^\top W V_{\Delta,t})V_{\Delta,t}^\top$
% Here $(U_{\Delta,t},V_{\Delta,t})$ are the rank-$r$ left and right singular vector matrices of $\Delta W $
% left and right singular vectors of $\Delta W_t$ from its SVD~\cite{hu2022lora}.
% and $\Delta U_t \Delta U_t^\top W \Delta V_t\Delta V_t^\top$ corresponds to the projection of $W$ onto the subspace spanned by $\Delta W_t$.

{%
\setlength{\textfloatsep}{0pt}      % 页顶/页底 table 与正文的距离（最关键）
\setlength{\intextsep}{1pt}         % 若 table 被放在正文中时的上下距离（备用）
\setlength{\floatsep}{1pt}          % 两个浮动体之间的距离（备用）
\setlength{\abovecaptionskip}{2pt}  % caption 与表格之间
\setlength{\belowcaptionskip}{8pt}  % caption 下方到正文
\captionsetup{skip=4pt}           % 若你用 caption 包，可用它更细调
\begin{table}[t]
\vspace{-5pt}
\centering
\caption{Scope effect in PECL: reporting $\Delta\operatorname{Acc}^{\mathrm{cls}}$ (vs.\ SGD), adapter-only perturbations consistently improve accuracy, whereas full-parameter perturbations degrade performance.}
\label{tab:lora_or_full_cls—show}
\footnotesize
\setlength{\tabcolsep}{3pt}
\renewcommand{\arraystretch}{0.88}
\begin{tabular}{@{}c c c c c@{}}
\toprule
\textbf{Scope} & \textbf{Strength} & \textbf{SeqLoRA} & \textbf{IncLoRA} & \textbf{OLoRA} \\
\midrule
% Base & --   & 69.07 & 71.29 & 71.41 \\
% \midrule
\multirow{3}{*}{Full}
 & 0.05 & -3.50 & -2.00 & -2.62 \\
 & 0.10 & -4.82 & -4.74 & -5.53 \\
 & 0.15 & -6.72 & -6.84 & -6.49 \\
\midrule
\multirow{3}{*}{LoRA}
 & 0.05 & 0.27 & 0.79 & 0.12 \\
 & 0.10 & 0.70 & 0.86 & 0.28 \\
 & 0.15 & 0.41 & 0.77 & 0.27 \\
\bottomrule
\end{tabular}
\end{table}
\vspace{-5pt}
}% end local group


To understand why the perturbation scope is decisive, we examine how the learned update $\Delta W_t$ interacts with the frozen weight $W^\mathrm{froz}_t$ and how dominant task-conditioned curvature modes localize relative to the geometry induced by $\Delta W_t$.
% To understand why the perturbation scope is decisive, we further analyze how the learned adapter update $\Delta W_t$ interacts with the frozen backbone weight $W$ and the task-conditioned curvature.
% To elucidate why the perturbation scope is decisive, we examine where the dominant task-conditioned curvature modes localize relative to the bilinear subspaces induced by the learned update $\Delta W_t$ and by the frozen backbone weight $W$.
% To elucidate why the perturbation scope is decisive, we analyze how the learned adapter update $\Delta W_t$ interacts with the frozen backbone weight $W$ and the task-conditioned curvature.
% Let $\Delta W_t = U_{\Delta,t} S_{\Delta,t} V_{\Delta,t}^\top$ be the rank-$r$ truncated SVD, where
% Let $U_{\Delta,t}$ and $V_{\Delta,t}^\top$ are the left and right singular vectors of the SVD decomposition of $\Delta W$.
% and the "projection” of matrix $W$ onto the subspace spanned by $\Delta W$ is
% $U_{\Delta,t}(U_{\Delta,t}^\top W V_{\Delta,t})V_{\Delta,t}^\top.$
% Frobenius-orthogonal projection of $W$ onto the bilinear subspace induced by $(U_{\Delta,t},V_{\Delta,t})$ 
Following~\citet{hu_lora_2021}, we adopt the feature amplification factor
$
\alpha_t=\frac{\|\Delta W_t\|_F}{\|U_{\Delta,t}^\top W^{\mathrm{froz}}_t V_{\Delta,t}\|_F},
$ to measure how much of task-specific directions are amplified by $\Delta W$, where $U_{\Delta,t}$ and $V_{\Delta,t}^\top$ are the left and right singular vectors of the SVD decomposition of $\Delta W$.
Similarly, to quantify whether dominant task-conditioned curvature modes concentrate more on the subspace induced by the learned update $\Delta W_t$ than on a backbone reference, we define the curvature projection ratio
$
c_t=\frac{R_{\Delta,t}}{R_{W^{\mathrm{froz}_t},t}}.
$
Here $R_{\Delta,t}$ and $R_{W,t}$ are localization scores of the top-$k$ task-conditioned curvature modes on the update-induced and backbone-reference subspaces, respectively.
The precise definitions of $R_{\Delta,t}, R_{W,t}$ are provided in Appendix~\ref{app:weight_projection}.
Larger $c_t$ indicates stronger curvature localization on the update-induced subspace relative to the backbone reference.
% the projection operators, as well as 
% Similarly, to quantify whether dominant task-conditioned curvature modes concentrate more on the subspace induced by the learned update $\Delta W_t$ than on a backbone reference subspace, 
% % Similarly, to compare how strongly dominant task-conditioned curvature modes concentrate more onto the adapter-induced subspace than onto the backbone reference subspace,
% % the adapter-induced subspace,
% % direction project onto the adapter-induced subspace versus a backbone reference subspace,
% % to assess whether dominant curvature modes have larger projection onto the adapter-induced subspace than onto the backbone reference subspace, 
% we define the curvature projection ratio 
% % $
% % c_t=\frac{\sum_{i=1}^k \lambda_{t, i}\left\|U_{\Delta, t}^{\top} {v_{t, i}} V_{\Delta, t} \right\|_F^2}{\sum_{i=1}^k \lambda_{t, i}\left\|U_{W, t}^{\top} v_{t, i} V_{W, t} \right\|_F^2},
% % $
% $
% c_t=\frac{R_{\Delta,t}}{R_{W,t}}.
% $
% where $R_{\Delta,t}$ and $R_{W,t}$ summarize the eigenvalue-weighted localization of the top-$k$ task-conditioned curvature modes onto  the  subspace induced by the learned update $\Delta W_t$ and a backbone reference subspace.
% after mapping each curvature eigen-direction to its effective weight-space perturbation. 
% The precise definitions of $R_{\Delta,t}, R_{W,t}$ and the projection operators, as well as the implementation details, are provided in Appendix~\ref{app:weight_projection}. 
% Larger $c_t$ indicates that dominant curvature modes exhibit stronger projected energy within the subspace induced by $\Delta W_t$ than within the backbone reference subspace.
% where ${(\lambda_{t,i},v_{t,i})}_{i=1}^k$ are the top-$k$ curvature eigenpairs at task $t$ in the full parameter space (Appendix~\ref{app:weight_projection}). 
% Larger $c_t$ indicates that dominant curvature modes have higher projected energy within subspace induced by $\Delta W_t$ than within the backbone reference subspace. 
% Moreover, let $\{(\lambda_{t,i},v_{t,i})\}_{i=1}^k$ be the top-$k$ eigenpairs of a positive semidefinite curvature operator at task $t$ computed in the full parameter space. We quantify curvature localization by eigenvalue-weighted projected-energy ratios onto $\mathcal S_{\Delta,t}$ and onto a backbone reference subspace, and report their preference ratio $c_t$ (Appendix~\ref{weight projection}).

% Larger $C_t$ indicate that the dominant curvature modes align more with the adapter-induced subspace spanned by  $(U_{\Delta,t},V_{\Delta,t})$.

Fig.~\ref{fig:evolute} shows consistent trends across the task sequence and aligns closely with the evolution of accuracy.
As the task index increases, the curvature projection ratio $c_t$ increases over tasks and remains above the reference level $c_t\approx 1$, suggesting that dominant task-conditioned curvature modes are concentrated in the adapter-induced subspace rather than in the backbone subspace.
The amplification factor $\alpha_t$ also rises, indicating that the learned update $\Delta W_t$ progressively amplifies task-specific directions.
Crucially, this structural strengthening is accompanied by the evolution of the perturbation-aware loss $\hat L+\mathrm{RS}^{(0),\Delta}$ and aligns with the degradation of $\operatorname{Acc}_t$.



{%
% {%
%
% \setlength{\intextsep}{4pt}
\setlength{\textfloatsep}{4pt}
% \setlength{\abovecaptionskip}{6pt}
% \setlength{\belowcaptionskip}{pt}
% \captionsetup{skip=10pt}
\begin{figure}[ht]
\centering
\includegraphics[width=0.95\linewidth]{figures_TRML/ACC_dataset_pair_Robustness_to_Corruptions_on_ImageNet-C_Robustness_to_Perturbation_on_ImageNet-P2.pdf}
% \includegraphics[width=0.9\linewidth]{figures_TRML/hebing_robutness2.pdf}

\vspace{-6pt}
\caption{\textbf{Robustness of perturbation types under CL} Classification average accuracy $\operatorname{Acc}_t^{\mathrm{cls}}$ over tasks evaluated under distribution shift on ImageNet-C  and ImageNet-P.
Across tasks and LoRA organizations, both $\mathrm{AS}^{(0)}$ and $\mathrm{AS}^{(1)}$ consistently outperform SGD, and the gap widens in later tasks. Here ImageNet-C/P denote the Tiny ImageNet variants; see Experimental Setup.}
\label{fig:robustness_imagenet_cp}
\vspace{-10pt}
\end{figure}
}%

Taken together, these diagnostics offer a mechanism-level account of Table~\ref{tab:lora_or_full_cls—show}. As training progresses, both the effective adaptation directions and the task-conditioned curvature that governs sensitivity concentrate within the task-permissible subspace. Perturbations restricted to the task-permissible subspace therefore act on the controllable geometry and can be counteracted by permissible LoRA updates, whereas full-scope perturbations introduce variations in the frozen model outside the permissible update subspace, creating an irreducible scope mismatch that accelerates performance degradation under the PECL constraint.
%scope
% Moreover, the perturbation type shapes the covariance geometry in LoRA coordinates, which steers the evolution of $\Delta W_t$ and biases the full-model solution $W_0+\Delta W_t$ toward flatter task-conditioned regions associated with better continual accuracy.

% Taken together, these diagnostics provide a mechanism-level explanation for Table~\ref{tab:lora_or_full_cls—show}:
% both the effective adaptation directions and the task-conditioned curvature that governs sensitivity become increasingly concentrated within the LoRA-admissible coordinates; therefore, perturbations restricted to the adapter scope act on the same controllable geometry whereas full-scope perturbations inject variations into frozen backbone that are outside the admissible update subspace, creating a scope mismatch that cannot be compensated by admissible LoRA updates under PECL constraint and accelerates performance degradation.

%  Under a fixed PECL scope, the perturbation design acts through the geometry within LoRA coordinates, which determines how $\Delta W_t$ evolves and, consequently, how the full-model solution $W_0+\Delta W_t$ bias toward flatter task-conditioned regions.


% Taken together, these diagnostics provide a mechanism-level explanation for Table~\ref{tab:lora_or_full_cls—show}: both adaptation and the curvature that governs sensitivity are increasingly localized within the LoRA-admissible coordinates in CL, so restricting perturbations to the adapter scope matches the controllable geometry, whereas full-scope perturbations inject variations into frozen backbone that are outside the admissible update subspace, creating a scope mismatch that cannot be compensated by admissible LoRA updates under PECL constraint and accelerates performance degradation.

% This strengthening of $\alpha_t$ and $c_t$ coincides with the growing performance gap: methods that maintain larger and increasing $\alpha_t$ and $c_t$ exhibit less loss and more stable $\operatorname{Acc}_t$,


% Fig.~\ref{fig:evolute} shows consistent trends across the task sequence: $\alpha_t$ increases with the task index, indicating that $\Delta W_t$ progressively strengthens task-specific directions. $C_t$ also increases over tasks and remains above the neutral reference level $c_t\approx 1$, suggesting that the dominant task-conditioned curvature becomes increasingly concentrated in the adapter-induced subspace rather than in the backbone subspace.






\noindent\textbf{Type of perturbations.}\quad
To isolate the effect of perturbation \emph{type} under PECL constraints, we apply three flatness optimizers to the \emph{trainable LoRA parameters only}:
$\mathrm{RS}^{(0)}_S(W,\Sigma)$, $\mathrm{AS}^{(0)}S(W,\rho)$, and $\mathrm{AS}^{(1)}S(W,\rho)$, with SGD as the baseline.
Across tasks on ImageNet-C and ImageNet-P (Fig.~\ref{fig:robustness_imagenet_cp}), both adversarial objectives consistently outperform SGD across the task sequence, and the gap widens in later tasks, indicating improved stability under continual adaptation. This trend is consistent with the final class-incremental results on ImageNet-R (Table~\ref{table:q1_ci_with_weight_flatness_reorg}), where moving from SGD to $\mathrm{RS}^{(0)}$, $\mathrm{AS}^{(0)}$, and then $\mathrm{AS}^{(1)}$ yields progressively higher $\operatorname{Acc}^{\mathrm{cls}}_{20}$ and $\operatorname{Acc}^{\mathrm{nme}}_{20}$.


% $\mathrm{AS}^{(0)}$ and $\mathrm{AS}^{(1)}$ outperform SGD throughout the sequence, and the advantage becomes more pronounced in later tasks, indicating improved stability under continual adaptation.
% This trend is consistent with the final class-incremental results in Table~\ref{table:q1_ci_with_weight_flatness_reorg}: across LoRA organizations, moving from SGD to $\mathrm{RS}^{(0)}$, $\mathrm{AS}^{(0)}$, and then $\mathrm{AS}^{(1)}$ yields progressively higher $\operatorname{ACC}{20}^{\mathrm{cls}}$ and $\operatorname{ACC}{20}^{\mathrm{nme}}$.
% Overall, stronger perturbation types lead to systematically better continual performance even when the update scope is fixed to LoRA, suggesting that the \emph{geometry} of perturbations inside the adapter coordinates matters beyond mere subspace restriction.


% To isolate the effect of perturbation \emph{type} under PECL constraints, we apply three representative flatness optimizers to the {trainable LoRA parameters only}: 
% $\mathrm{RS}^{(0)}_S(W,\Sigma)$, $\mathrm{AS}^{(0)}_S(W,\rho)$, and $\mathrm{AS}^{(1)}_S(W,\rho)$, with SGD as the baseline.
% Across tasks on ImageNet-C and ImageNet-P (Fig.~\ref{fig:robustness_imagenet_cp}), we observe a consistent performance ordering, where both $\mathrm{AS}^{(0)}$ and $\mathrm{AS}^{(1)}$ outperform SGD throughout the sequence, and the margin becomes more pronounced in later tasks, indicating improved stability under continual adaptation.
% % Fig.~\ref{fig:robustness_imagenet_cp} shows a consistent ordering across tasks on ImageNet-C and ImageNet-P: both $\mathrm{AS}^{(0)}$ and $\mathrm{AS}^{(1)}$ dominate SGD throughout the sequence, and the gap becomes more visible in later tasks, indicating improved stability under continual adaptation. 
% The same ordering is reflected by the final class-incremental performance and weight-space geometry in Table~\ref{table:q1_ci_with_weight_flatness_reorg}. Across LoRA configurations, moving from SGD to $\mathrm{RS}^{(0)}$, $\mathrm{AS}^{(0)}$, and then $\mathrm{AS}^{(1)}$ yields progressively higher $\operatorname{Acc}_{20}^{\mathrm{cls}}$ and $\operatorname{Acc}_{20}^{\mathrm{nme}}$, while simultaneously reducing sharpness and curvature diagnostics, including $\lambda{\max}(H)$ and $\operatorname{tr}(H)$, and similarly for Fisher. Overall, perturbation types that more strongly suppress adapter-subspace sharpness achieve better continual performance, consistent with the sharpness-dependent empirical term in Theorem~\ref{thm:pecl-bound}.


% We next examine how perturbation \emph{type} changes curvature and the loss surface under the same PECL scope.
% Quantitative diagnostics in Table~\ref{table:q1_ci_with_weight_flatness_reorg} show that stronger optimizers reduce sharpness and curvature statistics measured on the trainable LoRA parameters, including $\lambda_{\max}(H)$ and $\operatorname{tr}(H)$, and similarly for Fisher.
% Interestingly, a one-dimensional probe {within the LoRA subspace} (top of Fig.~\ref{fig:perturb_type_landscape}) can yield loss curves that appear comparably flat, which by itself may understate the differences among perturbation types.
% However, when the same trained models are probed under {full-parameter} perturbations (bottom of Fig.~\ref{fig:perturb_type_landscape}), the profiles separate markedly, revealing distinct sensitivity of the effective model $W_0+\Delta W$ in the full space.
% Taken together, these results indicate that, even though learning is restricted to LoRA coordinates, different perturbation geometries reshape the \emph{full-model} landscape of $W_0+\Delta W$ for task, consistent with the sharpness-dependent empirical term in Theorem~\ref{thm:pecl-bound}.

% Under the same LoRA-restricted update scope, perturbation type still induces distinct curvature and loss-surface properties. Table~\ref{table:q1_ci_with_weight_flatness_reorg} shows that stronger objectives reduce LoRA-restricted sharpness and curvature diagnostics (Hessian and Fisher). Together, these results support the bound perspective that, at a fixed scope, type acts primarily through the perturbation geometry inside $\mathcal W_{\Delta,t}$, thereby reshaping the effective landscape.
% Beyond accuracy, perturbation type induces a coherent set of geometric effects.
% % under the same LoRA-restricted update scope.
% As predicted by the bound, changing type modifies the perturbation geometry within $\mathcal W_{\Delta,t}$ and primarily acts by suppressing task-conditioned sharpness on the permissible parameters.
% Empirically, Table~\ref{table:q1_ci_with_weight_flatness_reorg} shows a consistent ordering across objectives: stronger types yield smaller LoRA-restricted sharpness and lower curvature diagnostics (Hessian and Fisher).
% This reduction in subspace sharpness is accompanied by a corresponding change in the \emph{full-parameter} landscape.
% As visualized in Fig.~\ref{fig:scope_type2}, the solutions obtained with increasingly adversarial objectives exhibit flatter loss profiles under the full-parameter probe along the leading eigen-direction of the full curvature.
% , demonstrating reduced global sensitivity of $W_0+\Delta W_t$.
% Taken together, the table and the loss-landscape visualization form a consistent chain of evidence that perturbation type, through lowering sharpness within $\mathcal W_{\Delta,t}$, reshapes both restricted curvature and the resulting full-model geometry, which explains the observed performance gains in continual adaptation.
Beyond accuracy, perturbation type induces distinct curvature and landscape properties under the task-permissible subspace.
% restricted update scope.
Table~\ref{table:q1_ci_with_weight_flatness_reorg} shows that moving from $\mathrm{RS}^{(0)}$ to $\mathrm{AS}^{(0)}$ and $\mathrm{AS}^{(1)}$ progressively reduces the subspace sharpness and curvature diagnostics (Hessian and Fisher), indicating that more concentrated perturbation geometries better suppress subspace sharpness within $\mathcal W_{\Delta,t}$.
Consistently, Fig.~\ref{fig:scope_type2} visualizes the resulting {full-parameter} landscape: solutions trained with stronger objectives exhibit flatter profiles along the leading eigen-direction of the full curvature.

% Together, these results support the bound perspective that, type acts primarily through the perturbation geometry inside $\mathcal W_{\Delta,t}$, which lowers task-conditioned sharpness and translates into flatter full-model solutions, thereby improving continual robustness and accuracy.


% Fig.~\ref{fig:perturb_type_landscape} and Fig.~\ref{fig:evolute} further clarify \emph{how} different perturbation geometries translate into different full-model solutions even when updates are restricted to LoRA.
% Within a LoRA-restricted slice (top of Fig.~\ref{fig:perturb_type_landscape}), the loss profiles can appear similarly flat across perturbation types. However, evaluating the same model under full-parameter perturbations (bottom of Fig.~\ref{fig:perturb_type_landscape}) separates the curves, revealing distinct sensitivity of the effective model $W_0+\Delta W$. 
% Along a one-dimensional direction \emph{within the LoRA subspace} (top of Fig.~\ref{fig:perturb_type_landscape}), the loss curves can look uniformly flat across perturbation types.
% However, the quantitative flatness diagnostics in Table~\ref{table:q1_ci_with_weight_flatness_reorg} reveal clear differences even under the same LoRA restriction: stronger flatness optimizers achieve smaller sharpness and reduced curvature statistics when curvature is evaluated on the trainable LoRA parameters.
% Moreover, when the same trained models are probed under \emph{full-parameter} perturbations (bottom of Fig.~\ref{fig:perturb_type_landscape}), the profiles separate markedly, exposing distinct sensitivity of the effective model $W_0+\Delta W$ in the full parameter space.
% This separation is consistent with the trends of sharpness and projected-curvature statistics reported in Table~\ref{table:q1_ci_with_weight_flatness_reorg}, where stronger flatness optimizers achieve smaller full-model curvature measures and exhibit flatter full-space loss profiles.

% \textbf{why $\mathrm{AS}^{(1)}$ is so better in LoRA based CL?}

The most pronounced effect is the amplified advantage of $\mathrm{AS}^{(1)}$ in LoRA-based PECL. Unlike $\mathrm{AS}^{(0)}$, $\mathrm{AS}^{(1)}$ explicitly penalizes the neighborhood gradient norm and is therefore more sensitive to the dominant modes of the task-conditioned restricted curvature $C_{\Delta,t}$. In a low-dimensional, structured task-permissible subspace, this curvature-aware control is easier to realize and has larger impact: in SeqLoRA, only $\mathrm{AS}^{(1)}$ induces a persistent increase in the localization ratios $(c_t,\alpha_t)$ (Fig.~\ref{fig:evolute}), which coincides with the lowest loss and most stable $\operatorname{Acc}_t$. This coupling suggests that $\mathrm{AS}^{(1)}$ preferentially shapes the few permissible directions, simultaneously concentrating task-relevant curvature into the adapter-induced geometry and steering the effective model toward flatter task-conditioned regions. Consequently, under a fixed PECL scope, perturbation type determines the geometry of perturbations within LoRA coordinates, which in turn determines $\Delta W_t$ and the full-model loss landscape of $W_0+\Delta W_t$, explaining the consistent performance ordering.

% The most striking observation is the amplified advantage of $\mathrm{AS}^{(1)}$ in LoRA-based PECL.
% Compared with SGD and the other perturbation types, $\mathrm{AS}^{(1)}$ achieves the best class-incremental accuracy, the smallest LoRA-restricted sharpness and curvature diagnostics (Table~\ref{table:q1_ci_with_weight_flatness_reorg}), and the flattest full-space loss profiles (Fig.~\ref{fig:perturb_type_landscape}), with a consistently better trajectory across tasks (Fig.~\ref{fig:robustness_imagenet_cp}).
% Notably, such a pronounced separation is less evident under standard full fine-tuning in the original study~\cite{Zhang_GAM}, suggesting that first-order sharpness control becomes particularly effective when adaptation is \emph{subspace-constrained}.
% To understand this amplification, we compare $\mathrm{AS}^{(1)}$ against $\mathrm{AS}^{(0)}$ and SGD in the SeqLoRA setting.
% Among all methods, only $\mathrm{AS}^{(1)}$ induces a persistent increase in the curvature localization ratios $c_t$ and $\alpha_t$, which coincides with the highest and most stable $\operatorname{ACC}t$.
% This coupling supports the following mechanism: by explicitly penalizing the neighborhood gradient norm, $\mathrm{AS}^{(1)}$ is more sensitive to dominant modes of the task-conditioned restricted curvature $C{\Delta,t}$, and thus preferentially shapes the few admissible directions that most strongly affect $W_0+\Delta W$.
% Under LoRA, where optimization is restricted to a low-dimensional structured subspace, this curvature-aware control is easier to realize and has larger consequences: it simultaneously concentrates task-relevant curvature into the adapter-induced subspace and steers the effective model toward flatter task-conditioned regions.


% Consequently, at a fixed PECL scope, perturbation \emph{type} determines the perturbation geometry within LoRA coordinates, which in turn determines the learned updates $\Delta W_t$ and reshapes the \emph{full-model} loss landscape of $W_0+\Delta W_t$, ultimately driving the observed performance ordering.

% In particular, compared with SGD and the other perturbation types, $\mathrm{AS}^{(1)}$ exhibits a markedly stronger advantage in our LoRA-based PECL setting.
% This gap is consistently visible across multiple levels of evidence: $\mathrm{AS}^{(1)}$ achieves the best class-incremental accuracy and the smallest sharpness and curvature diagnostics on the trainable LoRA parameters in Table\ref{table:q1_ci_with_weight_flatness_reorg}, the most flat lossland in Fig\ref{fig:perturb_type_landscape}. and a more favorable accuracy trajectory over the entire task sequence in Fig\ref{fig:robustness_imagenet_c}. Notably, such a pronounced separation is not as evident in the originalstudy under standard full fine-tuning \cite{Zhang_GAM}, which suggests that the benefit of first-order sharpness can be substantially amplified under subspace-constrained adaptation.
% To figure it out, we compare $\mathrm{AS}^{(1)}$ with  $\mathrm{AS}^{(0)}$ and SGD in basic seqLoRA setting. Among all methods, only $\mathrm{AS}^{(1)}$ drives a persistent increase in the curvature localization ratio $c_t$ and $\alpha_t$, coincides with the highest and most stable $\operatorname{Acc}_t$. This indicates that $\mathrm{AS}^{(1)}$ not only concentrates dominant task-conditioned curvature more strongly into the adapter-induced subspace, but also drives optimization toward flatter task-conditioned solutions in the effective model $W_0+\Delta W$. Consequently, $\mathrm{AS}^{(1)}$ learns updates that are both more task-specific and better regularized in the full-model geometry, which is consistent with its consistently improved continual-learning performance. The influence of local curvature is easier to control, and therefore yields larger gains, when adaptation is confined to a small set of trainable coordinates. Under LoRA, where optimization is restricted to a low-dimensional, structured subspace, curvature-aware updates can be aligned with the few admissible high-impact directions, making curvature shaping both feasible and consequentia These results show that, under a fixed PECL scope, the choice of perturbation \emph{type} controls the perturbation geometry within the LoRA coordinates, which in turn determines the resulting $\Delta W_t$ and reshapes the \emph{full-model} landscape of $W_0+\Delta W_t$



% , indicating that the dominant task-conditioned curvature becomes progressively more concentrated in the adapter-induced subspace where $\Delta W_t$ operates.
% At the same time, $\mathrm{AS}^{(1)}$ maintains the max amplification factor $\alpha_t$ throughout training, sugggetings it learen task 


% compared with SGD, $\mathrm{AS}^{(1)}$ yields a larger and steadily increasing curvature localization ratio $c_t$, together with a noticeably flatter loss profile under full-parameter evaluation. This indicates that $\mathrm{AS}^{(1)}$ not only concentrates dominant task-conditioned curvature more strongly into the adapter-induced subspace, but also drives optimization toward flatter task-conditioned solutions in the effective model $W_0+\Delta W$. Consequently, $\mathrm{AS}^{(1)}$ learns updates that are both more task-specific and better regularized in the full-model geometry, which is consistent with its consistently improved continual-learning performance.
% These results show that, under a fixed PECL scope, the choice of perturbation \emph{type} controls the perturbation geometry within the LoRA coordinates, which in turn determines the resulting $\Delta W_t$ and reshapes the \emph{full-model} landscape of $W_0+\Delta W_t$.


% \noindent\textbf{why }\quad

% \noindent\textbf{Interaction with LoRA organization.}\quad
% We further investigate how the effect of perturbation {type} varies with the LoRA organization. 
% As shown in Table~\ref{table:q1_ci_with_weight_flatness_reorg} and  Fig.~\ref{fig:robustness_imagenet_cp}, the ordering of perturbation types is stable across LoRA organizations, but the performance gains vary in magnitude, indicating that perturbation geometry interacts with the organization-defined admissible update directions.
% In particular, SeqLoRA remains more robust on ImageNet-C/P, whereas IncLoRA and OLoRA exhibit larger accuracy drops under shift.
% Since the sharpness term is defined on the task adapter subspace, it is inherently sensitive to how different LoRA organization shapes the admissible scope $\mathcal W_{\Delta,t}$ across tasks. When $\mathcal W_{\Delta,t}$ is larger or expands over time, perturbations can probe a broader set of admissible directions, including more directions with non-negligible projected curvature.


\subsection{Ablation study on task length and adapter rank}
\label{sec:results-ablation}


{%
\setlength{\intextsep}{4pt}        % [h] 浮动体与正文的上下间距（最关键）
\setlength{\textfloatsep}{4pt}     % [t]/[b] 时浮动体与正文间距（备用）
\setlength{\abovecaptionskip}{1pt} % 图与caption之间（上方）间距
\setlength{\belowcaptionskip}{1pt} % caption下方到正文的间距
\captionsetup{skip=3pt}            % 你已在用：图与caption的距离（更细粒度）
\begin{figure}[t]
    \centering
\includegraphics[width=1\linewidth]{figures_TRML/ablation_rank_length_imagenetr_NME_T.pdf}
\vspace{-10pt}
    \caption{ Sensitivity to Adapter Capacity and Sequence Difficulty
    % Ablation on adapter rank and task sequence length on ImageNet-R.
% We report the final NME-based accuracy $\operatorname{Acc}_T^{\mathrm{nme}}$ while varying (a) the adapter rank $r$ and (b) the task length $T$.
    }
    \label{fig:ablation}
\end{figure}
}



Ablation on adapter rank and task sequence length on ImageNet-R.
We study how adapter capacity and sequence difficulty affect PECL by varying the LoRA rank $r$ and the task sequence length $T$, and report the final NME accuracy $\operatorname{Acc}_T^{\mathrm{nme}}$ in Fig.~\ref{fig:ablation}.
Across all settings, the relative ordering is stable: flatness-aware optimizers consistently outperform SGD.
Accuracy improves from small $r$ but quickly saturates at moderate ranks. This supports our theoretical view that PECL primarily requires a low-dimensional task-permissible update subspace to capture task-specific directions, so enlarging $r$ beyond this intrinsic dimension adds capacity with limited benefit.
%while leaving the bound-relevant geometry largely unchanged




\section{Conclusion}

We develop a sequential hierarchical PAC-Bayes framework for LoRA-based PECL, showing that both the sharpness-dependent empirical term and the KL complexity are supported on the task-permissible subspace. Guided by the bound, we decomposed perturbation design into scope and type and showed empirically that subspace-aligned perturbations, especially those emphasizing curvature-aware directions, reduce perturbation-aware risk along the task sequence and translate into higher continual accuracy.

% In LoRA-Based CL, sharpness should be defined on the task-admissible LoRA subspace, the perturbation type modulates the strength of subspace-aligned flatness control and its associated performance gains.






\section*{Impact Statement}

% Authors are \textbf{required} to include a statement of the potential broader
% impact of their work, including its ethical aspects and future societal
% consequences. This statement should be in an unnumbered section at the end of
% the paper (co-located with Acknowledgements -- the two may appear in either
% order, but both must be before References), and does not count toward the paper
% page limit. In many cases, where the ethical impacts and expected societal
% implications are those that are well established when advancing the field of
% Machine Learning, substantial discussion is not required, and a simple
% statement such as the following will suffice:

% ``This paper presents work whose goal is to advance the field of Machine
% Learning. There are many potential societal consequences of our work, none
% which we feel must be specifically highlighted here.''

% The above statement can be used verbatim in such cases, but we encourage
% authors to think about whether there is content which does warrant further
% discussion, as this statement will be apparent if the paper is later flagged
% for ethics review.

% This paper presents work whose goal is to advance the field of Machine
% Learning and Continual Learning. 
% \section*{Impact Statement}
The rapid growth in the size of pre-trained models has driven a surge in parameter-efficient fine-tuning methods, among which LoRA is the most widely used. In this work, we study the impact of optimizing for flat minima in this setting. We show that flat-minima optimization can significantly reduce forgetting in a continual learning scenario, helping maintain performance on previously learned tasks while adapting efficiently to new ones. We both improve the theoretical understanding of this phenomenon and provide ample empirical results that support it.
% The work is primarily foundational and does not introduce new data sources or application-specific deployment pipelines. 
% Potential downstream impacts depend on how continual learning systems are applied; as with other modeling advances, misuse in sensitive domains could amplify existing risks if deployed without appropriate validation, monitoring, and governance.

% There are many potential societal consequences of our work, none
% which we feel must be specifically highlighted here.

% In the unusual situation where you want a paper to appear in the
% references without citing it in the main text, use \nocite
% \nocite{langley00}

\bibliography{liying.bib}
\bibliographystyle{icml2026}

%%%%%%%%%%%%%%%%%%%%%%%%%%%%%%%%%%%%%%%%%%%%%%%%%%%%%%%%%%%%%%%%%%%%%%%%%%%%%%%
%%%%%%%%%%%%%%%%%%%%%%%%%%%%%%%%%%%%%%%%%%%%%%%%%%%%%%%%%%%%%%%%%%%%%%%%%%%%%%%
% APPENDIX
%%%%%%%%%%%%%%%%%%%%%%%%%%%%%%%%%%%%%%%%%%%%%%%%%%%%%%%%%%%%%%%%%%%%%%%%%%%%%%%
%%%%%%%%%%%%%%%%%%%%%%%%%%%%%%%%%%%%%%%%%%%%%%%%%%%%%%%%%%%%%%%%%%%%%%%%%%%%%%%
\newpage
\appendix
\onecolumn


% \section{You \emph{can} have an appendix here.}

% You can have as much text here as you want. The main body must be at most $8$
% pages long. For the final version, one more page can be added. If you want, you
% can use an appendix like this one.

% The $\mathtt{\backslash onecolumn}$ command above can be kept in place if you
% prefer a one-column appendix, or can be removed if you prefer a two-column
% appendix.  Apart from this possible change, the style (font size, spacing,
% margins, page numbering, etc.) should be kept the same as the main body.
%%%%%%%%%%%%%%%%%%%%%%%%%%%%%%%%%%%%%%%%%%%%%%%%%%%%%%%%%%%%%%%%%%%%%%%%%%%%%%%
%%%%%%%%%%%%%%%%%%%%%%%%%%%%%%%%%%%%%%%%%%%%%%%%%%%%%%%%%%%%%%%%%%%%%%%%%%%%%%%
%\section{Appendix}

\appendix

\section{Scope and Summary of Notation}
\label{app:notation}

This appendix collects the main symbols used in the paper and in the appendices, following (approximately) the order in which they appear.

\begin{itemize}
\item $t\in\{1,\dots,T\}$: task index in the continual learning sequence. % optional
% \item $T$: total number of tasks.
\item $S_t$: training set of task $t$.
\item $m_t := |S_t|$: number of training samples in $S_t$.
\item $D_t$: data distribution of task $t$.
\item $z=(x,y)\sim D_t$: an input label pair drawn from $D_t$.
% \item $f_W$: predictor (model) parameterized by $W$.
\item $\ell(W;z)$: bounded loss evaluated at parameters $W$ on sample $z$.
\item $L_{D_t}(W)=\mathbb E_{z\sim D_t}[\ell(W;z)]$: population risk on task $t$. 
\item $\widehat L_{S_t}(W)=\frac{1}{m_t}\sum_{z\in S_t}\ell(W;z)$: empirical risk on $S_t$.

% \item $W\in\mathbb R^{d}$: full parameter vector (or a vectorized weight block) of the model.

% \item $d$: dimensionality of the ambient parameter (or increment) space.


\item $\rho$: perturbation radius for adversarial sharpness.
% \item $\epsilon$: adversarial perturbation with $\|\epsilon\|\le \rho$.

% \item $\varepsilon\sim\mathcal N(0,\Sigma)$: Gaussian perturbation variable.
\item $\epsilon$ and $\varepsilon$: adversarial and Gaussian perturbations, respectively, where $\|\epsilon\|\le \rho$ and $\varepsilon\sim\mathcal N(0,\Sigma)$.

% \item $\Sigma\succeq 0$: covariance matrix for Gaussian perturbations.
% \item $\varepsilon\sim\mathcal N(0,\Sigma_t)$: Gaussian perturbation used to define expected (random) sharpness. 
% \item $\nabla_W \hat L_S(W)$: gradient of empirical risk with respect to $W$.
% \item $\mathrm{AS}_S^{(0)}(W,\rho)$: adversarial zeroth-order sharpness.
% \item $\mathrm{AS}_S^{(1)}(W,\rho)$: adversarial first-order sharpness.
\item $\mathrm{AS}^{(0)}_S(W,\rho)$: adversarial zeroth-order sharpness from SAM,
$\mathrm{AS}^{(0)}_S(W,\rho)=\max_{\|\epsilon\|\le\rho}\big[\widehat L_S(W+\epsilon)-\widehat L_S(W)\big]$.

\item $\mathrm{RS}_S^{(0)}(W,\Sigma)$: random zeroth-order sharpness, i.e., the expected empirical loss elevation under $\varepsilon\sim\mathcal N(0,\Sigma)$:
$\mathrm{RS}_S^{(0)}(W,\Sigma)=\mathbb E_{\varepsilon\sim\mathcal N(0,\Sigma)}[\hat L_S(W+\varepsilon)-\hat L_S(W)]$.
\item $\mathrm{AS}^{(1)}_S(W,\rho)$: adversarial first-order sharpness  from GAM,
$\mathrm{AS}^{(1)}_S(W,\rho)=\max_{\|\epsilon\|\le\rho}\|\nabla_W\widehat L_S(W+\epsilon)\|$.

\item $P$: prior distribution over parameters (hypotheses) in PAC-Bayes analysis.
\item $Q$: posterior distribution over parameters (hypotheses) in PAC-Bayes analysis.
% \item $\mathrm{KL}(Q\Vert P)$: Kullback--Leibler divergence measuring posterior complexity relative to the prior.
% \item $\delta\in(0,1]$: confidence level in high-probability bounds. 

\item $P_t$: task prior at step $t$ 
% (a distribution over model parameters restricted to the current task). 
\item $Q_t(\cdot \mid P_t,S_t)$: task posterior after learning task $t$, possibly conditioned on $P_t$ and $S_t$.
\item $Q_t^\Delta \ll P_t^\Delta$: absolute continuity assumption
\item $\mathcal P_t$: hyper-posterior (a distribution over task priors), forming a time-varying hyper-process . 
\item $\mathcal F_{t-1}$: filtration ($\sigma$-algebra) capturing the history up to step $t-1$ in the martingale argument.

\item $W_0$: pre-trained initialization before continual adaptation.
\item $W_t$: model parameters after training on tasks $1$ to $t$.
\item $W_t^{\mathrm{froz}}$: frozen component when learning task $t$ (the part not trainable during task $t$).
\item $\Delta W_t$: trainable update at task $t$ under parameter-efficient adaptation.
\item $W_t = W_t^{\mathrm{froz}} + \Delta W_t$: decomposition of current parameters into frozen and trainable parts. 





% \item $\mathrm{LoRA}$: low-rank adaptation parameterization of a weight block, typically $\Delta W_t = B_tA_t$ with rank $r$.
% \item $A_t\in\mathbb R^{r\times d_{\text{in}}},\; B_t\in\mathbb R^{d_{\text{out}}\times r}$: LoRA factors for task $t$ (shapes may vary by layer/block).
% \item $r$: LoRA rank (adapter dimension).


\item $\mathcal W_{\Delta,t}\subseteq\mathbb R^d$: the task-admissible linear update subspace at step $t$, containing all feasible increments induced by LoRA while the backbone is frozen.


\item $\widehat W_t$: the full-model parameter vector after learning task $t$ (a point estimate in the ambient parameter space).
\item $\widehat{\Delta W}_t$: the learned trainable increment at task $t$ under PEFT/LoRA, satisfying $\widehat W_t = W_t^{\mathrm{froz}} + \widehat{\Delta W}_t$.

 


 
% \item $\mathcal B(\cdot)$: Borel sigma-algebra (and its trace on subspaces/affine sets). 
% \item $T_t(\Delta)=W_t^{\mathrm{froz}}+\Delta$: translation map from the update space to the affine model class. 
% \item $T_{t\#}\mu$: pushforward measure induced by $T_t$, defined by $(T_{t\#}\mu)(A)=\mu(T_t^{-1}(A))$. 
\item $P_t^\Delta,\;Q_t^\Delta$: prior/posterior supported on $(\mathcal W_{\Delta,t},\mathcal B(\mathcal W_{\Delta,t}))$.
% \item $P_,\; Q_t$: induced full-model prior/posterior supported on $W_t^{\mathrm{froz}}+\mathcal W_{\Delta,t}$.
% \item $Q_t^\Delta \ll P_t^\Delta$: absolute continuity assumption.
% \item $\frac{dQ}{dP}$: Radon--Nikodym derivative appearing in the KL invariance proof. 

% \item $\Sigma_t$: covariance defining a Gaussian perturbation distribution on $\mathcal W_{\Delta,t}$. 



\item $\mathrm{RS}_t^\Delta(\widehat W_t,\Sigma_t)$: subspace sharpness with perturbations supported on $\mathcal W_{\Delta,t}$, defined as
$\mathrm{RS}_t^\Delta(\widehat W_t,\Sigma_t)=\mathbb E_{\varepsilon\sim\mathcal N(0,\Sigma_t)}[\hat L_{S_t}(\widehat W_t+\varepsilon)-\hat L_{S_t}(\widehat W_t)]$,
where $\varepsilon$ is restricted to $\mathcal W_{\Delta,t}$, $\Sigma_t$ is defined on $\mathcal{W}_{\Delta,t}$
% (equivalently, $\mathrm{supp}(\varepsilon)\subseteq\mathcal W_{\Delta,t}$).


% \item $\widehat W_t$: posterior mean (or a point estimate) used when specializing $Q_t$ to a Gaussian.
% \item $\mathrm{RS}_t^\Delta(\widehat W_t,\Sigma_t)$:  random sharpness(expected sharpness) restricted to $\mathcal W_{\Delta,t}$ .

% \item $a_{tj}$: accuracy on task $j$ after training up to task $t$.
% \item $\mathrm{Acc}_t=\frac{1}{t}\sum_{j=1}^t a_{tj}$: average accuracy across the first $t$ tasks.
% \item $\mathrm{Acc}_t^{\mathrm{cls}}$: class-incremental classification accuracy at step $t$.
% \item $\mathrm{Acc}_t^{\mathrm{ncm}}$: nearest-class-mean accuracy (prototype-based) reflecting representation quality. 

\item $\alpha_t$ : feature amplification factor

% \item $c_t$: curvature projection ratio 

% \item $\lambda>0$: exponential tilting parameter used in the moment generating function (MGF) step of the proof~\ref{proof:theorem all}. 
% \item $B_T:=\sum_{t=1}^T (m_t-1)$: aggregate sample-size factor in Proof~\ref{proof:theorem all}. 
% \item $\omega_t := (m_t-1)/B_T$: sample-size weight used to form the weighted generalization gap.
% \item $\Delta_t$: step-$t$ generalization gap (defined in ~\ref{proof:theorem all}).
% \item $\overline\Delta := \sum_{t=1}^T \alpha_t \Delta_t$: size-weighted average gap. 
% \item $g_t(W):=\lambda\alpha_t\big(L_{D_t}(W)-\widehat L_{S_t}(W)\big)$: MGF test function at step $t$. 
% \item $h_t(P)$: hyper-level functional used in the Donsker--Varadhan (DV) change-of-measure step. 
% \item $\mathcal F_{t-1}$: filtration ($\sigma$-algebra) capturing the history up to step $t-1$ in the martingale argument. 
% \item $X_t$: exponential increment defined by the conditional MGF bound at step $t$.
% \item $\mathcal M_t:=\prod_{s=1}^t X_s$: nonnegative supermartingale used to obtain a high-probability bound.

% \item $\mathcal W_{\Delta,t}\subseteq\mathbb R^d$: permissible update subspace at task $t$ (task-admissible increment space). 
% \item $\Pi_{\Delta,t}$: orthogonal projector onto $\mathcal W_{\Delta,t}$ when $\mathcal W_{\Delta,t}$ is a linear subspace. 
% \item $J_{\mathrm{LoRA}}$: Jacobian of the LoRA map (used to define the local tangent increment space), with $\mathcal W_{\Delta,t}=\mathrm{Im}(J_{\mathrm{LoRA}})$.


% \item $\Delta W_t=U_{\Delta,t}S_{\Delta,t}V_{\Delta,t}^\top$: rank-$r$ truncated SVD of the learned update . 
% \item $\mathcal S_{\Delta,t}=\{U_{\Delta,t}ZV_{\Delta,t}^\top: Z\in\mathbb R^{r\times r}\}$: bilinear diagnostic subspace induced by the truncated SVD of $\Delta W_t$. 

% % \item $\Pi_{\Delta,t}^{\mathrm{SVD}}(M):=U_{\Delta,t}(U_{\Delta,t}^\top M V_{\Delta,t})V_{\Delta,t}^\top$: Frobenius-orthogonal projection onto $\mathcal S_{\Delta,t}$ (used only for curvature localization). 
% \item $\Pi_{\mathcal S,t}(M):=U_{\Delta,t}(U_{\Delta,t}^\top M V_{\Delta,t})V_{\Delta,t}^\top$: Frobenius-orthogonal projection onto $\mathcal S_{\Delta,t}$ .


% % \item $(U_{W,t},V_{W,t})$: rank-$r$ truncated SVD subspace of the corresponding backbone weight block (Appendix A.4.2). 
% % \item $\mathcal S_{W,t}=\{U_{W,t}ZV_{W,t}^\top: Z\in\mathbb R^{r\times r}\}$: backbone reference bilinear subspace. 
% % \item $\Pi_{W,t}^{\mathrm{SVD}}(M):=U_{W,t}(U_{W,t}^\top M V_{W,t})V_{W,t}^\top$: projection onto $\mathcal S_{W,t}$. 

% \item $(U_{W,t},V_{W,t})$: rank-$r$ truncated SVD bases of the corresponding backbone weight block.
% \item $\mathcal S_{W,t}=\{U_{W,t}ZV_{W,t}^\top: Z\in\mathbb R^{r\times r}\}$: backbone reference bilinear subspace.
% \item $\Pi_{\mathcal S_W,t}(M):=U_{W,t}(U_{W,t}^\top M V_{W,t})V_{W,t}^\top$: Frobenius-orthogonal projection onto $\mathcal S_{W,t}$.

% \item $\{(\lambda_{t,i},v_{t,i})\}_{i=1}^k$: top-$k$ eigenpairs of a PSD curvature operator (e.g., GGN/Fisher) at task $t$. 
% \item $(\delta W^{(0)}_{t,i},\delta A_{t,i},\delta B_{t,i})$: decomposition of eigen-direction $v_{t,i}$ into backbone and LoRA-factor components. 
% \item $\delta W_{t,i}=\delta W^{(0)}_{t,i}+(\delta B_{t,i}A_t+B_t\delta A_{t,i})$: effective weight-space perturbation induced by $(\delta W^{(0)}_{t,i},\delta A_{t,i},\delta B_{t,i})$. 
% \item $\rho_{\Delta,t,i}=\frac{\|\Pi^{\mathrm{SVD}}_{\Delta,t}(\delta W_{t,i})\|_F^2}{\|\delta W_{t,i}\|_F^2}$: per-mode projection ratio onto $\mathcal S_{\Delta,t}$.
% \item $\rho_{W,t,i}=\frac{\|\Pi_{W,t}^{\mathrm{SVD}}(\delta W_{t,i})\|_F^2}{\|\delta W_{t,i}\|_F^2}$: per-mode projection ratio onto $\mathcal S_{W,t}$.
\item $R_{\Delta,t}=\frac{\sum_{i=1}^k\lambda_{t,i}\rho_{\Delta,t,i}}{\sum_{i=1}^k\lambda_{t,i}}$: eigenvalue-weighted localization on $\mathcal S_{\Delta,t}$.
\item $R_{W,t}=\frac{\sum_{i=1}^k\lambda_{t,i}\rho_{W,t,i}}{\sum_{i=1}^k\lambda_{t,i}}$: eigenvalue-weighted localization on $\mathcal S_{W,t}$.
\item $c_t=R_{\Delta,t}/R_{W,t}$: curvature projection ratio comparing $\mathcal S_{\Delta,t}$ against the backbone reference. 
\end{itemize}


\section{Proofs}
In this appendix we provide complete proofs of the main theorems and lemmas from the main paper.

\subsection{Proof of Theorem \ref{thm:pathwise_pac}}
\label{proof:theorem all}

\paragraph{Preliminary: a supermartingale tail bound.}
Let $(\mathcal M_t)_{t\ge0}$ be a nonnegative supermartingale adapted to $(\mathcal F_t)$ with $\mathbb E[\mathcal M_0]\le 1$.
Then for any $\delta\in(0,1)$,
\[
\mathbb P\!\left(\mathcal M_T \ge \tfrac{1}{\delta}\right)
\le \delta,
\]
by Markov's inequality and $\mathbb E[\mathcal M_T]\le \mathbb E[\mathcal M_0]\le 1$.
Equivalently, with probability at least $1-\delta$, $\log \mathcal M_T \le \log \tfrac{1}{\delta}$.
This is the fixed-time form of the classical supermartingale maximal inequality  \citet{yangshuxue}.


\begin{theorem}[Pathwise PAC--Bayes with time-varying hyper-posterior]
\label{the: full}
For any $\delta\in(0,1]$, with probability at least $1-\delta$ over $S_{1:T}$,
\begin{equation}
\begin{aligned} 
& \frac{1}{T} \sum_{t=1}^T \underbrace{\mathbb{E}_{P_t \sim \mathcal{P}_t} \mathbb{E}_{W \sim Q_t(\cdot \mid P_t, S_t)} L_{\mathcal{D}_t}(W)}_{\text{test risk at step } t} 
\le \ \frac{1}{T} \sum_{t=1}^T \underbrace{\mathbb{E}_{P_t \sim \mathcal{P}_t} \mathbb{E}_{W \sim Q_t(\cdot \mid P_t, S_t)} \hat{L}_{S_t}(W)}_{\text{empirical risk at step } t} \\ 
& + \sqrt{\frac{ \sum_{t=1}^T \mathbb{E}_{P_t \sim \mathcal{P}_t} \mathrm{KL}\big(Q_t(\cdot \mid P_t, S_t) \,\|\, P_t\big) + \sum_{t=1}^T \mathrm{KL}(\mathcal{P}_t \,\|\, \mathcal{P}_{t-1}) + \log\frac{1}{\delta} }{2 \sum_{t=1}^T (m_t - 1)}}. 
\end{aligned} 
\end{equation}
\end{theorem}

\begin{proof}
    


% \subsection*{A.1\quad Two-level variational inequalities}
% \item $B_T:=\sum_{t=1}^T (m_t-1)$: aggregate sample-size factor in Proof~\ref{proof:theorem all}. 
% \item $\omega_t := (m_t-1)/B_T$: sample-size weight used to form the weighted generalization gap.
Let $B_T=\sum_{t=1}^T (m_t-1)$ denote an aggregate sample-size factor and define $\omega_t=(m_t-1)/B_T$ which serves as a sample-size weight satisfying $\sum_{t=1}^T\omega_t=1$.
% $B_T$ is 

% \paragraph{Goal quantity.} 


Define the step-$t$ generalization gap
\[
\Delta_t
:=\mathbb E_{P_t\sim\mathcal P_{t-1}}\,
\mathbb E_{W\sim Q_t(\cdot\mid P_t,S_t)}
\big(L_{\mathcal D_t}(W)-\hat L_{S_t}(W)\big),
\]
and the size-weighted gap $\bar\Delta:=\sum_{t=1}^T \omega_t \Delta_t$.
Fix $\lambda>0$ and set
\(
g_t(W):=\lambda\omega_t\big(L_{\mathcal D_t}(W)-\hat L_{S_t}(W)\big).
\)

\paragraph{Step 1 Model level.}
% For any fixed $P$ and $S_t$,
For any fixed prior distribution over the parameters $P$ and posterior distribution $Q$, and only current task dataset $S_t$,   
\begin{equation}
\label{eq:dv-w}
\mathbb E_{W\sim Q_t(\cdot\mid P,S_t)} g_t(W)
\;\le\;
\mathrm{KL}\!\big(Q_t(\cdot\mid P,S_t)\,\|\,P\big)
+\log\mathbb E_{W\sim P}\!\big[e^{g_t(W)}\big].
\end{equation}
Taking $\mathbb E_{S_t}$ and using the standard Hoeffding--MGF bound for bounded losses,
\begin{equation}
\label{eq:mgf-mean-gap}
\mathbb E_{S_t}\log\mathbb E_{W\sim P} e^{g_t(W)}
\;\le\; \frac{(\lambda\omega_t)^2}{8\,(m_t-1)}.
\end{equation}
Therefore,
\begin{equation}
\label{eq:model-step}
\mathbb E_{S_t}\mathbb E_{W\sim Q_t(\cdot\mid P,S_t)} g_t(W)
\;\le\;
\mathbb E_{S_t}\mathrm{KL}\!\big(Q_t(\cdot\mid P,S_t)\,\|\,P\big)
+\frac{(\lambda\omega_t)^2}{8\,(m_t-1)}.
\end{equation}
 % (change of measure in $P$-space)
\paragraph{Step 2 Hyper level.} 
Define
\[
h_t(P):=\mathbb E_{S_t}\mathbb E_{W\sim Q_t(\cdot\mid P,S_t)} g_t(W)
-\mathbb E_{S_t}\mathrm{KL}\!\big(Q_t(\cdot\mid P,S_t)\,\|\,P\big).
\]
By \eqref{eq:model-step}, $h_t(P)\le \frac{(\lambda\omega_t)^2}{8\,(m_t-1)}$ for all $P$.
Applying Donsker–Varadhan variational representation of the KL divergence~\cite{NEURIPS2021_54a367d6} to change the measure from $\mathcal P_{t-1}$ to $\mathcal P_t$ gives
\begin{equation}
\label{eq:dv-p}
\log\mathbb E_{P\sim\mathcal P_t}\!\big[e^{h_t(P)}\big]
\;\le\; \mathrm{KL}(\mathcal P_t\|\mathcal P_{t-1})
+\log\mathbb E_{P\sim\mathcal P_{t-1}}\!\big[e^{h_t(P)}\big]
\;\le\; \mathrm{KL}(\mathcal P_t\|\mathcal P_{t-1})
+\frac{(\lambda\omega_t)^2}{8\,(m_t-1)}.
\end{equation}
Using $\log\mathbb E[e^X]\ge \mathbb E[X]$ yields
\begin{equation}
\label{eq:p-expect}
\mathbb E_{P_t\sim\mathcal P_t} h_t(P_t)
\;\le\; \mathrm{KL}(\mathcal P_t\|\mathcal P_{t-1})
+\frac{(\lambda\omega_t)^2}{8\,(m_t-1)}.
\end{equation}
Unfolding $h_t(P)$ and interchanging expectations,
\begin{equation}
\label{eq:two-term}
\mathbb E_{S_t}\Big[
\lambda\omega_t\Delta_t
-\omega_t\,\mathbb E_{P_t\sim\mathcal P_{t-1}}\mathrm{KL}\!\big(Q_t(\cdot\mid P_t,S_t)\,\|\,P_t\big)
\Big]
\;\le\;
\mathrm{KL}(\mathcal P_t\|\mathcal P_{t-1})
+\frac{(\lambda\omega_t)^2}{8\,(m_t-1)}.
\end{equation}

% \subsection*{A.2\quad Conditional MGF and supermartingale}
\paragraph{Step 3 Martingale}
% From \eqref{eq:two-term} and the tower property of conditional expectation, 







Using Eq.\eqref{eq:two-term} and the tower property of conditional expectation (iterated expectation) with respect to $\mathcal F_{t-1}$, we obtain the desired conditional exponential bound~\cite{yangshuxue}.
% The conditional exponential bound
\begin{equation}
\label{eq:cond-mgf}
\mathbb E\!\left[
\exp\!\Big(
\lambda\omega_t\Delta_t
-\omega_t\,\mathbb E_{P_t\sim\mathcal P_{t-1}}\mathrm{KL}(Q_t\|P_t)
-\mathrm{KL}(\mathcal P_t\|\mathcal P_{t-1})
-\frac{(\lambda\omega_t)^2}{8\,(m_t-1)}
\Big)\,\Big|\,\mathcal F_{t-1}\right]\;\le\;1.
\end{equation}


From the conditional exponential bound~\eqref{eq:cond-mgf}, define
\[
X_t
:=
\exp\!\Big(
\lambda\omega_t\Delta_t
-\omega_t\,\mathbb E_{P_t\sim\mathcal P_{t-1}}\mathrm{KL}(Q_t\|P_t)
-\mathrm{KL}(\mathcal P_t\|\mathcal P_{t-1})
-\frac{(\lambda\omega_t)^2}{8\,(m_t-1)}
\Big).
\]
Assume $\omega_t$ is deterministic or $\mathcal F_{t-1}$-measurable. Then $X_t\ge 0$ is $\mathcal F_t$-measurable and accoring Eq.~\eqref{eq:cond-mgf}
\[
\mathbb E[X_t\mid \mathcal F_{t-1}] \le 1.
\]
Define the process
\[
\mathcal{M}_0:=1,\qquad
\mathcal{M}_t:=\prod_{s=1}^t X_s.
\]
It follows that $(\mathcal M_t)_{t\ge 0}$ is a nonnegative supermartingale, since
\[
\mathbb E[\mathcal M_t\mid \mathcal F_{t-1}]
=
\mathcal M_{t-1}\,\mathbb E[X_t\mid \mathcal F_{t-1}]
\le \mathcal M_{t-1}.
\]
In particular, $\mathbb E[\mathcal M_T]\le \mathbb E[\mathcal M_0]=1$.
By Markov's inequality,
\[
\mathbb P\!\left(\mathcal M_T \ge \tfrac{1}{\delta}\right)
\le \delta,
\]
so with probability at least $1-\delta$ we have $\log \mathcal M_T \le \log\tfrac{1}{\delta}$.
% Expanding $\log \mathcal M_T = \sum_
Expanding $\log \mathcal M_T = \sum_{t=1}^T \log X_t$ yields~\eqref{eq:hpp}.






% Define $X_t$ to be the exponential increment inside \eqref{eq:cond-mgf} and
% \[
% \mathcal{M}_0:=1,\qquad
% \mathcal{M}_t:=\prod_{s=1}^t X_s.
% \]
% Then $(\mathcal{M}_t)_{t\ge0}$ is a nonnegative supermartingale, so
% $\mathbb E[\mathcal{M}_T]\le 1$.
% By Markov's inequality, with probability at least $1-\delta$,

% so with probability at least $1-\delta$ we have $\log \mathcal M_T \le \log\tfrac{1}{\delta}$

\begin{equation}
\label{eq:hpp}
\sum_{t=1}^T \lambda\omega_t\Delta_t
\;\le\;
\sum_{t=1}^T \omega_t\,\mathbb E_{P_t\sim\mathcal P_{t-1}}\mathrm{KL}(Q_t\|P_t)
\;+\; \sum_{t=1}^T \mathrm{KL}(\mathcal P_t\|\mathcal P_{t-1})
\;+\; \frac{\lambda^2}{8}\sum_{t=1}^T \frac{\omega_t^2}{(m_t-1)}
\;+\; \log\tfrac1\delta.
\end{equation}

% \subsection*{A.3\quad Size-weighted self-normalization and optimization}

By $\omega_t=(m_t-1)/B_T$,
\[
\sum_{t=1}^T \frac{\omega_t^2}{(m_t-1)}
=\sum_{t=1}^T \frac{(m_t-1)^2/B_T^2}{(m_t-1)}
=\frac{1}{B_T}.
\]
Hence, \eqref{eq:hpp} reads, for all $\lambda>0$ and with probability at least $1-\delta$,
\begin{equation}
\label{eq:gap-preopt}
\bar\Delta
\;\le\;
\frac{1}{\lambda}\Big(
\sum_{t=1}^T \omega_t\,\mathbb E_{P_t\sim\mathcal P_{t-1}}\mathrm{KL}(Q_t\|P_t)
+ \sum_{t=1}^T \mathrm{KL}(\mathcal P_t\|\mathcal P_{t-1})
+ \log\tfrac1\delta\Big)
\;+\;
\frac{\lambda}{8B_T}.
\end{equation}
Optimizing the right-hand side over $\lambda>0$ (take $\lambda^\star=\sqrt{8B_T\,A}$ with
$A:=\sum_{t=1}^T \omega_t\,\mathbb E_{P_t\sim\mathcal P_{t-1}}\mathrm{KL}(Q_t\|P_t)
+ \sum_{t=1}^T \mathrm{KL}(\mathcal P_t\|\mathcal P_{t-1}) + \log\tfrac1\delta$) yields
\begin{equation}
\label{eq:weighted-final}
\bar\Delta
\;\le\;
\sqrt{\frac{A}{2B_T}}
\;\le\;
\sqrt{\frac{
\sum_{t=1}^T \mathbb E_{P_t\sim\mathcal P_{t-1}}\mathrm{KL}(Q_t\|P_t)
+ \sum_{t=1}^T \mathrm{KL}(\mathcal P_t\|\mathcal P_{t-1})
+ \log\tfrac1\delta}{2\,\sum_{t=1}^T (m_t-1)}}.
\end{equation}
The last inequality uses $\sum_t \omega_t x_t \le \sum_t x_t$ for nonnegative $(x_t)$.

\paragraph{Step 4 From weighted to the statement of Theorem~\ref{the: full}.}
The theorem controls the unweighted average $\frac1T\sum_{t=1}^T \Delta_t$ and uses the smoothed empirical risk on the right-hand side.
When $m_t$ are homogeneous, $\omega_t\equiv 1/T$ and $\bar\Delta=\frac1T\sum_t\Delta_t$.
For heterogeneous $m_t$, we keep the same right-hand side and upper-bound the numerator $\sum_t \omega_t \mathbb E_{P_t}\mathrm{KL}(Q_t\|P_t)$ by $\sum_t \mathbb E_{P_t}\mathrm{KL}(Q_t\|P_t)$, which yields exactly the denominator $2\sum_t(m_t-1)$ in the theorem.

% \subsection*{A.4\quad Gaussian smoothing (LoRA) corollary}

If $Q_t(\cdot\mid P_t,S_t)=\mathcal N(\hat W_t,\Sigma_t)$ (e.g., supported on an task-permissible subspace), then
\[
\mathbb E_{W\sim Q_t(\cdot\mid P_t,S_t)}\hat L_{S_t}(W)
=\mathbb E_{\varepsilon\sim\mathcal N(0,\Sigma_t)}\hat L_{S_t}(\hat W_t+\varepsilon),
\]
which converts the empirical term into the smoothed empirical loss (expected sharpness) used in the main text; substituting this identity into \eqref{eq:weighted-final} yields the empirical term in Theorem~\ref{the: full}.

Theorem~A.1 provides a tight complexity penalty of the form.
Using $\sqrt{x+y+z}\le \sqrt{x}+\sqrt{y}+\sqrt{z}$ for $x,y,z\ge 0$ yields the decomposed bound, which is the form stated in Theorem~\ref{thm:pathwise_pac} for interpretability.

\end{proof}

\subsection{Proof of Lemma~\ref{lem:kl-decomp-pecl}}
\label{proof:Lemma1}

Let $\mathcal W_{\Delta,t}\subseteq\mathbb R^d$ be a linear subspace and $W_t^\mathrm{froz}+\mathcal W_{\Delta,t}$ be its affine translation. 
We endow both sets with the trace Borel $\sigma$-algebra induced from $\mathbb R^d$. 
Define the translation map
\[
T_t:(\mathcal W_{\Delta,t},\mathcal B(\mathcal W_{\Delta,t}))
\to (W_t^{\mathrm{froz}}+\mathcal W_{\Delta,t},\mathcal B(W_t^{\mathrm{froz}}+\mathcal W_{\Delta,t})),
\qquad
T_t(\Delta)=W_t^{\mathrm{froz}}+\Delta .
\]
For any probability measure $\mu$ on $\mathcal W_{\Delta,t}$, its pushforward
$T_{t\#}\mu$ is the measure on $W_t^{\mathrm{froz}}+\mathcal W_{\Delta,t}$
defined by $(T_{t\#}\mu)(A)=\mu(T_t^{-1}(A))$ for all measurable $A$.
Accordingly, we define the induced full-model prior and posterior supported on
$W_t^{\mathrm{froz}}+\mathcal W_{\Delta,t}$ as
\[
P_t := T_{t\#} P_t^\Delta,
\qquad
Q_t := T_{t\#} Q_t^\Delta,
\]
where $P_t^\Delta$ and $Q_t^\Delta$ are probability measures on
$(\mathcal W_{\Delta,t},\mathcal B(\mathcal W_{\Delta,t}))$.
% $\mathcal B(\mathcal W_{\Delta,t})\ =\{B\cap \mathcal W_{\Delta,t}:\ B\in \mathcal B(\mathbb R^d)\},\quad
% \mathcal B(W^{\mathrm{froz}}_t+\mathcal W_{\Delta,t})\ =\{B\cap (W^{\mathrm{froz}}_t+\mathcal W_{\Delta,t}):\ B\in \mathcal B(\mathbb R^d)\}.$
% These coincide with the Borel $\sigma$-algebras generated by the subspace topologies on $\mathcal W_{\Delta,t}$ and $W^{\mathrm{froz}}_t+\mathcal W_{\Delta,t}$, respectively, so both are Borel spaces.
% To formalize the probabilistic relationship between the LoRA subspace and the full model space, we define the translation map:
% % \begin{align*}
% $T_t: (\mathcal W_{\Delta,t},\mathcal B(\mathcal W_{\Delta,t})) \to (W^{\mathrm{froz}}_t+\mathcal W_{\Delta,t},\mathcal B(W^{\mathrm{froz}}_t+\mathcal W_{\Delta,t})), 
% T_t(\Delta)=W^{\mathrm{froz}}_t+\Delta,$
% % \end{align*}
% which is a Borel isomorphism between the subspace and affine space.  The full-space prior and posterior are defined as pushforwards onto the affine subspace by
% $
% P=(T_t)_{\text{\#}} P^{\Delta}_t,
% Q=(T_t)_{\text{\#}} Q^{\Delta}_t,
% $
% where $P_t^{\Delta}$ and $Q_t^{\Delta}$ are probability measures on $(\mathcal{W}_{\Delta,t}, \mathcal{B}(\mathcal{W}_{\Delta,t}))$.

% \begin{lemma}[KL decomposition with a frozen backbone]
% \label{equation: subspace-full}
% Let $W\in\mathbb R^d$ be fixed and let $P^{\Delta},Q^{\Delta}$ be probability measures on $(\mathcal W_\Delta,\mathcal B(\mathcal W_\Delta))$ with $Q^{\Delta}\ll P^{\Delta}$. Define their pushforwards onto the affine subspace by
% \[
% P:=(T_W)_{\text{\#}} P^{\Delta},\qquad
% Q:=(T_W)_{\text{\#}} Q^{\Delta}.
% \]
% Then $Q\ll P$ and
% \[
% \mathrm{KL}\left(Q\middle|P\right)
% =\mathrm{KL}\left(Q^{\Delta}\middle|P^{\Delta}\right).
% \]
% Equivalently, identifying ${W}\times\mathcal W_\Delta$ with $W+\mathcal W_\Delta$ via $\Phi$, one may write
% $P=\Phi_{\text{\#}}(\delta_W\otimes P^{\Delta})$ and
% $Q=\Phi_{\text{\#}}(\delta_W\otimes Q^{\Delta})$, and the same equality holds.
% \end{lemma}

\begin{lemma}[KL in PECL is adapter-subspace KL]
% Assume $Q_t^\Delta \ll P_t^\Delta$. 
% Assume $Q_t^\Delta \ll P_t^\Delta$, meaning $Q_t^\Delta(A)=0$ for any measurable set $A$ such that $P_t^\Delta(A)=0$.
Assume $Q_t^\Delta \ll P_t^\Delta$ (i.e., $Q_t^\Delta$ is absolutely continuous w.r.t.\ $P_t^\Delta$), meaning $Q_t^\Delta(A)=0$ for any measurable set $A$ such that $P_t^\Delta(A)=0$.
Then $Q_t \ll P_t$ and
$$
\mathrm{KL}\bigl(Q_t\,\Vert\, P_t\bigr)
\;=\;
\mathrm{KL}\bigl(Q_t^\Delta \,\Vert\, P_t^\Delta\bigr).
$$
\end{lemma}

\begin{proof}
  % Absolute continuity is preserved by pushforward under measurable maps, hence $Q_t\ll P_t$. Since $T_t$ is a Borel isomorphism, 
  Since $T_t$ is measurable, absolute continuity is preserved under pushforward, hence
$Q_t=T_{t\#}Q_t^\Delta \ll T_{t\#}P_t^\Delta=P_t$.
Moreover, $T_t$ is a measurable bijection with measurable inverse. 
  By the Radon-Nikodym chain rule
 \[
 \frac{dQ_t}{dP_t}(T_t(\Delta))=\frac{dQ_t^{\Delta}}{dP_t^{\Delta}}(\Delta)\quad\text{for }P^{\Delta}_t\text{-a.e.\ }\Delta.
 \]

\begin{equation*}
     \begin{aligned}
 \mathrm{KL}\left(Q_t\middle|P_t\right)
 &= \int_{W+\mathcal W_{\Delta,t}}
 \log\Big(\frac{dQ_t}{dP_t}(y)\Big)
 Q_t(dy)\\
 &= \int_{\mathcal W_{\Delta,t}}
 \log\Big(\frac{dQ_t}{dP_t}\big(T_t\Delta\big)\Big)
 Q_t^{\Delta}(d\Delta) \quad \text{(} Q_t=(T_t)_{\text{\#}}Q_t^{\Delta}\text{)} 
 \\
 &= \int_{\mathcal W_{\Delta,t}}
 \log\Big(\frac{dQ_t^{\Delta}}{dP_t^{\Delta}}(\Delta)\Big)
 Q_t^{\Delta}(d\Delta) \\
 &= \mathrm{KL}\left(Q_t^{\Delta}\middle|P_t^{\Delta}\right).
 \end{aligned}
\end{equation*} 
\end{proof}

% \section{ Proof of Lemma~\ref{lemma: subspace-full}}
% \begin{lemma}[Sharpness decomposition in PECL]
% \label{lemma: Sharpness in PECL}
% Let $W_0\in\mathbb R^d$ be fixed and let the trainable update lie in the adapter subspace $\mathcal W_\Delta\subseteq\mathbb R^d$. Consider a Gaussian posterior on $(\mathcal W_\Delta,\mathcal B(\mathcal W_\Delta))$,
%  $
%  Q_t^{\Delta}=\mathcal N(\hat{\Delta W}_t,\Sigma_t),
%  $
%  with mean $\hat{\Delta W}_t\in\mathcal W_\Delta$ and covariance operator $\Sigma_t\succeq 0$ acting on $\mathcal W_\Delta$ (viewed in $\mathbb R^d$, this covariance is supported in $\mathcal W_\Delta$). Define the full-space posterior on ${W_0}\times\mathcal W_\Delta$ by
%  $
%  Q_t:=\delta_{W_0}\otimes Q_t^{\Delta},
%  $
% and identify ${W_0}\times\mathcal W_\Delta$ with the affine subspace $W_0+\mathcal W_\Delta$ via $(w,\Delta)\mapsto w+\Delta$. Let $\hat L_{S_t}:\mathbb R^d\to\mathbb R$ be measurable and integrable under $Q_t$. Then
%  \[
%  \mathbb{E}_{W\sim Q_t}\big[\hat{L}_{S_t}(W)\big]
%  = \mathbb{E}_{\varepsilon\sim\mathcal N(0,\Sigma_t)}
%  \big[\hat{L}_{S_t}(W_0+\hat{\Delta W}_t+\varepsilon)\big]
%  = \hat{L}_{S_t}(W_0+\hat{\Delta W}_t) +\mathrm{RS}_t^{\Delta}(\hat{\Delta W}_t,\Sigma_t) \]
%    where the subspace expected sharpness is
%    \[
%    \mathrm{RS}_t^{\Delta}(\hat{\Delta W}_t,\Sigma_t)
%    = \mathbb{E}_{\varepsilon\sim\mathcal N(0,\Sigma_t)}
%    \big[\hat{L}_{S_t}(W_0+\hat{\Delta W}_t+\varepsilon)
%   - \hat{L}_{S_t}(W_0+\hat{\Delta W}_t)\big].
%      \]
% \end{lemma}
% \begin{proof}
%     Sampling $W\sim Q_t$ is equivalent to sampling $\varepsilon\sim\mathcal N(0,\Sigma_t)$ on $\mathcal W_\Delta$ and setting $W=W_0+\hat{\Delta W}_t+\varepsilon$. Taking expectations and adding-subtracting $\hat L_{S_t}(W_0+\hat{\Delta W}_t)$ inside the expectation gives the decomposition. 
% \end{proof}

\subsection{Proof of Lemma~\ref{lem:subspace-loss-pecl}}
\label{proof:lemma2}

% \begin{equation*}
% \label{eq:subspace-loss-pecl}
% \begin{aligned}
% &\mathbb E_{W\sim Q_t}\bigl[\hat L_{S_t}(W)\bigr] \\
% &= \mathbb{E}_{\varepsilon\sim\mathcal N(0,\Sigma_t)}
%     \big[\hat{L}_{S_t}(W_\mathrm {frozen }+\hat{\Delta W}_t+\varepsilon)\big],\varepsilon\in\mathcal W_\Delta \\
% & =\hat L_{S_t}(W_{\mathrm{frozen}} + \hat{\Delta W}_t)
% +
% \mathrm{RS}_t^\Delta(\hat{\Delta W}_t,\Sigma_t),\varepsilon\in\mathcal W_\Delta.
% \end{aligned}
% \end{equation*}
We specialize the posterior to a Gaussian supported on $\mathcal W_\Delta$. For task $t$, we consider a Gaussian posterior \( Q_t^{\Delta} = \mathcal{N}(\hat{\Delta W}_t, \Sigma_t) \) on the measurable space \( (\mathcal{W}_\Delta, \mathcal{B}(\mathcal{W}_\Delta)) \), where \( \hat{\Delta W}_t \in \mathcal{W}_\Delta \) is the mean and \( \Sigma_t\) is a covariance on the task-permissible subspace \( \mathcal{W}_\Delta \). Sampling from \( \Delta W \sim Q_t^{\Delta} \) is equivalent to 
$
\Delta W = \hat{\Delta W}_t + \varepsilon, \quad \varepsilon \sim \mathcal{N}(0, \Sigma_t) \ \text{on } \mathcal{W}_\Delta.
$ 

\begin{lemma}
    The posterior expected empirical loss decomposes as
\begin{equation*}
    \mathbb{E}_{W\sim Q_t}\big[\hat{L}_{S_t}(W)\big]
= \mathbb{E}_{\varepsilon\sim\mathcal N(0,\Sigma_t)}
    \big[\hat{L}_{S_t}(W_\mathrm {frozen }+\hat{\Delta W}_t+\varepsilon)\big]
= \hat{L}_{S_t}(W_\mathrm {frozen }+\hat{\Delta W}_t)
  + \mathrm{RS}_t^{\Delta}(\hat{\Delta W}_t,\Sigma_t),
\end{equation*}
where the random sharpness is
$$
\mathrm{RS}_t^{\Delta}(\hat{\Delta W}_t,\Sigma_t)
:= \mathbb{E}_{\varepsilon\sim\mathcal N(0,\Sigma_t)}
   \big[\hat{L}_{S_t}(W_\mathrm {frozen }+\hat{\Delta W}_t+\varepsilon)
       - \hat{L}_{S_t}(W_\mathrm {frozen }+\hat{\Delta W}_t)\big],\quad
\varepsilon\sim\mathcal N(0,\Sigma_t)\ \text{on } \mathcal W_\Delta$$
\end{lemma}


\begin{proof}
Recall that the adapter posterior is a Gaussian supported on the update subspace,
$
Q_t^\Delta = \mathcal N(\hat{\Delta W}_t,\Sigma_t)
$
on $(\mathcal W_\Delta,\mathcal B(\mathcal W_\Delta))$.
Hence we can represent a draw $\Delta W \sim Q_t^\Delta$ as
$
\Delta W = \hat{\Delta W}_t + \varepsilon
$
with $\varepsilon \sim \mathcal N(0,\Sigma_t)$ on $\mathcal W_\Delta$.

Under the PECL constraint, the full parameter is obtained by translating the adapter update by the frozen part,
$
W = W^{\mathrm{froz}} + \Delta W.
$
Therefore, for the empirical risk,
\begin{align*}
\mathbb E_{W \sim Q_t}\big[\hat L_{S_t}(W)\big]
&= \mathbb E_{\Delta W \sim Q_t^\Delta}\big[\hat L_{S_t}(W^{\mathrm{froz}}+\Delta W)\big] \\
&= \mathbb E_{\varepsilon \sim \mathcal N(0,\Sigma_t)}
\big[\hat L_{S_t}(W^{\mathrm{froz}}+\hat{\Delta W}_t+\varepsilon)\big].
\end{align*}
Finally, add and subtract the mean point loss to obtain
\begin{align*}
\mathbb E_{\varepsilon \sim \mathcal N(0,\Sigma_t)}
\big[\hat L_{S_t}(W^{\mathrm{froz}}+\hat{\Delta W}_t+\varepsilon)\big]
&=
\hat L_{S_t}(W^{\mathrm{froz}}+\hat{\Delta W}_t) \\
&\quad +
\mathbb E_{\varepsilon \sim \mathcal N(0,\Sigma_t)}
\Big[
\hat L_{S_t}(W^{\mathrm{froz}}+\hat{\Delta W}_t+\varepsilon)
-
\hat L_{S_t}(W^{\mathrm{froz}}+\hat{\Delta W}_t)
\Big].
\end{align*}
The second term is exactly the random Sharpness on the subspace
\[
\mathrm{RS}_t^\Delta(\hat{\Delta W}_t,\Sigma_t)
:=
\mathbb E_{\varepsilon \sim \mathcal N(0,\Sigma_t)}
\Big[
\hat L_{S_t}(W^{\mathrm{froz}}+\hat{\Delta W}_t+\varepsilon)
-
\hat L_{S_t}(W^{\mathrm{froz}}+\hat{\Delta W}_t)
\Big],
\]
completes the proof.
\end{proof}



\subsection{Proof of Theorem \ref{thm:pecl-bound}}
\label{proof:theorem}

% \subsection{Proof of }


\begin{theorem}[Sharpness-aware PAC--Bayes bound for PECL]
% \label{thm:pecl-bound}

For any $\delta\in(0,1]$, with probability at least $1-\delta$ over $S_1,\ldots,S_T$,
\begin{equation}\label{eq:pecl-main}
\begin{aligned}
&\frac{1}{T}\sum_{t=1}^T
\mathbb{E}_{P_t\sim\mathcal P_{t-1}}\mathbb{E}_{W\sim Q_t}L_{\mathcal{D}_t}(W)
\ \le\
\ \frac{1}{T}\sum_{t=1}^T
\mathbb{E}_{P_t\sim\mathcal P_{t-1}}
\mathbb{E}_{\varepsilon \sim \mathcal{N}(0, \Sigma_t)}[\hat{L}_{S_t}(W_0 + \widehat{\Delta W}_t + \varepsilon)
\\
&\quad +\sqrt{\frac{
\displaystyle \sum_{t=1}^T \mathbb{E}_{P_t^\Delta\sim\mathcal P_{t-1}}\!\Big[\mathrm{KL}\!\big(Q_t^\Delta\|P_t^\Delta\big)\Big]
\;+\;\displaystyle \sum_{t=1}^T \mathrm{KL}\!\big(\mathcal P_t\|\mathcal P_{t-1}\big)
\;+\;\log\!\frac{1}{\delta}
}{2\,\displaystyle\sum_{t=1}^T (m_t-1)}}.
\end{aligned}
\end{equation}
\end{theorem}



\begin{proof}
The bound follows by specializing Theorem \ref{the: full} to the PECL setting through two key decompositions:
By Lemma~\ref{lem:kl-decomp-pecl}, the full-space KL terms reduce to their adapter-space counterparts:
\[
\mathrm{KL}(Q_t \| P_t) = \mathrm{KL}(Q_t^\Delta \| P_t^\Delta).
\]
% \quad
% \mathrm{KL}(\mathcal{P}_t \| \mathcal{P}_{t-1}) = \mathrm{KL}(\mathcal{P}_t^\Delta \| \mathcal{P}_{t-1}^\Delta)
% \textbf{Step 2: Sharpness decomposition via Lemma 6.}  
Under the Gaussian posterior $Q_t^\Delta = \mathcal{N}(\Delta\hat{W}_t, \Sigma_t)$, Lemma \ref{lem:subspace-loss-pecl} gives:
\[
\mathbb{E}_{W \sim Q_t}[\hat{L}_{S_t}(W)] 
= \mathrm{E}_{\varepsilon \sim \mathcal{N}(0, \Sigma_t)}{L}_{S_t}(W^{\mathrm{froz}}_t+\widehat{\Delta W}_t + \varepsilon)
 = \hat{L}_{S_t}(W^{\mathrm{froz}}_t + \widehat{\Delta{W}_t}) + \mathrm{RS}_t^\Delta(\Delta\widehat{W}_t, \Sigma_t),
\]
where the perturbations are confined to the task-permissible subspace.

% \textbf{Step 3: Substitution into general bound.}
Substituting these decompositions into Lemma \ref{the: full} yields the final PECL-specific bound, where all complexity and sharpness terms are now restricted to the task-permissible subspace $\mathcal{W}_\Delta$.
% Substituting these decompositions into Theorem 2 yields the final PECL-specific bound, where all complexity and sharpness terms are now restricted to the trainable LoRA subspace $\mathcal{W}_\Delta$.

\end{proof}



\section{Bridging the Permissible Update Subspace and the LoRA Implementation}
\label{app:lora_subspace_bridge}

This appendix clarifies how the abstract permissible update subspace
$\mathcal W_{\Delta,t}$ used in the PAC--Bayesian analysis relates to the concrete LoRA
parameterization implemented in practice.
The key point is that, throughout the paper, $\mathcal W_{\Delta,t}$ denotes a
\emph{linear subspace of the ambient weight-increment space} $\mathbb R^d$ with
$d=\dim(\mathrm{vec}(\Delta W))$, rather than the rank-$r$ factor coordinate space of $(A,B)$.
This distinction is essential because the LoRA mapping $\theta_\Delta=(A,B)\mapsto \Delta W=BA$
is bilinear and therefore does not define a global linear set of reachable increments in weight space.
Our PAC--Bayesian analysis is local and only requires a linear model of permissible \emph{increment directions}
around the learned solution, which is precisely what the Jacobian provides.

% \paragraph{Abstract object in the bound: a linear increment subspace.}
% In the main text, PECL is modeled as restricting task-$t$ updates to a linear subspace
% $\mathcal W_{\Delta,t}\subseteq\mathbb R^d$, leading to the affine model class
% $
% W_{\mathrm{froz}}+\mathcal W_{\Delta,t}.
% $
% All priors, posteriors, and perturbation distributions appearing in
% Theorem~\ref{thm:pathwise_pac} are defined on $(\mathcal W_{\Delta,t},\mathcal B(\mathcal W_{\Delta,t}))$
% and pushed forward to the affine class by translation.
% Concretely, letting $T_{W_{\mathrm{froz}}}(\Delta)=W_{\mathrm{froz}}+\Delta$, we work with
% \[
% P_t := (T_{W_{\mathrm{froz}}})_{\#} P_t^{\Delta},
% \qquad
% Q_t := (T_{W_{\mathrm{froz}}})_{\#} Q_t^{\Delta},
% \]
% where $P_t^\Delta$ and $Q_t^\Delta$ are measures supported on $\mathcal W_{\Delta,t}$.
% Lemma~\ref{lem:kl-decomp-pecl} shows that
% \[
% \mathrm{KL}\!\left(Q_t\Vert P_t\right)
% =
% \mathrm{KL}\!\left(Q_t^{\Delta}\Vert P_t^{\Delta}\right),
% \]
% so under PECL the KL contribution is determined entirely by the distributions supported on
% $\mathcal W_{\Delta,t}$.

% \paragraph{Implemented object: LoRA paramaters.}
For a given adapted weight block (or layer) $\ell$, LoRA introduces trainable coordinates
$\theta_{\Delta,\ell}=(A_\ell,B_\ell)$ and defines the effective increment
\[
\Delta W_\ell(A_\ell,B_\ell)= B_\ell A_\ell,
\qquad
A_\ell\in\mathbb R^{r\times n_\ell},\;
B_\ell\in\mathbb R^{m_\ell\times r}.
\]
Let $\hat\theta_{\Delta,t}=(\hat A_t,\hat B_t)$ denote the learned LoRA parameters after training task $t$.
Consider a small coordinate perturbation
$\delta\theta_\Delta=(\delta A,\delta B)$ around $\hat\theta_{\Delta,t}$, and define
\[
g(\theta_\Delta):=\mathrm{vec}(\Delta W(\theta_\Delta))\in\mathbb R^d,
\qquad d=m_\ell n_\ell .
\]
A first-order expansion yields
\begin{equation}
\label{eq:lora_local_linearization}
g(\hat\theta_{\Delta,t}+\delta\theta_\Delta)
=
g(\hat\theta_{\Delta,t})
+
J_{\mathrm{LoRA}}(\hat\theta_{\Delta,t})\,\delta\theta_\Delta
+
R(\delta\theta_\Delta),
\end{equation}
where $J_{\mathrm{LoRA}}(\hat\theta_{\Delta,t})$ is the Jacobian of $g$ evaluated at $\hat\theta_{\Delta,t}$,
and $\|R(\delta\theta_\Delta)\|_2=O(\|\delta\theta_\Delta\|_2^2)$ as $\|\delta\theta_\Delta\|_2\to 0$.

Crucially, for LoRA the Jacobian image admits an explicit characterization.
Writing $\delta\theta_\Delta=(\delta A,\delta B)$, the induced first-order increment variation in weight space is
\[
\delta(\Delta W_\ell)
=
s_\ell(\delta B\,\hat A_t+\hat B_t\,\delta A),
\]
and hence the Jacobian image at the learned coordinates is the linear set
\begin{equation}
\label{eq:WDelta_Jacobian_image_explicit}
\mathrm{Im}\!\big(J_{\mathrm{LoRA}}(\hat\theta_{\Delta,t})\big)
=
\Big\{
\mathrm{vec}\!\big(s_\ell(\delta B\,\hat A_t+\hat B_t\,\delta A)\big)
:\;
\delta A\in\mathbb R^{r\times n_\ell},\;
\delta B\in\mathbb R^{m_\ell\times r}
\Big\}
\subseteq \mathbb R^d .
\end{equation}
This is a \emph{local subspace in increment space}, not a global reachable set, because it is defined
by the first-order linearization around $\hat\theta_{\Delta,t}$.

% \paragraph{Identification used by the theory: $\mathcal W_{\Delta,t}$ as a local tangent model.}
The abstract subspace $\mathcal W_{\Delta,t}$ in Theorem~\ref{thm:pathwise_pac} is interpreted as a local linear
model of permissible increment \emph{directions} around the task-$t$ solution in the ambient increment space.
Accordingly, we instantiate it via the LoRA-induced tangent space
\begin{equation}
\label{eq:WDelta_t_identification}
\mathcal W_{\Delta,t}
\equiv
\mathrm{Im}\!\big(J_{\mathrm{LoRA}}(\hat\theta_{\Delta,t})\big),
\end{equation}
and define the PAC--Bayesian distributions directly on this linear subspace, as required by
Lemma~\ref{lem:kl-decomp-pecl} and Theorem~\ref{thm:pathwise_pac}.

% \paragraph{How coordinate perturbations induce increment perturbations.}
In implementation, perturbations are applied in the LoRA coordinate space.
Under the local linearization~\eqref{eq:lora_local_linearization}, a coordinate perturbation
$\delta\theta_\Delta=(\delta A,\delta B)$ induces an increment perturbation
\[
\delta\Delta
:=
g(\hat\theta_{\Delta,t}+\delta\theta_\Delta)-g(\hat\theta_{\Delta,t})
\approx
J_{\mathrm{LoRA}}(\hat\theta_{\Delta,t})\,\delta\theta_\Delta,
\]
whose support lies in $\mathrm{Im}(J_{\mathrm{LoRA}}(\hat\theta_{\Delta,t}))=\mathcal W_{\Delta,t}$.
For instance, if $\delta\theta_\Delta\sim\mathcal N(0,\Sigma_{\theta,t})$ and the perturbation radius is chosen
so that the remainder term is negligible with high probability, then the induced increment perturbation is approximately Gaussian in $\mathbb R^d$,
\begin{equation}
\label{eq:Sigma_pushforward}
\delta\Delta \approx \mathcal N(0,\Sigma_t),
\qquad
\Sigma_t := J_{\mathrm{LoRA}}(\hat\theta_{\Delta,t})\,\Sigma_{\theta,t}\,
J_{\mathrm{LoRA}}(\hat\theta_{\Delta,t})^\top,
\end{equation}
and is supported on $\mathcal W_{\Delta,t}$.

% \paragraph{Takeaway.}
The bound is formulated on $\mathcal W_{\Delta,t}\subset\mathbb R^d$ because it is a statement about perturbations
and posterior complexity in \emph{weight-increment space} under PECL constraints.
LoRA is a concrete mechanism that realizes such a constraint.
Although the mapping $(A,B)\mapsto BA$ is bilinear and does not define a global linear class of increments,
its local linearization around the learned coordinates induces the tangent increment subspace
$\mathrm{Im}(J_{\mathrm{LoRA}}(\hat\theta_{\Delta,t}))$, which provides an implementation-level realization of
$\mathcal W_{\Delta,t}$ while keeping the KL analysis in Lemma~\ref{lem:kl-decomp-pecl} intact.


\section{Feature Amplification and Curvature Projection Ratio}
\label{app:weight_projection}

This appendix details two diagnostic measures used to interpret why the perturbation scope matters in LoRA-based continual learning: the feature amplification factor $\alpha_t$ and the curvature projection preference $c_t$. 
These two diagnostics are computed on a single adapted block (layer/module) in experiments rather than the full parameter vector.



\subsection{Feature Amplification Factor}
\label{app:feature_amplification}

Let $\Delta W_t$ denote the learned weight update for a given adapted block at task $t$ (implemented as $\Delta W_t = B_t A_t$).
Let
$
\Delta W_t = U_{\Delta,t} S_{\Delta,t} V_{\Delta,t}^\top
$
be its rank-$r$ truncated SVD. The pair $(U_{\Delta,t},V_{\Delta,t})$ induces the bilinear subspace
$
\mathcal S_{\Delta,t}=\{U_{\Delta,t} Z V_{\Delta,t}^\top : Z\in\mathbb{R}^{r\times r}\}.
$
Note that $\mathcal{S}_{\Delta, t}$ is an empirical diagnostic subspace defined from the learned update $\Delta W_t$ (via its truncated SVD) and is used only for curvature localization. $\mathcal{W}_{\Delta, t}=\operatorname{Im}\left(J_{\text {LoRA }}\right)$ is the local tangent space of LoRA-induced increments in $\operatorname{vec}(\Delta W)$-space, whereas $\mathcal{S}_{\Delta, t}$ captures only the core component that preserves the row and column subspaces of $\Delta W_t$; hence $\mathcal{S}_{\Delta, t}$ is an empirical diagnostic subspace of the current solution and is not assumed to coincide with $\mathcal{W}_{\Delta, t}$.
Under the Frobenius inner product, the orthogonal projection of a matrix $M$ onto $\mathcal S_{\Delta,t}$ is
% $
% \Pi_{\Delta,t}(M)=U_{\Delta,t}(U_{\Delta,t}^\top M V_{\Delta,t})V_{\Delta,t}^\top.
% $
$
\Pi_{\mathcal S,t}(M)=U_{\Delta,t}(U_{\Delta,t}^\top M V_{\Delta,t})V_{\Delta,t}^\top.
$
Equivalently, since $U_{\Delta,t}$ and $V_{\Delta,t}$ have orthonormal columns,
% $
% \|\Pi_{\Delta,t}(M)\|_F = \|U_{\Delta,t}^\top M V_{\Delta,t}\|_F.
% $
$\|\Pi_{\mathcal S,t}(M)\|_F = \|U_{\Delta,t}^\top M V_{\Delta,t}\|_F.
$



Following~\cite{hu_lora_2021}, we define the feature amplification factor
\[
\alpha_t=\frac{\|\Delta W_t\|_F}{\|U_{\Delta,t}^\top W V_{\Delta,t}\|_F}
=\frac{\|\Delta W_t\|_F}{\|\Pi_{\mathcal S,t}(W)\|_F},
\]
%=\frac{\|\Delta W_t\|_F}{\|\Pi_{\Delta,t}(W)\|_F},
which compares the magnitude of the learned update to the backbone energy restricted to the same bilinear subspace induced by $\Delta W_t$.
Larger $\alpha_t$ indicates that, relative to the backbone component available within $\mathcal S_{\Delta,t}$, the learned update contributes a larger effective change along the same row and column subspaces.

\subsection{Curvature Projection Ratio}
\label{app:curvature_preference}

We next quantify whether dominant task-conditioned curvature modes concentrate more on $\mathcal S_{\Delta,t}$ than on a backbone reference bilinear subspace.

% \paragraph{Curvature eigenpairs in the LoRA-augmented parameter space.}
Let $\{(\lambda_{t,i}, v_{t,i})\}_{i=1}^k$ denote the top-$k$ eigenpairs of a positive semidefinite curvature operator at task $t$ (e.g., GGN/Fisher) computed in the parameter space that includes the backbone block $W$ and the LoRA factors $(A_t,B_t)$ (with $W$ frozen during training).
% \paragraph{Mapping eigenvectors to effective weight-space directions.}
Each eigenvector $v_{t,i}$ is a direction in the LoRA-augmented parameter space.
We decompose it into the corresponding blocks and denote the induced local perturbations by
$(\delta W_{t,i}^{(0)}, \delta A_{t,i}, \delta B_{t,i})$,
where $\delta W_{t,i}^{(0)}$ is reshaped to match the backbone weight block.
Under the first-order pushforward of the LoRA parametrization, this direction induces an effective weight-space perturbation
\[
\delta W_{t,i}
=
\delta W_{t,i}^{(0)} + \big(\delta B_{t,i}A_t + B_t \delta A_{t,i}\big).
\]
% Each eigenvector $v_{t,i}$ is a direction in the full parameter space.
% We map it to an effective direction in the weight-matrix space by aggregating its components on the backbone block and the LoRA factors:
% \[
% dW_{t,i} = dW_{t,i}^{(0)} + s\big(dB_{t,i}A_t + B_t dA_{t,i}\big),
% \]
where $ \delta W_{t,i}^{(0)}$ is the component corresponding to the backbone weight block, $(\delta A_{t,i},\delta B_{t,i})$ are the components corresponding to the LoRA factors.
This mapping matches the implementation, where curvature directions are evaluated after being pushed forward to effective weight-space increments.

% \paragraph{Per-mode projection ratios.}
We define the per-mode projection ratio onto $\mathcal S_{\Delta,t}$ as
% $$
% \rho_{\Delta,t,i}
% =\frac{\|\Pi_{\Delta,t}(\delta W_{t,i})\|_F^2}{\|\delta W_{t,i}\|_F^2}.
% $$
$$
\rho_{\Delta,t,i}
=\frac{\|\Pi_{\mathcal S,t}(\delta W_{t,i})\|_F^2}{\|\delta W_{t,i}\|_F^2}.
$$
For the backbone reference, let $(U_{W,t},V_{W,t})$ be the rank-$r$ truncated SVD subspace of the corresponding backbone weight block (restricted to the same layer/block), inducing
$
\mathcal S_{W,t}=\{U_{W,t} Z V_{W,t}^\top : Z\in\mathbb{R}^{r\times r}\},
\quad
\Pi_{\mathcal S_W,t}(M)=U_{W,t}(U_{W,t}^\top M V_{W,t})V_{W,t}^\top,
$
%\Pi_{W,t}(M)=U_{W,t}(U_{W,t}^\top M V_{W,t})V_{W,t}^\top,
and define
% $$
% \rho_{W,t,i}
% =\frac{\|\Pi_{W,t}(\delta W_{t,i})\|_F^2}{\|\delta W_{t,i}\|_F^2}.
% $$
$$\rho_{W,t,i}
=\frac{\|\Pi_{\mathcal S_W,t}(\delta W_{t,i})\|_F^2}{\|\delta W_{t,i}\|_F^2}.
$$

% \paragraph{Eigenvalue-weighted aggregation and preference ratio.}
We aggregate the localization of the top-$k$ curvature modes by eigenvalue-weighted averages
\[
R_{\Delta,t}=\frac{\sum_{i=1}^k \lambda_{t,i}\rho_{\Delta,t,i}}{\sum_{i=1}^k \lambda_{t,i}},
\qquad
R_{W,t}=\frac{\sum_{i=1}^k \lambda_{t,i}\rho_{W,t,i}}{\sum_{i=1}^k \lambda_{t,i}},
\]
and report the curvature projection ratio
\[
c_t=\frac{R_{\Delta,t}}{R_{W,t}}.
\]
Larger $c_t$ indicates that dominant curvature modes have higher projected Frobenius energy within the bilinear subspace induced by $\Delta W_t$ than within the backbone reference bilinear subspace.






\section{Implementation Details}

% Our code is available here: \url{https://anonymous.4open.science/r/FlatnessICMLNoname9527223/}.
% Unless stated otherwise, we use SGD as the base optimizer, train each session for $20$ epochs and SAM radius $\rho = 0.05$. LoRA based variants are trained with learning rate $0.01$ without momentum or weight decay. The same training configuration is used across SeqLoRA, IncLoRA, and OLoRA, so that differences in performance can be attributed to the geometry of the adapter subspace and the choice of flatness optimizer rather than to changes in training budget.

% \subsection{Implementation Details}
\noindent\textbf{Data.}
We evaluate continual learning on ImageNet-R. For robustness under corruptions and perturbations, we use the Tiny ImageNet variants of ImageNet-C and ImageNet-P (200 classes) and construct ImageFolder-style splits that prevent leakage by grouping all corruption/perturbation variants from the same original image (or video) into a single train/test partition. All robustness evaluations use the same $224$-resolution preprocessing pipeline as the main experiments.

\noindent\textbf{Experimental setup.}
Unless stated otherwise, we use SGD as the base optimizer and train each task/session for $20$ epochs with batch size $128$ and a cosine learning-rate schedule. For sharpness-aware optimization, we set the SAM radius to $\rho=0.05$. LoRA-based variants are trained with learning rate $0.01$ and no momentum or weight decay. The same training configuration is used across SeqLoRA, IncLoRA, and OLoRA to isolate the effect of adapter geometry and perturbation design.

\noindent\textbf{LoRA configuration.}
Unless specified otherwise, the LoRA rank is set to $r=16$ for the main ImageNet-R experiments. For the Tiny ImageNet-C/P robustness experiments, we use $r=8$. In rank ablations, we sweep $r\in\{2,8,16\}$.

\noindent\textbf{Evaluation protocol.}
For accuracy reporting, we aggregate results over multiple random seeds. For sharpness and curvature measurements, we fix the random seed to $1993$ to ensure controlled comparisons and reduce estimator variance. Flatness/curvature evaluation is performed on a $10\%$ subset of the evaluation set with batch size $32$. Random sharpness uses $100$ Gaussian samples. Hessian trace and top-eigenvalue estimates are computed with $100$ Lanczos power iterations and $100$ trace samples. Loss-landscape visualizations use radius $1$ with $41$ points along a random basis with norm-filtered directions.

\noindent\textbf{Weight and curvature analysis.}
In the weight/curvature analysis, we include frozen parameters when computing full-parameter Hessian-vector products and use Lanczos with top-$k=16$ eigenvalues. Curvature localization and $\Delta W$ projections are computed on the last transformer block (block 11) and restricted to the $q$ and $v$ subspaces with row ranges $[0,768]$ and $[1536,2304]$, respectively. The localized Rayleigh estimator uses $128$ samples, and the curvature MVP backend is set to GGN. For $\Delta W$ and $\Delta W_{\mathrm{full}}$ projections, we use the QKV weight and its corresponding LoRA factors with rank $16$. The $W$--$\Delta W$ alignment analysis uses rank $16$, seed $42$, and the seqlora setting.




\subsection{How Fig.~\ref{fig:scope_type2} is Constructed: Curvature-Aligned 1D Loss Probes (Full vs.\ LoRA)}
\label{app:curv1d_probe}

\noindent\textbf{Task and data.}
All curves are produced in the class-incremental ImageNet-R setup.
For each task $t$, we take the reference model $\widehat W_t$ (the full model parameters after learning task $t$) and evaluate losses on a fixed subset of mini-batches sampled from the evaluation source used for flatness analysis (we use a fixed sampling ratio and a fixed random seed across methods to ensure comparability).

\noindent\textbf{Curvature operator and leading eigen-direction.}
Let $\mathcal L(\widehat W_t)$ denote the empirical loss averaged over the fixed batches.
We define a curvature operator $\mathbf C_t$ at $\widehat W_t$ using a curvature proxy (GGN in our implementation), accessed only through a matrix--vector product (MVP) routine. In our experiments, we use the GGN as the curvature proxy.
We estimate the leading eigenpair $(\lambda_t, v_t)$ by running Lanczos~\cite{ghorbaniInvestigationNeuralNet2019} algorithm on the MVP operator, i.e.,
$$
\mathbf C_t v_t \approx \lambda_t v_t,
\qquad \|v_t\|=1.
$$
This answers the practical question of how the eigenvectors are computed: we never form $\mathbf C_t$ explicitly, but recover its dominant eigen-direction from repeated MVP calls.

% \noindent\textbf{Full-parameter probe vs.\ LoRA-restricted probe.}
% To clarify what it means to ``probe curvature restricted to trainable parameters,'' we explicitly generate two curvature-aligned 1D loss curves by running eigen-direction extraction and loss sweeps on two different parameter sets.

% \begin{itemize}
% \item \textbf{Full-parameter curvature probe (bottom curve).}
% We define a curvature matrix--vector product (MVP) over the full parameter set \texttt{params\_full} (provided as \texttt{mvp\_full}).
% We then estimate the leading eigenpair $(\lambda_t^{\mathrm{full}}, v_t^{\mathrm{full}})$ by applying Lanczos/power iteration to \texttt{mvp\_full}.
% The resulting eigenvector is unflattened to parameter-shaped tensors and used to evaluate a 1D loss curve by perturbing the full parameter set along $v_t^{\mathrm{full}}$ while keeping all other components unchanged.

% \item \textbf{LoRA-restricted curvature probe (top curve).}
% Let $\mathcal I_\Delta$ denote the index set of trainable LoRA parameters (implemented as \texttt{params\_lora}, e.g., the $A/B$ factors).
% We define an MVP \emph{directly} on this LoRA parameter subspace via \texttt{mvp\_lora}.
% Running Lanczos/power iteration on \texttt{mvp\_lora} yields a leading eigenpair $(\lambda_t^\Delta, v_t^\Delta)$ supported only on LoRA coordinates.
% We then compute the 1D loss curve by perturbing \emph{only} \texttt{params\_lora} along $v_t^\Delta$, while keeping the frozen backbone parameters fixed throughout.
% % (Optionally: This restriction can be interpreted as probing the full-space operator restricted to the LoRA subspace, although the implementation uses a dedicated LoRA-subspace MVP.)
% \end{itemize}

\noindent\textbf{Full-parameter probe vs.\ LoRA-restricted probe.}
To make the notion of ``curvature probes restricted to trainable parameters'' explicit, we construct two curvature-aligned 1D loss curves by extracting dominant curvature directions on two different parameter domains and sweeping the loss along the resulting directions.

\begin{itemize}
% \item \textbf{Full-parameter curvature probe (bottom curve).}
% Let $\Theta_{\mathrm{full}}$ denote the ambient parameter space of the model (including the backbone and the LoRA parameters), and let $\widehat W_t\in\Theta_{\mathrm{full}}$ be the reference model after learning task $t$.
% We choose a curvature operator $\mathbf C_t^{\mathrm{full}}$ at $\widehat W_t$ ( GGN ) defined on $\Theta_{\mathrm{full}}$, accessed through its matrix--vector product.
% We estimate its leading eigenpair $(\lambda_t^{\mathrm{full}}, v_t^{\mathrm{full}})$ via Lanczos iteration, where $v_t^{\mathrm{full}}\in\Theta_{\mathrm{full}}$ and $\|v_t^{\mathrm{full}}\|=1$.
% We then evaluate the 1D loss curve by perturbing the full parameter vector along this direction:
% \[
% W_t(\alpha)=\widehat W_t+\alpha\, v_t^{\mathrm{full}},\qquad \alpha\in[-r_{\mathrm{full}},r_{\mathrm{full}}],
% \]
% and computing the empirical loss on a fixed set of mini-batches.

% \item \textbf{LoRA-restricted curvature probe (top curve).}
% Let $\Theta_{\Delta,t}$ denote the subspace of trainable LoRA parameters at task $t$ (the adapter coordinates), and keep the backbone parameters fixed at $\widehat W_t$.
% We define a curvature operator $\mathbf C_t^{\Delta}$ \emph{directly on} $\Theta_{\Delta,t}$ (again accessed through matrix--vector products), and estimate its leading eigenpair $(\lambda_t^{\Delta}, v_t^{\Delta})$ via Lanczos, where $v_t^{\Delta}\in\Theta_{\Delta,t}$ and $\|v_t^{\Delta}\|=1$.
% We then evaluate the 1D loss curve by perturbing \emph{only} the LoRA parameters along $v_t^{\Delta}$ while leaving the frozen backbone unchanged:
% \[
% W_t(\alpha)=\widehat W_t+\alpha\, \tilde v_t^{\Delta},\qquad \alpha\in[-r_{\Delta},r_{\Delta}],
% \]
% where $\tilde v_t^{\Delta}$ denotes the embedding of $v_t^{\Delta}$ into the full parameter space by setting all non-LoRA coordinates to zero.
\item \textbf{Full-parameter curvature probe (bottom curve).}
Let $\mathcal H_{\mathrm{full}}$ denote the chosen \emph{parameter set} used in the full-space probe (a collection of tensors that includes the LoRA parameters and, optionally, frozen backbone tensors).
We identify this parameter set with a Euclidean product space
\[
\Theta_{\mathrm{full}}
:=
\prod_{h\in \mathcal P_{\mathrm{full}}}\mathbb R^{\mathrm{shape}(h)},
\qquad
d_{\mathrm{full}}:=\sum_{h\in \mathcal H_{\mathrm{full}}}\mathrm{numel}(h),
\]
and use a fixed vectorization map $\mathrm{vec}_{\mathrm{full}}:\Theta_{\mathrm{full}}\to\mathbb R^{d_{\mathrm{full}}}$ with inverse $\mathrm{unvec}_{\mathrm{full}}$ that reshapes a vector back to the tensor list.
We define a GGN-based curvature operator $\mathbf C_t^{\mathrm{full}}$ on $\mathbb R^{d_{\mathrm{full}}}$ via matrix--vector products, estimate its leading eigenpair $(\lambda_t^{\mathrm{full}}, v_t^{\mathrm{full}})$ with $\|v_t^{\mathrm{full}}\|_2=1$, and then map $v_t^{\mathrm{full}}$ back to tensor-shaped perturbations $\mathrm{unvec}_{\mathrm{full}}(v_t^{\mathrm{full}})$.
The loss curve is evaluated by perturbing \emph{only} tensors in $\mathcal H_{\mathrm{full}}$ along this direction:
\[
W_t(\alpha)=\widehat W_t+\alpha\, \tilde v_t^{\mathrm{full}},
\qquad \alpha\in[-r_{\mathrm{full}},r_{\mathrm{full}}],
\]
where $\tilde v_t^{\mathrm{full}}$ denotes the perturbation embedded into the full model by adding the corresponding tensor increments on $\mathcal H_{\mathrm{full}}$ and leaving all other tensors unchanged.

\item \textbf{LoRA-restricted curvature probe (top curve).}
Let $\mathcal H_{\Delta,t}$ denote the \emph{trainable LoRA parameter set} at task $t$ (e.g., the LoRA $A/B$ factors), and define
\[
\Theta_{\Delta,t}
:=
\prod_{h\in \mathcal H_{\Delta,t}}\mathbb R^{\mathrm{shape}(h)},
\qquad
d_{\Delta,t}:=\sum_{h\in \mathcal H_{\Delta,t}}\mathrm{numel}(h),
\]
with the corresponding vectorization maps $\mathrm{vec}_{\Delta,t}$ and $\mathrm{unvec}_{\Delta,t}$.
We define a GGN-based curvature operator $\mathbf C_t^\Delta$ \emph{directly on} $\mathbb R^{d_{\Delta,t}}$ (through matrix--vector products restricted to $\mathcal H_{\Delta,t}$), and estimate its leading eigenpair $(\lambda_t^\Delta,v_t^\Delta)$ with $\|v_t^\Delta\|_2=1$.
We then evaluate the loss curve by perturbing \emph{only} the LoRA parameters along $\mathrm{unvec}_{\Delta,t}(v_t^\Delta)$ while keeping the backbone fixed:
\[
W_t(\alpha)=\widehat W_t+\alpha\, \tilde v_t^{\Delta},
\qquad \alpha\in[-r_{\Delta},r_{\Delta}],
\]
where $\tilde v_t^{\Delta}$ embeds the LoRA-only perturbation into the full model by setting all non-LoRA parameter increments to zero.

\end{itemize}

\noindent
In both probes, we use the same reference model $\widehat W_t$ and evaluate losses on the same fixed mini-batches; the only difference is whether the curvature eigen-direction is extracted and applied in the ambient parameter space or in the LoRA trainable subspace.
(Equivalently, the LoRA probe can be viewed as probing the ambient curvature restricted to $\Theta_{\Delta,t}$, although we implement it by defining the curvature operator directly on the LoRA coordinates.)



% \noindent\textbf{Full-parameter probe vs.\ LoRA-restricted probe.}
% We explicitly separate two probes to clarify what ``restricted to trainable parameters'' means.

% \begin{itemize}
% \item \textbf{Full-parameter curvature probe (bottom curve).}
% We run Lanczos on the curvature MVP defined over the \emph{full} parameter vector (including frozen backbone parameters).
% This yields a dominant direction $v_t^{\mathrm{full}}$ in the ambient parameter space.
% We then sweep the loss by perturbing the full model:
% $$
% W_t(\alpha) = \widehat W_t + \alpha\, v_t^{\mathrm{full}},
% \qquad \alpha\in[-r,r].
% $$

% \item \textbf{LoRA-restricted curvature probe (top curve).}
% Let $\mathcal I_\Delta$ denote the index set of trainable LoRA parameters (e.g., the $A/B$ factors), and let $\Pi_\Delta$ be the projection that keeps only coordinates in $\mathcal I_\Delta$ (all other coordinates are set to zero).
% We define the \emph{restricted} curvature MVP by composing the full MVP with this projection:
% $$
% u \mapsto \Pi_\Delta\big(\mathbf C_t(\Pi_\Delta(u))\big),
% $$
% which corresponds to the curvature restricted to the LoRA update subspace.
% Running Lanczos on this restricted operator yields $v_t^\Delta$ supported only on LoRA parameters.
% We then sweep the loss by perturbing \emph{only} the LoRA parameters while keeping the backbone fixed:
% $$
% W_t(\alpha) = \widehat W_t + \alpha\, v_t^\Delta,
% \qquad \text{with } (v_t^\Delta)_{\backslash \Delta}=0.
% $$
% \end{itemize}

\noindent\textbf{Loss curve evaluation and visualization.}
For a fixed grid $\{\alpha_i\}_{i=1}^n$ (controlled by a radius $r$ and number of points $n$), we evaluate
$$
\mathcal{L}_t(\alpha_i)
=
\frac{1}{|\mathcal B|}\sum_{B\in\mathcal B}\ell\!\left(W_t(\alpha_i);B\right),
$$
where $\mathcal B$ is the fixed set of mini-batches.
The resulting 1D curves $\{\alpha_i,\mathcal L_t(\alpha_i)\}$ are saved together with the corresponding directions and leading eigenvalues (for both the full and LoRA-restricted probes).
Finally, we overlay curves from different optimizers by reading their corresponding files and plotting $\mathcal L_t(\alpha)$ for direct comparison.

%(or $\mathcal L_t(\alpha)-\mathcal L_t(0)$) 






\section{Additional Experimental Results}


\subsection{Raw Task-Wise $\operatorname{ACC}^{\mathrm{cls}}_t$ on ImageNet-R and ImageNet-C for Fig.~\ref{fig:robustness_imagenet_cp}}

To complement Fig.~\ref{fig:robustness_imagenet_cp}, we provide the underlying raw task-wise $\operatorname{ACC}^{\mathrm{cls}}_t$ trajectories for ImageNet-C and ImageNet-P in the tables below.
All entries are {means over repeated runs} (averaged across the same set of random seeds used for the main figure).



\begin{table*}[ht]
\centering
\caption{Raw task-wise $\operatorname{ACC}^{\mathrm{cls}}_t$ trajectories ($t=1,\dots,20$) on ImageNet-C under different optimizers (mean over repeated runs).}
\label{tab:taskwise_acc_imagenetr}
\setlength{\tabcolsep}{2.2pt}
\renewcommand{\arraystretch}{0.92}
\scriptsize
\begin{tabular}{@{}l l *{20}{r}@{}}
\toprule
\textbf{Method} & \textbf{Opt.} &
\textbf{t01} & \textbf{t02} & \textbf{t03} & \textbf{t04} & \textbf{t05} &
\textbf{t06} & \textbf{t07} & \textbf{t08} & \textbf{t09} & \textbf{t10} &
\textbf{t11} & \textbf{t12} & \textbf{t13} & \textbf{t14} & \textbf{t15} &
\textbf{t16} & \textbf{t17} & \textbf{t18} & \textbf{t19} & \textbf{t20} \\
\midrule

\multirow{3}{*}{\textbf{SeqLoRA}}
& SGD
& 92.67 & 87.84 & 87.67 & 83.42 & 82.13 & 81.50 & 79.62 & 78.83 & 78.11 & 76.47
& 75.85 & 75.28 & 74.69 & 74.24 & 74.27 & 73.48 & 72.69 & 71.72 & 70.49 & 69.77 \\
& \quad $+\mathrm{AS}^{(0)}_S$
& 92.67 & 87.67 & 87.33 & 84.42 & 83.20 & 82.39 & 80.33 & 79.29 & 79.11 & 78.20
& 77.73 & 77.00 & 76.69 & 77.00 & 77.33 & 77.08 & 76.45 & 75.94 & 75.33 & 73.93 \\
& \quad $+\mathrm{AS}^{(1)}_S$
& 95.67 & 93.67 & 89.78 & 88.84 & 87.67 & 86.61 & 86.67 & 86.33 & 84.85 & 84.14
& 82.61 & 82.17 & 81.87 & 81.60 & 81.42 & 80.44 & 80.10 & 79.96 & 80.21 & 80.03 \\

\addlinespace[2pt]

\multirow{3}{*}{\textbf{IncLoRA}}
& SGD
& 92.33 & 89.84 & 88.22 & 86.58 & 83.40 & 81.22 & 80.38 & 78.46 & 75.89 & 73.77
& 73.21 & 73.61 & 73.02 & 72.95 & 71.62 & 70.17 & 68.82 & 66.30 & 64.86 & 61.77 \\
& \quad $+\mathrm{AS}^{(0)}_S$
& 92.33 & 90.83 & 90.45 & 85.92 & 87.00 & 83.39 & 81.81 & 79.75 & 78.22 & 77.87
& 76.79 & 73.47 & 72.03 & 73.43 & 71.07 & 71.44 & 70.78 & 69.63 & 67.96 & 67.98 \\
& \quad $+\mathrm{AS}^{(1)}_S$
& 96.00 & 92.17 & 88.33 & 86.92 & 85.87 & 85.11 & 86.33 & 85.17 & 84.78 & 84.03
& 82.85 & 82.56 & 81.54 & 81.52 & 80.64 & 79.73 & 79.59 & 80.32 & 79.58 & 79.47 \\

\addlinespace[2pt]

\multirow{3}{*}{\textbf{OLoRA}}
& SGD
& 92.33 & 88.50 & 87.78 & 85.59 & 83.00 & 81.28 & 80.57 & 78.58 & 75.67 & 74.10
& 73.55 & 73.36 & 73.46 & 73.33 & 71.87 & 71.40 & 70.57 & 68.33 & 67.44 & 65.12 \\
& \quad $+\mathrm{AS}^{(0)}_S$
& 92.33 & 88.67 & 90.22 & 85.67 & 86.33 & 83.00 & 81.48 & 79.83 & 78.26 & 78.33
& 77.00 & 73.97 & 71.77 & 72.50 & 70.51 & 70.36 & 69.33 & 68.54 & 66.21 & 66.98 \\
& \quad $+\mathrm{AS}^{(1)}_S$
& 96.00 & 91.17 & 88.00 & 86.42 & 85.47 & 85.83 & 86.33 & 84.96 & 84.19 & 83.07
& 82.30 & 81.53 & 81.46 & 80.69 & 80.27 & 78.75 & 79.14 & 78.98 & 79.28 & 79.35 \\

\bottomrule
\end{tabular}
\end{table*}


\begin{table*}[ht]
\centering
\caption{Raw task-wise $\operatorname{ACC}^{\mathrm{cls}}_t$ trajectories ($t=1,\dots,20$) on ImageNet-P under different optimizers (mean over repeated runs).}
\label{tab:taskwise_acc_new}
\setlength{\tabcolsep}{2.2pt}
\renewcommand{\arraystretch}{0.92}
\scriptsize
\begin{tabular}{@{}l l *{20}{r}@{}}
\toprule
\textbf{Method} & \textbf{Opt.} &
\textbf{t01} & \textbf{t02} & \textbf{t03} & \textbf{t04} & \textbf{t05} &
\textbf{t06} & \textbf{t07} & \textbf{t08} & \textbf{t09} & \textbf{t10} &
\textbf{t11} & \textbf{t12} & \textbf{t13} & \textbf{t14} & \textbf{t15} &
\textbf{t16} & \textbf{t17} & \textbf{t18} & \textbf{t19} & \textbf{t20} \\
\midrule

\multirow{3}{*}{\textbf{SeqLoRA}}
& SGD
& 96.17 & 93.42 & 92.00 & 89.88 & 88.17 & 87.86 & 86.43 & 84.92 & 85.57 & 84.63
& 83.88 & 83.71 & 82.40 & 82.02 & 83.08 & 81.95 & 80.09 & 80.17 & 78.01 & 72.76 \\
& \quad $+\mathrm{AS}^{(0)}_S$
& 95.83 & 92.67 & 92.28 & 89.54 & 88.47 & 88.58 & 87.90 & 87.29 & 87.63 & 86.82
& 85.80 & 85.60 & 84.57 & 84.21 & 84.51 & 83.58 & 83.30 & 83.70 & 82.32 & 80.95 \\
& \quad $+\mathrm{AS}^{(1)}_S$
& 99.17 & 94.92 & 95.00 & 92.71 & 91.60 & 91.92 & 89.55 & 88.69 & 89.29 & 89.10
& 87.70 & 87.71 & 87.68 & 87.62 & 88.19 & 87.51 & 86.89 & 87.02 & 86.93 & 86.08 \\

\addlinespace[2pt]

\multirow{3}{*}{\textbf{IncLoRA}}
& SGD
& 95.67 & 92.50 & 91.78 & 89.71 & 88.80 & 88.08 & 85.69 & 85.29 & 84.59 & 83.25
& 80.94 & 78.58 & 77.72 & 77.79 & 74.13 & 73.24 & 71.70 & 69.29 & 68.85 & 66.92 \\
& \quad $+\mathrm{AS}^{(0)}_S$
& 95.83 & 91.58 & 92.17 & 89.62 & 89.27 & 88.28 & 86.57 & 85.31 & 85.56 & 85.33
& 83.95 & 81.81 & 80.01 & 78.26 & 78.75 & 77.08 & 77.34 & 75.79 & 76.53 & 73.75 \\
& \quad $+\mathrm{AS}^{(1)}_S$
& 98.83 & 96.41 & 95.39 & 93.33 & 91.63 & 90.97 & 89.81 & 88.83 & 88.78 & 88.63
& 87.46 & 86.27 & 86.58 & 86.15 & 86.63 & 86.42 & 86.45 & 85.54 & 85.34 & 85.09 \\

\addlinespace[2pt]

\multirow{3}{*}{\textbf{OLoRA}}
& SGD
& 95.67 & 90.00 & 90.39 & 88.33 & 88.20 & 87.61 & 85.76 & 85.02 & 84.45 & 83.37
& 79.97 & 77.76 & 76.68 & 76.64 & 72.82 & 72.54 & 70.40 & 68.16 & 67.61 & 65.55 \\
& \quad $+\mathrm{AS}^{(0)}_S$
& 95.83 & 90.17 & 91.06 & 88.46 & 88.27 & 87.86 & 86.74 & 84.96 & 85.37 & 84.95
& 83.74 & 82.61 & 81.01 & 79.58 & 78.79 & 77.21 & 76.96 & 75.30 & 75.73 & 73.32 \\
& \quad $+\mathrm{AS}^{(1)}_S$
& 98.83 & 95.33 & 94.39 & 91.46 & 90.17 & 90.14 & 88.60 & 87.48 & 87.02 & 87.23
& 86.23 & 85.47 & 85.99 & 85.79 & 86.19 & 86.23 & 86.01 & 85.29 & 84.61 & 84.47 \\
\bottomrule
\end{tabular}
\end{table*}





\subsection{Computational Cost Analysis}


Table~\ref{tab:lora_opt_metric} reports the average per-task training time on ImageNet-R.
Across LoRA variants, adding sharpness regularization increases the computational cost relative to SGD, with $AS^{(1)}_S$ being consistently the most expensive due to the additional gradient-norm maximization, while $AS^{(0)}_S$ shows a moderate overhead.
In contrast, $RS^{(0)}_S$ incurs a smaller and more variable overhead, and can be comparable to SGD in some settings.


\begin{table}[hpb]
\centering
\caption{Average per-task training time (seconds) across methods and optimizers on ImageNet-R.}
\label{tab:lora_opt_metric}
\footnotesize
\setlength{\tabcolsep}{3pt}
\renewcommand{\arraystretch}{0.95}
\begin{tabular}{l cccc}
\toprule
\textbf{Method} & \textbf{SGD} & $+\mathrm{RS}^{(0)}_S$ & $+\mathrm{AS}^{(0)}_S$ & $+\mathrm{AS}^{(1)}_S$ \\
\midrule
SeqLoRA & $282.72 \pm 61.34$  & $244.95 \pm 6.38$  & $472.29 \pm 83.70$  & $670.20 \pm 217.92$ \\
IncLoRA & $310.54 \pm 141.00$ & $658.34 \pm 83.90$ & $777.91 \pm 239.29$ & $1247.56 \pm 5.05$ \\
OLoRA   & $412.99 \pm 146.84$ & $644.19 \pm 55.05$ & $628.44 \pm 181.29$ & $836.99 \pm 327.65$ \\
\bottomrule
\end{tabular}
\end{table}

% \begin{table}[hpb]
% \centering
% \caption{Average per‑task training time (seconds) across methods and optimizers on ImageNet‑R}
% \label{tab:lora_opt_metric}
% \setlength{\tabcolsep}{4pt}
% \renewcommand{\arraystretch}{1}
% \resizebox{\linewidth}{!}{
% \begin{tabular}{l ccccc}
% \toprule
% \textbf{Method} & \textbf{SGD} & $+RS^{(0)}_S$ & $+AS^{(0)}_S$ & $+AS^{(1)}_S$ \\
% \midrule
% SeqLoRA & 282.72 ± 61.34  & 244.95 ± 6.38  & 472.29 ± 83.70  & 670.20 ± 217.92    \\
% IncLoRA & 310.54 ± 141.00 & 658.34 ± 83.90 & 777.91 ± 239.29 & 1247.56 ± 5.05  3 \\
% OLoRA   & 412.99 ± 146.84 & 644.19 ± 55.05 & 628.44 ± 181.29 & 836.99 ± 327.65  \\
% \bottomrule
% \end{tabular}
% }
% \end{table}


\end{document}

% This document was modified from the file originally made available by
% Pat Langley and Andrea Danyluk for ICML-2K. This version was created
% by Iain Murray in 2018, and modified by Alexandre Bouchard in
% 2019 and 2021 and by Csaba Szepesvari, Gang Niu and Sivan Sabato in 2022.
% Modified again in 2023 and 2024 by Sivan Sabato and Jonathan Scarlett.
% Previous contributors include Dan Roy, Lise Getoor and Tobias
% Scheffer, which was slightly modified from the 2010 version by
% Thorsten Joachims & Johannes Fuernkranz, slightly modified from the
% 2009 version by Kiri Wagstaff and Sam Roweis's 2008 version, which is
% slightly modified from Prasad Tadepalli's 2007 version which is a
% lightly changed version of the previous year's version by Andrew
% Moore, which was in turn edited from those of Kristian Kersting and
% Codrina Lauth. Alex Smola contributed to the algorithmic style files.
